%% GENERATED by worksheets/make-problems.py --mode problems from the worksheet bodies. Do not edit.

\documentclass[11pt]{amsart}
\newif\ifsolutions\solutionstrue

%% ---- preamble ----
%% ======================================================================
%% Canonical preamble - Santiago Arango-Piñeros
%% Source of truth: ~/Dropbox/templates/preamble.tex
%% Copied into new projects by /newpaper. Fix things HERE, not only in the
%% project copy that happened to surface the problem.
%%
%% Merged from projects 15-APZB (base: Magma listings, LMFDB links, theorem
%% environments) and 8/9/12/13 (the shared macro core). Collisions resolved
%% in favour of the APZB definitions.
%% ======================================================================


\usepackage[T1]{fontenc}   % accented/ogonek names (\k{a}, \v{C}) need T1, not OT1
\usepackage[utf8]{inputenc}
\usepackage{booktabs,array}
\usepackage{amssymb}
\usepackage{amsmath}
\usepackage{amsthm}
\usepackage{stmaryrd}
\usepackage{fullpage}	% smaller margins %incompatible with ucthesis
\usepackage{amscd}   	% for commutative diagrams
\usepackage[all]{xy}    % for complicated commutative diagrams
\usepackage{comment} 	% for \begin{comment} ... \end{comment}
\usepackage{mathrsfs}   % for \mathscr
\usepackage{placeins}   % Defines a \FloatBarrier command, beyond which floats may not pass;
\usepackage[shortlabels]{enumitem} % key=value list options (leftmargin=*, label=...);
                                  % shortlabels keeps the old [(a)] / [(i)] syntax
\usepackage{mathtools}
\usepackage{colonequals}
\usepackage{xspace}
\usepackage{nicefrac}
\usepackage{booktabs}
\usepackage{caption}
% Table captions go below the table (as in every document here). Without this,
% caption assumes amsart's top placement and puts its skip on the wrong side,
% so the caption sits flush against the bottom rule.
\captionsetup[table]{position=below}
\usepackage{tikz}
\usetikzlibrary{cd}

\usepackage[alphabetic,backrefs,lite,msc-links]{amsrefs} % for bibliography
\usepackage{color}
\usepackage{tcolorbox}
\usepackage{colortbl}
\usepackage{setspace}
%\usepackage[dvipsnames, table, xcdraw]{xcolor}
\usepackage{hyperref}
\definecolor{headercolor}{RGB}{255,255,240}
\definecolor{mycitecolor}{RGB}{169,169,169}
\definecolor{mylinkcolor}{rgb}{1.0, 0.0, 0.5}
\definecolor{myurlcolor}{rgb}{1.0, 0.0, 0.5}
\hypersetup{colorlinks=true,
urlcolor=myurlcolor,
citecolor=mycitecolor,
linkcolor=mylinkcolor,
linktoc=page,
breaklinks=true}
\usepackage{url} 
\urlstyle{same}
\usepackage{cleveref} % automatic referencing
\usepackage{indentfirst} % to indent the first paragraph of a section
\usepackage[mathscr]{eucal} % nice caligraphy for sheaves and categories
\usepackage{colonequals} % correct way to write :=
\usepackage{soul} % to strikethrough words: \st{}
\usepackage{marvosym} % at work symbol
\usepackage{listings} % source-code printing
\usepackage{upquote}  % straight quotes inside listings (copy-pasteable code)

% ------------------------------------------------------------------------
% MAGMA CODE LISTINGS
% ------------------------------------------------------------------------
% listings has no built-in Magma language, so we define one below.  Usage:
%     \begin{magma} ... \end{magma}      display code
%     \mg{PolynomialRing}                inline code
% Inside a listing, (* ... *) escapes back to LaTeX, so comments can carry
% real mathematics:   // the curve (*$X_\QQ$*)

\definecolor{magmabuiltin}{rgb}{0.10,0.20,0.45} % muted navy
\definecolor{magmacomment}{gray}{0.45}
\definecolor{magmastring}{rgb}{0.10,0.40,0.25}  % muted green
\definecolor{magmarule}{gray}{0.72}

\lstdefinelanguage{Magma}{%
  sensitive=true,
  morecomment=[l]{//},
  morecomment=[s]{/*}{*/},
  morestring=[b]",
  % Reserved words of the language (Magma handbook, "Reserved Words"),
  % together with the angle-bracket constructors hom<>, sub<>, quo<>, ...
  morekeywords=[1]{%
    adj,and,assert,assert2,assert3,assigned,break,by,case,cat,catch,clear,
    cmpeq,cmpne,continue,declare,default,delete,diff,div,do,elif,else,end,eq,
    error,eval,exists,exit,false,for,forall,forward,fprintf,freeze,function,
    ge,gt,if,iload,import,in,intrinsic,is,join,le,load,local,lt,meet,mod,ne,
    not,notadj,notin,notsubset,or,print,printf,procedure,quit,random,read,
    readi,repeat,require,requirege,requirerange,restore,return,save,sdiff,
    select,subset,then,time,to,true,try,until,vprint,vprintf,vtime,when,
    where,while,xor,
    car,cop,elt,ext,hom,ideal,iso,lideal,map,ncl,pmap,quo,rec,recformat,
    rideal,sub},
  % Built-in intrinsics.  Extend this list as the paper grows.
  morekeywords=[2]{%
    Integers,Rationals,Reals,Complexes,IntegerRing,RationalField,RingOfIntegers,
    FiniteField,GF,NumberField,CyclotomicField,PolynomialRing,FunctionField,
    FieldOfFractions,AffineSpace,ProjectiveSpace,Scheme,Curve,EllipticCurve,
    HyperellipticCurve,Jacobian,GenusOneModel,cInvariants,bInvariants,
    Invariants,MinimalModel,WeierstrassModel,jInvariant,Discriminant,Conductor,
    Rank,Regulator,TorsionSubgroup,MordellWeilGroup,Generators,RationalPoints,
    Points,Genus,Degree,Dimension,BaseChange,BaseRing,Parent,Factorization,
    Factorisation,GCD,LCM,Valuation,Sqrt,Roots,Coefficients,Evaluate,Numerator,
    Denominator,Matrix,Vector,Eltseq,Domain,Codomain,Kernel,Image,Order,Index,
    IsSquare,IsPrime,IsIrreducible,IsIsomorphic,IsSingular,Automorphisms,
    AutomorphismGroup,DefiningEquations,Ideal,Set,Sequence,SetToSequence,
    Sprintf,Random,Abs,Max,Min,Support,Divisor,
    Basis,ChangeRing,Identity,IsNonSingular,IsomorphismData,Place,Places,
    RiemannRochSpace,TraceOfFrobenius,Type},
}

\lstdefinestyle{magmastyle}{%
  language=Magma,
  basicstyle=\ttfamily\footnotesize,
  keywordstyle=[1]\bfseries,                       % reserved words
  keywordstyle=[2]\color{magmabuiltin},            % intrinsics
  commentstyle=\itshape\color{magmacomment},
  stringstyle=\color{magmastring},
  % columns=fullflexible stops listings from letter-spacing the text onto a
  % fixed grid; keepspaces then preserves the indentation.
  columns=fullflexible,
  keepspaces=true,
  upquote=true,
  showspaces=false,
  showstringspaces=false,
  showtabs=false,
  tabsize=2,
  breaklines=true,
  breakatwhitespace=true,
  postbreak=\mbox{\textcolor{magmacomment}{$\hookrightarrow$}\space},
  frame=tb,
  framerule=0.4pt,
  rulecolor=\color{magmarule},
  framesep=6pt,
  xleftmargin=\parindent,
  aboveskip=1.1em,
  belowskip=1.1em,
  numbers=none,
  captionpos=b,
  escapeinside={(*}{*)},
}

\lstset{style=magmastyle}

\lstnewenvironment{magma}[1][]{\lstset{style=magmastyle,#1}}{}
% \mg{...} is for short identifiers and expressions.  Its argument is read as
% ordinary LaTeX, so for code containing #, %, ^ or _ use \lstinline!...!.
\newcommand{\mg}[1]{\lstinline[style=magmastyle]{#1}}



%%%%% Macros %%%%%

\newcommand{\pb}{\rule{.4pt}{5.4pt}\rule[5pt]{5pt}{.4pt}\llap{$\cdot$\hspace{1pt}}} % for cartesian diagrams
\newcommand{\defi}[1]{{\color{blue} \textsf{#1}}} 	% for defined terms
\newcommand{\cdef}[1]{\defi{#1}}    % for defined terms
\newcommand{\known}{{\operatorname{known}}}
\newcommand{\Abar}{{\overline{A}}}
\newcommand{\To}{\longrightarrow}
\newcommand{\nichts}{\ensuremath{\left.\right.}}
\newcommand{\ssigma}{{\! \nichts^\sigma \!}}
\newcommand{\YY}{{\mathcal Y}}
\newcommand{\prim}{{\operatorname{prim}}}
\newcommand{\fake}{{\operatorname{fake}}}
\newcommand{\conts}{{\operatorname{conts}}}
\newcommand{\wedgeisom}{{\;\stackrel{\wedge}\isom}\;}
\newcommand{\x}{\mathsf{x}}
\newcommand{\y}{\mathsf{y}}
\newcommand{\z}{\mathsf{z}}
\newcommand{\w}{\mathsf{w}}
\renewcommand{\t}{\mathsf{t}}
\newcommand{\xbar}{{\bar{x}}}
\newcommand{\ybar}{{\bar{y}}}
\newcommand{\zbar}{{\bar{z}}}
\newcommand{\Xtilde}{{\tilde{X}}}
\newcommand{\pitilde}{{\tilde{\pi}}}
\newcommand{\subs}{{\operatorname{subset}}}
\newcommand{\umtp}[1]{#1^\times\!/(#1^\times)^3}
\newcommand{\brk}[1]{\left\{ #1 \right\}}



% Characters
\newcommand{\Aff}{{\mathbb A}}
\newcommand{\C}{{\mathbb C}}
\newcommand{\F}{{\mathbb F}}
\newcommand{\G}{{\mathbb G}}
\newcommand{\bbH}{{\mathbb H}}
\newcommand{\PP}{{\mathbb P}}
\newcommand{\Q}{{\mathbb Q}}
\newcommand{\QQ}{\Q}
\newcommand{\R}{{\mathbb R}}
\newcommand{\Z}{{\mathbb Z}}
\newcommand{\ZZ}{\Z}
\newcommand{\ZS}{\Z_{\calS}}
\newcommand{\OKS}{\OO_{K,\calS}}
\newcommand{\OLS}{\OO_{L,\calS}}
\newcommand{\N}{{\mathbb{N}}}
\newcommand{\Qbar}{{\overline{\Q}}}
\newcommand{\Zhat}{{\widehat{\Z}}}
\newcommand{\Qhat}{{\widehat{\Q}}}
\newcommand{\Zbar}{{\overline{\Z}}}
\newcommand{\kbar}{{\overline{k}}}
\newcommand{\Kbar}{{\overline{K}}}
\newcommand{\ksep}{{k^{\operatorname{sep}}}}
\newcommand{\Ksep}{{K^{\operatorname{sep}}}}
\newcommand{\Fbar}{{\overline{\F}}}
\newcommand{\Adeles}{{\mathbf A}}
\newcommand{\kk}{{\mathbf k}}
\newcommand{\pp}{{\mathfrak p}}
\newcommand{\mfp}{\pp}
\newcommand{\mfP}{{\mathfrak P}}
\newcommand{\frakq}{{\mathfrak q}}
\newcommand{\qq}{{\mathfrak q}}
\newcommand{\mm}{{\mathfrak m}}


% mathcal characters
\newcommand{\calA}{{\mathcal A}}
\newcommand{\calB}{{\mathcal B}}
\newcommand{\calC}{{\mathcal C}}
\newcommand{\calD}{{\mathcal D}}
\newcommand{\calE}{{\mathcal E}}
\newcommand{\calF}{{\mathcal F}}
\newcommand{\calG}{{\mathcal G}}
\newcommand{\calH}{{\mathcal H}}
\newcommand{\calI}{{\mathcal I}}
\newcommand{\calJ}{{\mathcal J}}
\newcommand{\calK}{{\mathcal K}}
\newcommand{\calL}{{\mathcal L}}
\newcommand{\calM}{{\mathcal M}}
\newcommand{\calN}{{\mathcal N}}
\newcommand{\calO}{{\mathcal O}}
\newcommand{\calP}{{\mathcal P}}
\newcommand{\calQ}{{\mathcal Q}}
\newcommand{\calR}{{\mathcal R}}
\newcommand{\calS}{{\mathcal S}}
\newcommand{\calT}{{\mathcal T}}
\newcommand{\calU}{{\mathcal U}}
\newcommand{\mcU}{\calU}
\newcommand{\calV}{{\mathcal V}}
\newcommand{\calW}{{\mathcal W}}
\newcommand{\calX}{{\mathcal X}}
\newcommand{\calY}{{\mathcal Y}}
\newcommand{\calZ}{{\mathcal Z}}
\newcommand{\FF}{{\mathcal F}}
\newcommand{\GG}{{\mathcal G}}
\newcommand{\II}{{\mathcal I}}
\newcommand{\JJ}{{\mathcal J}}
\newcommand{\LL}{{\mathcal L}}
\newcommand{\OO}{{\mathcal O}}
\renewcommand{\O}{\OO}

% mathscr characters
\newcommand{\scrA}{{\mathscr A}}
\newcommand{\scrB}{{\mathscr B}}
\newcommand{\scrC}{{\mathscr C}}
\newcommand{\scrD}{{\mathscr D}}
\newcommand{\scrE}{{\mathscr E}}
\newcommand{\scrF}{{\mathscr F}}
\newcommand{\scrG}{{\mathscr G}}
\newcommand{\scrH}{{\mathscr H}}
\newcommand{\scrI}{{\mathscr I}}
\newcommand{\scrJ}{{\mathscr J}}
\newcommand{\scrK}{{\mathscr K}}
\newcommand{\scrL}{{\mathscr L}}
\newcommand{\scrM}{{\mathscr M}}
\newcommand{\scrN}{{\mathscr N}}
\newcommand{\scrO}{{\mathscr O}}
\newcommand{\scrP}{{\mathscr P}}
\newcommand{\scrQ}{{\mathscr Q}}
\newcommand{\scrR}{{\mathscr R}}
\newcommand{\scrS}{{\mathscr S}}
\newcommand{\scrT}{{\mathscr T}}
\newcommand{\scrU}{{\mathscr U}}
\newcommand{\scrV}{{\mathscr V}}
\newcommand{\scrW}{{\mathscr W}}
\newcommand{\scrX}{{\mathscr X}}
\newcommand{\X}{\scrX}
\newcommand{\scrY}{{\mathscr Y}}
\newcommand{\scrZ}{{\mathscr Z}}

\def\Q{\mathbb{Q}}
\def\C{\mathbb{C}}
\def\R{\mathbb{R}}
\def\P{\mathbb{P}}
\def\Pone{\P^{1}}
\def\A{\mathbb{A}}
% Affine space.  \AA was LaTeX's text-mode Angstrom sign; \textAA still gives it.
% Must be \renewcommand, not \def: lib/gen-macros.el in the elliptic-curves
% notes only reads \newcommand/\renewcommand, so a \def never reaches MathJax.
\let\textAA\AA
\renewcommand{\AA}{{\mathbb A}}
\def\Z{\mathbb{Z}}
\def\Qb{\bar{\mathbb{Q}}}
\def\kb{\bar{k}}
\def\hb{\hat{h}}


%%%%% Math operators %%%%%

\DeclareMathOperator{\fix}{fix}
\DeclareMathOperator{\ns}{ns}
\DeclareMathOperator{\spc}{sp}
\DeclareMathOperator{\sfr}{sf} 
\DeclareMathOperator{\tr}{tr} 
\DeclareMathOperator{\trdeg}{tr deg}
\DeclareMathOperator{\coker}{coker} 
\DeclareMathOperator{\rk}{rk}
\DeclareMathOperator{\Char}{char} 
\DeclareMathOperator{\inv}{inv}
\DeclareMathOperator{\nil}{Nil} 
\DeclareMathOperator{\im}{im}
\DeclareMathOperator{\re}{Re} 
\DeclareMathOperator{\End}{End}
\DeclareMathOperator{\Jac}{Jac} 
\DeclareMathOperator{\Prym}{Prym} 
\DeclareMathOperator{\END}{\bf End}
\DeclareMathOperator{\Lie}{Lie} \DeclareMathOperator{\Hom}{Hom}
%\DeclareMathOperator{\HOM}{\bf Hom} 
\DeclareMathOperator{\HOM}{\scrH om} 
\DeclareMathSymbol{\minus}{\mathbin}{AMSa}{"39}
\DeclareMathOperator{\Aut}{Aut}
\DeclareMathOperator{\Isom}{Isom}
\DeclareMathOperator{\Gal}{Gal}
% Frobenius and Verschiebung are symbols, not operators applied to an argument:
% ordinary atoms, so that \Frob^2 - [a]\Frob + [q] = 0 keeps its binary spacing.
\newcommand{\Frob}{\mathsf{F}}    % p^r-power Frobenius isogeny E \to E^{(p^r)}
\newcommand{\Versch}{\mathsf{V}}  % Verschiebung, the dual of Frobenius
\DeclareMathOperator{\Ind}{Ind}
\DeclareMathOperator{\Inf}{Inf}
\DeclareMathOperator{\Cor}{Cor}
\DeclareMathOperator{\Res}{Res}
\DeclareMathOperator{\Norm}{Norm} 
\DeclareMathOperator{\Br}{Br}
\DeclareMathOperator{\Gr}{Gr} 
\DeclareMathOperator{\cd}{cd}
\DeclareMathOperator{\scd}{scd} 
\DeclareMathOperator{\Sel}{Sel}
\DeclareMathOperator{\Cl}{Cl} 
\DeclareMathOperator{\divv}{div}
\DeclareMathOperator{\ord}{ord} 
\DeclareMathOperator{\Sym}{Sym}
\DeclareMathOperator{\Div}{Div} 
\DeclareMathOperator{\Pic}{Pic}
\DeclareMathOperator{\Num}{Num} 
\DeclareMathOperator{\PIC}{\bf Pic}
\DeclareMathOperator{\Spec}{Spec} 
\DeclareMathOperator{\SPEC}{\bf Spec} 
\DeclareMathOperator{\Proj}{Proj}
\DeclareMathOperator{\Frac}{Frac} 
\DeclareMathOperator{\rank}{rank}
\DeclareMathOperator{\CaCl}{CaCl}
\DeclareMathOperator{\Mod}{Mod}
\DeclareMathOperator{\Open}{Open}
\DeclareMathOperator{\Coh}{Coh}
\DeclareMathOperator{\id}{id}
\DeclareMathOperator{\eq}{eq}
\DeclareMathOperator{\Th}{th}
\DeclareMathOperator{\cond}{cond}
\DeclareMathOperator{\lcm}{lcm}



% Categories
\newcommand{\Schemes}{\operatorname{\bf Schemes}}
\newcommand{\Sch}{\operatorname{\bf Sch}}
\newcommand{\Sets}{\operatorname{\bf Sets}}
%\newcommand{\Ab}{\operatorname{\bf Ab}}
\newcommand{\Ab}{\operatorname{Ab}}
\newcommand{\Groups}{\operatorname{\bf Groups}}

% Text subscripts, superscripts
\newcommand{\perf}{{\operatorname{perf}}}
\newcommand{\op}{{\operatorname{op}}}
\newcommand{\an}{{\operatorname{an}}}
\newcommand{\ab}{{\operatorname{ab}}}
\newcommand{\smooth}{{\operatorname{smooth}}}
\newcommand{\sing}{{\operatorname{sing}}}
\newcommand{\CS}{\operatorname{\bf CS}}
\newcommand{\Zar}{{\operatorname{Zar}}}
\newcommand{\fl}{{\operatorname{f\textcompwordmark l}}}
\newcommand{\et}{{\operatorname{et}}}
\newcommand{\Az}{{\operatorname{Az}}}
\newcommand{\ET}{{\operatorname{\bf \acute{E}t}}}
\newcommand{\tors}{{\operatorname{tors}}}
\newcommand{\red}{{\operatorname{red}}}
\newcommand{\unr}{{\operatorname{unr}}}
\newcommand{\tH}{{\operatorname{th}}}


\newcommand{\GalQ}{{\Gal}(\Qbar/\Q)}
\newcommand{\Cech}{\v{C}ech\xspace}
\newcommand{\HH}{{\operatorname{H}}}
\newcommand{\HHcech}{{\check{\HH}}}
\newcommand{\HHat}{{\hat{\HH}}}
\newcommand{\M}{\operatorname{M}}
\newcommand{\GL}{\operatorname{GL}}
\newcommand{\PGL}{\operatorname{PGL}}
\newcommand{\GSp}{\operatorname{GSp}}
\newcommand{\SL}{\operatorname{SL}}
\newcommand{\Sp}{\operatorname{Sp}}
\newcommand{\PSL}{\operatorname{PSL}}

\newcommand{\surjects}{\twoheadrightarrow}
\newcommand{\injects}{\hookrightarrow}
\newcommand{\isom}{\simeq}
\newcommand{\notdiv}{\nmid}
\newcommand{\del}{\partial}
\newcommand{\Intersection}{\bigcap}
\newcommand{\intersect}{\cap}
\newcommand{\Union}{\bigcup}
\newcommand{\union}{\cup}
\newcommand{\tensor}{\otimes}
\newcommand{\directsum}{\oplus}
\newcommand{\Directsum}{\bigoplus}

\newcommand{\isomto}{\overset{\sim}{\rightarrow}}
\newcommand{\isomfrom}{\overset{\sim}{\leftarrow}}
\newcommand{\leftexp}[2]{{\vphantom{#2}}^{#1}{#2}}
\newcommand{\sigmaiota}{{\leftexp{\sigma}{\iota}}}
\newcommand{\sigmaphi}{{\leftexp{\sigma}{\phi}}}
\newcommand{\tauphi}{{\leftexp{\tau}{\phi}}}
\newcommand{\sigmatauphi}{{\leftexp{\sigma\tau}{\phi}}}


\def\bfb{{\bf b}}
\def\bfv{{\bf v}}
\def\bfx{{\bf x}}
\def\bfy{{\bf y}}
\def\bfz{{\bf z}}

%!!!!!!!!!!!!!!!!!!!!!!!!!!!!!!!!!!!!!!!!!!!!!!!!!!!!!!!!!!!!!!!!!!!!!
%!!!!!!!!!!!!!!!!!!!!!!!!!!!!!!!!!!!!!!!!!!!!!!!!!!!!!!!!!!!!!!!!!!!!!

%From RSZB
\newcommand{\Magma}{\texttt{Magma}\xspace}
\newcommand{\magmacmd}[1]{\texttt{#1}}
\newcommand{\smallmat}[4]{\left[\begin{smallmatrix} #1 & #2 \\ #3 & #4 \end{smallmatrix}\right]}
\newcommand{\eclabel}[1]{\href{https://www.lmfdb.org/EllipticCurve/Q/#1}{\texttt{#1}}}
\newcommand{\ecnflabel}[2]{\href{https://www.lmfdb.org/EllipticCurve/#1-#2}{\texttt{#2}}}
\newcommand{\mflabel}[1]{\href{https://www.lmfdb.org/ModularForm/GL2/Q/holomorphic/#1}{\texttt{#1}}}
\newcommand{\nflabel}[1]{\href{https://www.lmfdb.org/NumberField/#1}{\texttt{#1}}}
\newcommand{\repolink}[2]{\href{https://github.com/AndrewVSutherland/ell-adic-galois-images/blob/main/#1/#2}{\texttt{#2}}}
\newcommand{\glink}[1]{\hyperlink{#1}{\texttt{#1}}}
\newcommand{\gtarget}[1]{\hypertarget{#1}{\texttt{#1}}}
\newcommand{\glabel}[1]{\texttt{#1}} % use for group links with no target
\newcommand{\XNgdis}{X(N)}
\newcommand{\JNgdis}{J(N)}
\newcommand{\XNconn}{X(N)'}
\newcommand{\JNconn}{J(N)'}

\DeclareMathOperator{\alg}{alg}

\newcommand{\kron}[2]{\left(\frac{#1}{#2}\right)}
\DeclareMathOperator{\disc}{disc}
\DeclareMathOperator{\Tr}{Tr}
\DeclareMathOperator{\trd}{trd} % reduced trace of a quaternion algebra
\DeclareMathOperator{\nrd}{nrd} % reduced norm
                               
\setcounter{MaxMatrixCols}{20}

%%% Uncomment these to number equations by section
\numberwithin{equation}{subsection}
\newtheorem{theorem}{Theorem}[section]
\numberwithin{theorem}{subsection}
\newtheorem{lemma}[theorem]{Lemma}
\newtheorem{corollary}[theorem]{Corollary}
\newtheorem{proposition}[theorem]{Proposition}

\theoremstyle{definition}
\newtheorem{definition}[theorem]{Definition}
\newtheorem{question}[theorem]{Question}
\newtheorem{answer}[theorem]{Answer}
\newtheorem{fact}[theorem]{Fact}
\newtheorem{conjecture}[theorem]{Conjecture}
\newtheorem{example}[theorem]{Example}
\newtheorem{algorithm}[theorem]{Algorithm}
%\newtheorem{para}[subsection]{}

%\theoremstyle{remark}
\newtheorem{remark}[theorem]{Remark}
\newtheorem{remarks}[theorem]{Remarks}
\newtheorem{warning}[theorem]{Warning}

\newcommand{\labelpar}[1]{\refstepcounter{subsection}\label{#1}\thesubsection.}
%% Ending here.



% %%% Uncomment these to number equations by section
% \numberwithin{equation}{section}
% \newtheorem{theorem}[equation]{Theorem}
% \newtheorem{lemma}[equation]{Lemma}
% \newtheorem{corollary}[equation]{Corollary}
% \newtheorem{proposition}[equation]{Proposition}

% \theoremstyle{definition}
% \newtheorem{definition}[equation]{Definition}
% \newtheorem{question}[equation]{Question}
% \newtheorem{answer}[equation]{Answer}
% \newtheorem{fact}[equation]{Fact}
% \newtheorem{conjecture}[equation]{Conjecture}
% \newtheorem{example}[equation]{Example}
% \newtheorem{para}[equation]{}

% \theoremstyle{remark}
% \newtheorem{remark}[equation]{Remark}
% \newtheorem{remarks}[equation]{Remarks}
% \newtheorem{warning}[equation]{Warning}
%\newcommand{\labelpar}[1]{\refstepcounter{equation}\label{#1}\theequation.}
% %% Ending here.


%%% Uncomment these to have unnumbered equations
%\newtheorem*{theorem}{Theorem}
%\newtheorem*{lemma}{Lemma}
%\newtheorem*{corollary}{Corollary}
%\newtheorem*{proposition}{Proposition}
%
%\theoremstyle{definition}
%\newtheorem*{definition}{Definition}
%\newtheorem*{question}{Question}
%\newtheorem*{conjecture}{Conjecture}
%\newtheorem*{example}{Example}
%
%\theoremstyle{remark}
%\newtheorem*{remark}{Remark}
%\newtheorem*{remarks}{Remarks}
%\newtheorem*{warning}{Warning}
%%% Ending here.

%\newtheorem{theorem}{Theorem}[section]
%\newtheorem{lemma}[theorem]{Lemma}
%\newtheorem{corollary}[theorem]{Corollary}
%\newtheorem{proposition}[theorem]{Proposition}
%
%\theoremstyle{definition}
%\newtheorem{definition}[theorem]{Definition}
%\newtheorem{question}[theorem]{Question}
%\newtheorem{conjecture}[theorem]{Conjecture}
%\newtheorem{example}[theorem]{Example}
%
%\theoremstyle{remark}
%\newtheorem{remark}[theorem]{Remark}
%\newtheorem{remarks}[theorem]{Remarks}
%\newtheorem{warning}[theorem]{Warning}



%% colors
% see https://latexcolor.com/ for more options

\definecolor{santicolor}{rgb}{1.0, 0.0, 0.16}
\definecolor{dzbcolor}{rgb}{0.0, 0.45, 0.73}

%% comments

\newcommand{\margin}[2]{\normalsize{{#2 \footnote{{#2 #1}}}{\marginpar[{#2\hfill\tiny\thefootnote$\rightarrow$}]
{{#2 $\leftarrow$\tiny\thefootnote}}}}}

% santi
\newcommand{\sapedit}[1]{{\color{santicolor} \sf  #1}}
\newcommand{\santi}[1]{{\color{santicolor} \MineSign \, \sf (Santi) #1}}
\newcommand{\Santi}[1]{\margin{(Santi) #1}{\color{santicolor}}}

% dzb
\newcommand{\dzbedit}[1]{{\color{dzbcolor} \sf  #1}}
\newcommand{\dzb}[1]{{\color{dzbcolor} \MineSign \, \sf (David) #1}}
\newcommand{\Dzb}[1]{\margin{(David) #1}{\color{dzbcolor}}}

%% footnotes
\renewcommand{\thefootnote}{\roman{footnote}}

% TODO NOTES
\usepackage[textsize=tiny]{todonotes}
\usepackage{xargs, xpatch}
\setlength{\marginparwidth}{2cm}
\newcommand{\td}[1]{\todo[fancyline,backgroundcolor=santicolor!25,bordercolor=santicolor]{#1}}
\newcommand{\tdi}[1]{\todo[inline,fancyline,backgroundcolor=santicolor!25,bordercolor=santicolor]{#1}}
\usepackage{marginnote}
\newcommand{\sidebar}[1]{\marginnote{\color{magenta}\scriptsize{$\ll$#1$\gg$}}}







%% ======================================================================
%% Shared macro core, merged in from the 8/9/12/13 lineage.
%% Definitions that clashed with the base above were dropped, and macros
%% that clash with LaTeX kernel commands use \renewcommand.
%% ======================================================================

% ---- merged in from projects/9-counting-prim-sols/tex/frontmatter.tex ----
\newcommand{\AbGrps}{\texttt{AbGrps}} % category of abelian groups
\DeclareMathOperator{\Autsch}{\mathbf{Aut}} % Automorphism group scheme
\newcommand{\CC}{\mathscr{C}} % rigidification of Fermat stack
\newcommand{\DD}{\mathbf{D}}
\newcommand{\Fp}{\mathbb{F}_p}
\newcommand{\Ga}{\mathbb{G}_\text{a}}
\newcommand{\Gm}{\mathbb{G}_\text{m}}
\newcommand{\GrpSch}{\texttt{GrpSch}} % category of schemes
\DeclareMathOperator{\Ht}{Ht}
\DeclareMathOperator{\Id}{Id} % Id
\renewcommand{\Im}{\operatorname{Im}}
\DeclareMathOperator{\Nm}{Nm} % norm
\DeclareMathOperator{\PProj}{\mathbf{Proj}} % stacky Proj
\newcommand{\Qp}{{\Q_p}} % for p-adics
\renewcommand{\Re}{\operatorname{Re}}
\renewcommand{\SS}{\mathscr{S}} % mathscr S
\DeclareMathOperator{\Stab}{Stab} % Stabilizer
\newcommand{\XX}{\mathscr{X}} % mathscr X
\newcommand{\abc}{(a,b,c)}
\newcommand{\belyi}{Belyi } % Belyi
\newcommand{\belyis}{Belyi's } % Belyi's
\newcommand{\bfa}{\mathbf{a}}
\newcommand{\bfc}{\mathbf{c}}
\newcommand{\bfe}{\mathbf{e}}
\newcommand{\bfj}{\mathbf{j}}
\newcommand{\bfs}{\mathbf{s}}
\newcommand{\bfw}{\mathbf{w}}
\newcommand{\bfX}{\mathbf{X}}
\newcommand\cHH{\check{\mathrm{H}}}
\DeclareMathOperator{\chr}{char} % characteristic
\newcommand{\coeff}{\mathbf{sfc}} % Fermat coefficient
\newcommand{\conprod}[1]{\stackon{$\times$}{\tiny $#1$}} % contracted product
\newcommand{\dd}{\,\mathrm{d}} % "d" of differential
\DeclareMathOperator{\den}{den}
% \div is LaTeX's division sign; the 8/9/12/13 papers use it for the principal
% divisor. \let ... \relax first, or \DeclareMathOperator refuses to redefine it.
\let\div\relax
\DeclareMathOperator{\div}{div} % principal divisor (\divv is the same operator)
\renewcommand{\emptyset}{\varnothing}
\newcommand{\fppf}{\texttt{fppf}} % faithfully flat and locally of finite presentation
\renewcommand{\geq}{\geqslant}
\newcommand{\gfe}{A\x^{a} + B\y^{b} + C\z^{c}}
\newcommand{\ggfe}{c_1\x_1^{e_1} + c_2\x_2^{e_2} + c_3\x_3^{e_3}}
\newcommand{\ideal}[1]{\langle #1 \rangle} % ideal
\renewcommand{\leq}{\leqslant}
\newcommand{\mand}{\text{ and }} % "and" in mathmode
\newcommand{\mcE}{\mathcal{E}} % mathcal E
\newcommand{\mcF}{\mathcal{F}} % mathcal F
\newcommand{\mcH}{\mathcal{H}} % mathcal H
\newcommand{\mcI}{\mathcal{I}} % mathcal I
\newcommand{\mcL}{\mathcal{L}} % mathcal L
\newcommand{\mcO}{\mathcal{O}} % sheaf O
\newcommand{\mcP}{\mathcal{P}} % mathcal P
\newcommand{\mcQ}{\mathcal{Q}} % mathcal Q
\newcommand{\mcR}{\mathcal{R}} % mathcal R
\newcommand{\mcS}{\mathcal{S}} % mathcal S
\newcommand{\mcT}{\mathcal{T}} % mathcal T
\newcommand{\mcV}{\mathcal{V}} % mathcal V
\newcommand{\mcX}{\mathcal{X}} % mathcal X
\newcommand{\md}{\operatorname{mod}} % mod without parenthesis
\newcommand{\mfa}{\mathfrak{a}}
\newcommand{\mfb}{\mathfrak{b}}
\newcommand{\mfc}{\mathfrak{c}}
\newcommand{\mfm}{\mathfrak{m}} % ideal m
\newcommand{\mfor}{\text{ for }} % "for" in mathmode
\newcommand{\mforall}{ \text{ for all }} % "for all" in mathmode
\newcommand{\mfpb}{\bar{\mathfrak{p}}}
\newcommand{\mif}{\text{ if }} % "if" in mathmode
\renewcommand{\mod}{\md}
\newcommand{\mof}{\text{ of }} % "of" in mathmode
\newcommand{\nfield}[1]{\href{https://www.lmfdb.org/NumberField/#1}{\texttt{#1}}}
\DeclareMathOperator{\num}{num}
\newcommand{\oalpha}{\bar{\alpha}}
\newcommand{\paren}[1]{\left( #1 \right)} % parenthesis
\newcommand{\pibar}{\bar{\pi}}
\DeclareMathOperator{\rad}{rad}
\DeclareMathOperator{\rig}{rig}
\newcommand{\rmp}{\mathrm{p}} % roman p
\newcommand{\s}{\mathsf{s}}
\newcommand{\sfC}{\mathsf{C}} % mathsf C
\newcommand{\sfG}{\mathsf{G}} % mathsf G
\newcommand{\sfH}{\mathsf{H}} % mathsf H
\newcommand{\sfI}{\mathsf{I}} % mathsf I
\newcommand{\sfK}{\mathsf{K}} % mathsf K
\newcommand{\sfM}{\mathsf{M}} % mathsf M
\newcommand{\sfS}{\mathsf{S}} % mathsf S
\newcommand{\sfX}{\mathsf{X}}
\newcommand{\sheafP}{\mathcal{P}} % reserved for a torsor
\DeclareMathOperator{\sign}{sign}
\newcommand{\spcite}[1]{\cite[\href{https://stacks.math.columbia.edu/tag/#1}{Tag
    #1}]{stacks-project}}
\newcommand{\sqbrk}[1]{\left[ #1 \right]} % square brackets
\renewcommand{\th}{\text{th}}
\newcommand{\thickslash}{\mathbin{\!\!\pmb{\fatslash}}}
\newcommand{\tri}{\triangle} % extended triangle group
\newcommand{\tribar}{\bar\triangle} % triangle group
\newcommand{\ur}{\texttt{ur}} % unramified
\DeclareMathOperator{\vol}{vol} % volume

% ---- merged in from projects/13-counting-euclidean-fermat/overleaf/preamble.tex ----
\newcommand{\Fpbar}{\overline{\mathbb{F}}_p}
\newcommand{\Gw}[1][]{\mathbb{G}_{\bfw  #1}}
\newcommand{\Reg}{\operatorname{Reg}}
\newcommand{\Vol}{\operatorname{Vol}}
\newcommand{\Xone}{X(\Q)_{1}}
\newcommand{\Zp}{{\Z_p}} % for p-adic numbers
\newcommand{\bad}{\mfb}
\newcommand{\chiabc}{\tfrac1a+\tfrac1b+\tfrac1c-1}
\newcommand{\chie}{\tfrac{1}{e_1} + \tfrac{1}{e_2} + \tfrac{1}{e_3} - 1}
\DeclareMathOperator{\covol}{covol} % volume
\DeclareMathOperator{\defect}{defect}
\newcommand{\eee}{(e_1, e_2, e_3)}
\DeclareMathOperator{\logHt}{ht}
\newcommand{\mcD}{\mathcal{D}} % mathcal D
\newcommand{\pr}{\mathrm{pr}}
\newcommand{\sfY}{\mathsf{Y}}
\newcommand{\sfu}{\mathsf{u}}
\newcommand{\sfv}{\mathsf{v}}
\newcommand{\sfw}{\mathsf{w}}


%% ======================================================================
%% amsrefs @misc fix
%% ----------------------------------------------------------------------
%% Two amsrefs bugs conspire to silently drop URLs from @misc entries:
%%
%%  1. amsrefs aliases  misc -> miscellaneous -> the `book' BibSpec, and the
%%     book spec has NO url line. So `url = {...}' in an @misc entry is
%%     parsed, accepted, and then never typeset.
%%  2. amsra.bst (which amsrefs forces -- it does \let\bibliographystyle
%%     \@gobble, so any \bibliographystyle{...} in your document is ignored)
%%     emits `howpublished' as the key `how', which amsrefs never defines.
%%     That one at least warns: "rkeyval Warning: Unknown key: \bib'how".
%%
%% Neither is an error, so a bibliography can lose every @misc URL and still
%% compile clean. Below: register `how', and give miscellaneous entries the
%% book spec plus `how' and `url' lines.
%%
%% NB the url line uses \texttt, not \url. Applying \url here re-reads text
%% whose catcodes amsrefs has already fixed, which recurses until TeX dies
%% with "input stack size exceeded". \nolinkurl fails the same way. \texttt
%% does not insert breakpoints, so a very long bare url= may overflow the
%% line; a url inside `howpublished' is already \url{}-wrapped and breaks
%% normally.
%% ======================================================================
\DefineSimpleKey{bib}{how}
\BibSpec{miscellaneous}{%
    +{}  {\PrintPrimary}                {transition}
    +{,} { \textit}                     {title}
    +{.} { }                            {part}
    +{:} { \textit}                     {subtitle}
    +{,} { \PrintEdition}               {edition}
    +{}  { \PrintEditorsB}              {editor}
    +{,} { \PrintTranslatorsC}          {translator}
    +{,} { \PrintContributions}         {contribution}
    +{,} { }                            {series}
    +{,} { \voltext}                    {volume}
    +{,} { }                            {publisher}
    +{,} { }                            {organization}
    +{,} { }                            {address}
    +{,} { \PrintDateB}                 {date}
    +{,} { }                            {status}
    +{,} { }                            {how}
    +{,} { \texttt}                     {url}
    +{}  { \parenthesize}               {language}
    +{}  { \PrintTranslation}           {translation}
    +{;} { \PrintReprint}               {reprint}
    +{.} { }                            {note}
    +{.} {}                             {transition}
    +{}  {\SentenceSpace \PrintReviews} {review}
}

\makeatletter\def\toclevel@part{0}\makeatother
\setcounter{tocdepth}{1}

\usepackage{eso-pic}
\AddToShipoutPictureFG*{%
  \put(\LenToUnit{1.1cm},\LenToUnit{.5\paperheight}){%
    \rotatebox{90}{\makebox[0pt][c]{\small\sffamily\color{black!50}%
      DRAFT \textemdash{} last edited October 8, 2026}}}}

%% ---- shared machinery of the worksheets ----

\tcbuselibrary{breakable}
\usepackage{pdfcol}
\usepackage{needspace}

\usepackage[shortlabels]{enumitem}
\setlist[itemize,enumerate]{leftmargin=*}
\setlist[enumerate,1]{label=(\arabic*), ref=(\arabic*), leftmargin=0pt, itemindent=*}
\setlist[enumerate,2]{label=(\alph*), ref=(\arabic{enumi})(\alph*)}

\numberwithin{equation}{section}

\newtheorem{problemenv}{Problem}
\newcounter{wsstatement}[section]
\renewcommand{\thewsstatement}{\thesection.\arabic{wsstatement}}
\newtheorem{wsexample}[wsstatement]{Example}
\newtheorem{wswarning}[wsstatement]{Warning}
\newtheorem{wsremark}[wsstatement]{Remark}
\crefname{problemenv}{Problem}{Problems}
\Crefname{problemenv}{Problem}{Problems}
\crefname{wsexample}{Example}{Examples}
\Crefname{wsexample}{Example}{Examples}
\crefname{wswarning}{Warning}{Warnings}
\Crefname{wswarning}{Warning}{Warnings}
\crefname{wsremark}{Remark}{Remarks}
\Crefname{wsremark}{Remark}{Remarks}

\theoremstyle{plain}
\newtheorem{wstheorem}[wsstatement]{Theorem}
\newtheorem{wscorollary}[wsstatement]{Corollary}
\newtheorem{wsproposition}[wsstatement]{Proposition}
\newtheorem{wslemma}[wsstatement]{Lemma}
\theoremstyle{definition}
\newtheorem{wsdefinition}[wsstatement]{Definition}
\crefname{wstheorem}{Theorem}{Theorems}
\Crefname{wstheorem}{Theorem}{Theorems}
\crefname{wscorollary}{Corollary}{Corollaries}
\Crefname{wscorollary}{Corollary}{Corollaries}
\crefname{wsproposition}{Proposition}{Propositions}
\Crefname{wsproposition}{Proposition}{Propositions}
\crefname{wslemma}{Lemma}{Lemmas}
\Crefname{wslemma}{Lemma}{Lemmas}
\crefname{wsdefinition}{Definition}{Definitions}
\Crefname{wsdefinition}{Definition}{Definitions}

\newenvironment{problem}[1][]%
  {\needspace{6\baselineskip}%
   \begin{tcolorbox}[breakable,use color stack,colback=white,colframe=black!55,
      boxrule=0.5pt,left=6pt,right=6pt,top=5pt,bottom=5pt]%
   \begin{problemenv}[#1]\hfill\par\ignorespaces}%
  {\end{problemenv}\end{tcolorbox}}

\definecolor{wssolution}{RGB}{25,45,95}
\ifsolutions
  \newenvironment{solution}%
    {\par\smallskip
     \begin{list}{}{\leftmargin=0pt\rightmargin=0pt\topsep=2pt\partopsep=0pt
       \parsep=\parskip\itemsep=0pt}%
     \color{wssolution}%
     \item[]{\small\scshape\color{black!55}Solution.}\enspace\ignorespaces}%
    {\end{list}\par\smallskip}
\else
  \excludecomment{solution}
\fi

\SetSymbolFont{stmry}{bold}{U}{stmry}{m}{n}
\newcommand{\ans}[1]{{\bfseries\boldmath #1}}

\newenvironment{solremark}%
  {\par\smallskip\small\noindent\textit{Remark.}\enspace\ignorespaces}%
  {\par}

\ifsolutions
  \newcommand{\hint}[1]{}
\else
  \newcommand{\hint}[1]{\footnote{\emph{Hint:} #1}}
\fi

% Hints stay in the problems file, as in the worksheets.
\renewcommand{\hint}[1]{\footnote{\emph{Hint:} #1}}

%% ---- macros of individual worksheets ----
\newcommand{\Cdbar}{\overline{C}_{d}}
\newcommand{\Cd}{C_d}
\newcommand{\trilin}[3]{\langle #1, #2, #3 \rangle}
\DeclareMathOperator{\discrd}{discrd}

\title{Math 714: problems}
\author{Your name here}
\date{\today}

\begin{document}

\begin{abstract}
% TODO: add the course's submission instructions after the first sentence.
Math 714 is the graduate course on \textit{The Arithmetic Of Elliptic Curves}
  taught at UMass Amherst. I taught this class in the Fall of 2026. The
  students read ahead, and we used the lecture time to work on the problems in
  this document. Each part below is one problem set (corresponding to one
  meeting), with its summaries, definitions, results and problems. Each problem
  is preceded by the definitions and results it uses. Write your solution to each item in the
  \textsc{Solution} block under it.
\end{abstract}
\maketitle

\subsection*{How to use this file}
\begin{itemize}
  \item Fill in your name above, and compile with pdf\LaTeX{}. Run it twice, so that the
        references and the table of contents are up to date. The file is self-contained:
        it needs no other file.
  \item Write your solutions inside the solution blocks, and keep the problem text.
  \item Problems are numbered $1, 2, \dots$ through the whole document: cite
        ``Problem~$n$'' in a submission, and use \verb|\cref| or \verb|\ref| with labels
        for your own references. Definitions and results are numbered by section.
\end{itemize}
% TODO: submission format, deadline, collaboration policy (add as items above).

\subsection*{An example}
A problem is a framed box with a title. Each item of the problem is followed by an empty
solution block, as in the example below (it is not part of the problem set: it is printed
as it appears in the file).
\begin{verbatim}
\begin{problem}[A toy problem]
  \label{prob:toy}
  \begin{enumerate}
    \item\label{it:toy-sum} Compute $1 + 2 + \dots + n$.
\begin{solution}
% Write your solution here.
\end{solution}
    \item Use \ref{it:toy-sum} to compute $1 + 2 + \dots + 100$.
\begin{solution}
% Write your solution here.
\end{solution}
  \end{enumerate}
\end{problem}
\end{verbatim}
You write between \verb|\begin{solution}| and \verb|\end{solution}|, and these two lines
must stay alone on their lines. A finished item looks like this:
\begin{verbatim}
\begin{solution}
\ans{$n(n+1)/2$.} Add the sum to itself in reverse order. Each of the
$n$ columns adds up to $n + 1$, so twice the sum is $n(n+1)$:
\[
  2\sum_{k=1}^{n} k = n(n+1).
\]
\begin{solremark}
Gauss is said to have found this as a schoolboy.
\end{solremark}
\end{solution}
\end{verbatim}
Here \verb|\ans{...}| states the answer in bold, in one line, before the argument.
\verb|\begin{solremark}| sets a side comment apart from the argument. Refer to a problem or
to one of its items by label, as in \verb|\Cref{prob:toy}| (which prints ``Problem~$n$'') or
\verb|\ref{it:toy-sum}|, and never type a number by hand. Give your own labels a prefix
such as \verb|my:|, so that they do not clash with the labels of the problem sets.

\subsection*{A note on how this was made}
The problems were chosen and curated by Santiago Arango-Pi\~neros. The first
draft typesetting and the solutions were written by Claude, and then reviewed
by the author. Corrections and recommendations are welcome: send them to
\begin{center}
\href{mailto:santiago.arango.pineros@gmail.com}{\texttt{santiago.arango.pineros@gmail.com}}.
\end{center}

\tableofcontents

\part[Selmer elliptic curves]{Selmer elliptic curves}
\label{selmer:part}
\noindent
The problems below study a single family of plane cubics and use it to
reconstruct, by hand, some of the main results of \cite{Silverman09}*{\S\S
  III.1--III.4}: Weierstrass equations, the group law, the Riemann--Roch
construction of a Weierstrass model, and a first endomorphism. \Cref{selmer:sec:magma} is a
set of worked \Magma\ examples to read and re-run afterwards, with three short
variations to try at the end.

\subsection*{Standing notation}
\begin{itemize}
  \item A field $k$ with $\Char k \neq 3$, an algebraic closure $\kbar$ of $k$, and
        $d \in k^{\times}$.
  \item The cubic form $F \coloneqq X^3 + Y^3 - dZ^3$, the plane cubic
        $\Cd \coloneqq \brk{F = 0} \subset \PP^2_k$, and its rational point
        $O \coloneqq (1:-1:0)$.
  \item A primitive cube root of unity $\zeta \in \kbar$.
  \item A \defi{Selmer elliptic curve}: a pair $(\Cd, O)$ of this shape, after
        \cite{Selmer-51}. \Cref{selmer:sec:name} explains the name.
  \item Projective coordinates $X, Y, Z$, as usual, and lowercase letters for affine
        coordinates.
  \item The base change $\Cdbar \coloneqq \Cd \times_k \kbar$.
\end{itemize}

\section{Warm-up: a plane cubic and a rational point on it}
\label{selmer:sec:warmup}

\begin{wsdefinition}[{Elliptic curves \cite{Silverman09}*{\S III.3}}]
\label{selmer:def:elliptic}
An \defi{elliptic curve} over $k$ is a pair $(E,O)$ with $E$ a smooth, projective,
geometrically irreducible curve of genus one over $k$ and $O \in E(k)$ a fixed rational
point.
\end{wsdefinition}

\begin{problem}[It is an elliptic curve]
\label{selmer:prob:nice}
\begin{enumerate}
  \item\label{selmer:it:nice-O} Check that $O \in \Cd(k)$.
\begin{solution}
% Write your solution here.
\end{solution}
  \item\label{selmer:it:nice-smooth} Compute the three partial derivatives of $F$ and deduce that $\Cd$ is
        smooth, using $\Char k \neq 3$ and $d \neq 0$. Where is each hypothesis
        used?
\begin{solution}
% Write your solution here.
\end{solution}
  \item\label{selmer:it:nice-char3} Suppose instead that $\Char k = 3$ (only for this problem). Show
        that over $\kbar$ the curve $\Cd$ is a \emph{triple
        line}.\hint{Frobenius.}
\begin{solution}
% Write your solution here.
\end{solution}
  \item\label{selmer:it:nice-irreducible} Show that a smooth plane curve is geometrically
        irreducible.\hint{B\'ezout.}
\begin{solution}
% Write your solution here.
\end{solution}
  \item\label{selmer:it:nice-elliptic} Conclude from the degree--genus formula that $(\Cd, O)$ is an elliptic
        curve over $k$.
\begin{solution}
% Write your solution here.
\end{solution}
\end{enumerate}
\end{problem}

The group law in \Cref{selmer:prob:flex}\ref{selmer:it:flex-grouplaw} is the chord--tangent
construction of \cite{Silverman09}*{\S III.2}.

\begin{wsproposition}[{The group law and $\Pic^0$ \cite{Silverman09}*{Proposition~III.3.4}}]
\label{selmer:prop:pic}
The map $P \mapsto [P - O]$ identifies the chord--tangent group law on $\Cd(\kbar)$ with
addition in $\Pic^0 \Cdbar$.
\end{wsproposition}

\begin{problem}[The base point is an inflection point]
\label{selmer:prob:flex}
\begin{enumerate}
  \item\label{selmer:it:flex-infinity} Show that the line $Z = 0$ meets $\Cd$ in three \emph{distinct} points
        of $\Cd(\kbar)$, and write them down. For which $k$ are all three in
        $\Cd(k)$?
\begin{solution}
% Write your solution here.
\end{solution}
  \item\label{selmer:it:flex-tangent} Compute the tangent line $L$ to $\Cd$ at $O$, and show that
        $L \cap \Cd = 3O$. (So $O$ is an \defi{inflection point}, or
        \defi{flex}.)
\begin{solution}
% Write your solution here.
\end{solution}
  \item\label{selmer:it:flex-weierstrass} In a Weierstrass model, which line plays the role of $L$? What is
        different here?
\begin{solution}
% Write your solution here.
\end{solution}
  \item\label{selmer:it:flex-grouplaw} Let $\oplus$ denote the group law on $\Cd$ with identity $O$, defined
        by the chord--tangent construction.
        Using \ref{selmer:it:flex-tangent}, show that for $P,Q,R \in \Cd(\kbar)$,
        \[
          P \oplus Q \oplus R = O
          \iff
          P + Q + R \text{ is cut out on } \Cd \text{ by a line.}
        \]
\begin{solution}
% Write your solution here.
\end{solution}
\end{enumerate}
\end{problem}

\section{The group law}
\label{selmer:sec:grouplaw}

\begin{problem}[Negation, and $2$-torsion]
\label{selmer:prob:neg}
\begin{enumerate}
  \item\label{selmer:it:neg-formula} Show that $\ominus(X_0:Y_0:Z_0) = (Y_0:X_0:Z_0)$ for every point
        $(X_0:Y_0:Z_0) \in \Cd(\kbar)$. There are two things to check:
        that the right-hand side lies on $\Cd$, and that it is the third
        intersection point of $\Cd$ with the line through $O$ and $(X_0:Y_0:Z_0)$.
\begin{solution}
% Write your solution here.
\end{solution}
  \item\label{selmer:it:neg-2torsion} Deduce that the nontrivial $2$-torsion of $\Cd$ is cut out by the line
        $X = Y$, and that $\#\Cd[2](\kbar) = 4$ when $\Char k \neq 2$.
\begin{solution}
% Write your solution here.
\end{solution}
  \item\label{selmer:it:neg-rational-2torsion} Take $k = \Q$. For which $d$ does $\Cd$ have a nontrivial rational
        $2$-torsion point? Write it down, and give the smallest positive
        integer $d$ for which this happens.
\begin{solution}
% Write your solution here.
\end{solution}
\end{enumerate}
\end{problem}

Part~\ref{selmer:it:tri-expansion} of \Cref{selmer:prob:trilinear} is its heart. Parts
\ref{selmer:it:tri-chord} and \ref{selmer:it:tri-doubling} are then mechanical. If time is short,
take \ref{selmer:it:tri-expansion} on faith and come back to its proof, but do not skip
\ref{selmer:it:tri-tangent}: it is what makes \ref{selmer:it:tri-doubling} painless. The identity of
\ref{selmer:it:tri-expansion} is the polarization of a cubic form. For a diagonal cubic the
trilinear form can be written down by hand.

\begin{problem}[The trilinear form, and the addition formulas]
\label{selmer:prob:trilinear}
For vectors $P = (X_P,Y_P,Z_P)$, $Q$ and $R$ in $\kbar^3$, define
\begin{equation}
\label{selmer:eq:tri}
  \trilin{P}{Q}{R} \coloneqq X_PX_QX_R + Y_PY_QY_R - d\,Z_PZ_QZ_R .
\end{equation}
\begin{enumerate}
  \item\label{selmer:it:tri-symmetric} Check that \eqref{selmer:eq:tri} is symmetric in $P,Q,R$, linear in each
        argument separately, and satisfies $\trilin{P}{P}{P} = F(P)$.
\begin{solution}
% Write your solution here.
\end{solution}
  \item\label{selmer:it:tri-expansion} Prove the \emph{expansion identity}: for scalars $s,t$,
        \begin{equation}
        \label{selmer:eq:exp}
          F(sP + tQ) = s^3\trilin{P}{P}{P} + 3s^2t\,\trilin{P}{P}{Q}
                     + 3st^2\,\trilin{P}{Q}{Q} + t^3\trilin{Q}{Q}{Q}.
        \end{equation}
        Where do the $3$'s come from?\hint{expand
        $\trilin{sP+tQ}{sP+tQ}{sP+tQ}$ by trilinearity into eight terms and
        collect using symmetry.}
\begin{solution}
% Write your solution here.
\end{solution}
  \item\label{selmer:it:tri-tangent} Prove that $3\trilin{P}{P}{Q} = \nabla F(P) \cdot Q$, and deduce that for
        $P \in \Cd(\kbar)$,
        \[
          \trilin{P}{P}{Q} = 0 \iff Q \text{ lies on the tangent line to } \Cd
          \text{ at } P.
        \]
\begin{solution}
% Write your solution here.
\end{solution}
  \item\label{selmer:it:tri-chord} Let $P_1 = (X_1:Y_1:Z_1)$ and $P_2 = (X_2:Y_2:Z_2)$ be distinct points of
        $\Cd(\kbar)$, and parametrize the line through them as $sP_1 + tP_2$. Using \ref{selmer:it:tri-expansion}, show that the third intersection
        point is at $(s:t) = (B:-A)$ where
        $A = \trilin{P_1}{P_1}{P_2}$ and $B = \trilin{P_1}{P_2}{P_2}$, and hence that
        \begin{equation}
        \label{selmer:eq:chord}
          P_1 \oplus P_2 = (BY_1 - AY_2 : BX_1 - AX_2 : BZ_1 - AZ_2).
        \end{equation}
        Why can $A$ and $B$ not both vanish?
\begin{solution}
% Write your solution here.
\end{solution}
  \item\label{selmer:it:tri-doubling} Let $P = (X_0:Y_0:Z_0) \in \Cd(\kbar)$ and put $Q \coloneqq (Y_0^2:-X_0^2:0)$. Verify that $Q$
        lies on the tangent to $\Cd$ at $P$, and use \ref{selmer:it:tri-expansion} and \ref{selmer:it:tri-tangent} to derive
        \begin{equation}
        \label{selmer:eq:doubling}
          [2]P = \bigl(-Y_0(2X_0^3+Y_0^3) : X_0(X_0^3+2Y_0^3) : Z_0(X_0^3-Y_0^3)\bigr).
        \end{equation}
\begin{solution}
% Write your solution here.
\end{solution}
  \item\label{selmer:it:tri-check} Check \eqref{selmer:eq:doubling} against \Cref{selmer:prob:neg}: what does it give on
        the line $X = Y$? (The line $X = 0$ is taken up in \Cref{selmer:prob:flexes}.)
\begin{solution}
% Write your solution here.
\end{solution}
  \item\label{selmer:it:tri-char3} Where does the whole method break down if $\Char k = 3$? Compare your
        answer with \Cref{selmer:prob:nice}\ref{selmer:it:nice-char3}.
\begin{solution}
% Write your solution here.
\end{solution}
\end{enumerate}
\end{problem}

\begin{problem}[Arithmetic on $C_7$]
\label{selmer:prob:d7}
Let $d = 7$, so that $P \coloneqq (2:-1:1) \in C_7(\Q)$ because $2^3 + (-1)^3 = 7$.
\begin{enumerate}
  \item\label{selmer:it:d7-double} Compute $[2]P$ using \eqref{selmer:eq:doubling}, and simplify. Verify your
        answer by exhibiting an identity between cubes of integers.
\begin{solution}
% Write your solution here.
\end{solution}
  \item\label{selmer:it:d7-triple} Compute $[3]P$ using \eqref{selmer:eq:chord} applied to $P$ and $[2]P$, and
        verify it the same way.
\begin{solution}
% Write your solution here.
\end{solution}
  \item\label{selmer:it:d7-quadruple} Compute $[4]P$.
\begin{solution}
% Write your solution here.
\end{solution}
  \item\label{selmer:it:d7-order} The three rational points you have found have numerators and denominators of
        rapidly increasing size. What does this suggest about the order of $P$
        in $C_7(\Q)$? (You are not asked to prove it.)
\begin{solution}
% Write your solution here.
\end{solution}
\end{enumerate}
\end{problem}

In \Cref{selmer:prob:flexes}, parts \ref{selmer:it:flexes-hessian-count}--\ref{selmer:it:flexes-exist} are the
general theory and \ref{selmer:it:flexes-selmer}--\ref{selmer:it:flexes-rational-3torsion} the payoff for
our family. If time is short, take \ref{selmer:it:flexes-exist} on faith and go to
\ref{selmer:it:flexes-selmer}--\ref{selmer:it:flexes-rational-3torsion}. Part \ref{selmer:it:flexes-nine} is the
most satisfying: two completely different counts turn out to
agree. Part \ref{selmer:it:flexes-hessian-count} uses B\'ezout's theorem
\cite{Hartshorne77}*{\S I.7}.

\begin{wsproposition}[The Hessian criterion]
\label{selmer:prop:hessian}
Let $C = V(F) \subset \PP^2_{\kbar}$ be a smooth plane curve of degree $m$ over
a field of characteristic $0$ or $> 3$, and let
\[
  H(F) = \det\!\paren{\frac{\del^2 F}{\del X_i \, \del X_j}}_{0 \leq i,j \leq 2}.
\]
Then a point $P \in C(\kbar)$ is an inflection point if and only if $H(F)(P) = 0$.
This is classical. Characteristic three is taken up in
\Cref{selmer:prob:flexes}\ref{selmer:it:flexes-char3} below.
\end{wsproposition}

\begin{problem}[Inflection points and $3$-torsion]
\label{selmer:prob:flexes}
\begin{enumerate}
  \item\label{selmer:it:flexes-hessian-count} Let $C = V(F)$ be as in \Cref{selmer:prop:hessian}.
        Show that $H(F)$ is a form of degree $3(m-2)$, so of degree $3$ when
        $m = 3$. Assuming $C$ is a smooth plane cubic, explain why
        $C \not\subseteq V(H(F))$, and conclude from B\'ezout's
        theorem that $C$ has exactly $9$ inflection
        points \emph{counted with multiplicity}.
\begin{solution}
% Write your solution here.
\end{solution}
  \item\label{selmer:it:flexes-3torsion} Now forget the Hessian. Let $C$ be a smooth plane cubic over
        $\kbar$ and suppose $O \in C(\kbar)$ is a flex, so that $(C,O)$ is an
        elliptic curve with the group law of \Cref{selmer:prob:flex}\ref{selmer:it:flex-grouplaw}. Show that
        \[
          P \text{ is a flex} \iff [3]P = O ,
        \]
        so that the set of flexes is exactly $C[3](\kbar)$.
\begin{solution}
% Write your solution here.
\end{solution}
  \item\label{selmer:it:flexes-exist} Show that a smooth plane cubic over $\kbar$ always \emph{has} a
        flex, without using the Hessian.\hint{for any $P_0 \in C(\kbar)$
        the class $\mcO_C(1) - 3[P_0]$ has degree $0$, and multiplication by
        $3$ is surjective on $\Pic^0 C$.} Deduce that, with no flex singled
        out, the set of flexes is a torsor under $\Pic^0(C)[3]$. In particular
        the number of flexes is $\#C[3](\kbar)$.
\begin{solution}
% Write your solution here.
\end{solution}
  \item\label{selmer:it:flexes-nine} Combine \ref{selmer:it:flexes-3torsion} and \ref{selmer:it:flexes-exist} with the structure of $C[3]$ in characteristic
        $\neq 3$ to conclude that a smooth plane cubic has exactly $9$ flexes.
        What does comparing this with \ref{selmer:it:flexes-hessian-count} tell you about the intersection
        $C \cap V(H(F))$?
\begin{solution}
% Write your solution here.
\end{solution}
  \item\label{selmer:it:flexes-selmer} Specialize to $\Cd$. Compute $H(F)$ for $F = X^3 + Y^3 - dZ^3$ and
        deduce that, as subsets of $\Cd(\kbar)$,
        \[
          \Cd[3](\kbar) = \paren{\Cd \cap \brk{XYZ = 0}}(\kbar).
        \]
        Write down all nine flexes.
\begin{solution}
% Write your solution here.
\end{solution}
  \item\label{selmer:it:flexes-doubling} Check \ref{selmer:it:flexes-selmer} against the group law directly: use \eqref{selmer:eq:doubling} to
        show that every point $P = (0:Y_0:Z_0) \in \Cd(\kbar)$ satisfies $[2]P = \ominus P$.
\begin{solution}
% Write your solution here.
\end{solution}
  \item\label{selmer:it:flexes-rational-3torsion} Which of the nine flexes lie in $\Cd(k)$? Treat the three lines
        separately, and say what happens for $k = \Q$.
\begin{solution}
% Write your solution here.
\end{solution}
  \item\label{selmer:it:flexes-char3} ($\star$) In characteristic three, $\#C[3](\kbar)$ is $3$ or $1$
        according as $C$ is ordinary or supersingular, so a smooth plane cubic
        has $3$ or $1$ flexes. How is that compatible with \ref{selmer:it:flexes-hessian-count}?
\begin{solution}
% Write your solution here.
\end{solution}
\end{enumerate}
\end{problem}

In \Cref{selmer:prob:cm} you may use the following two results. Part
\ref{selmer:it:cm-predict} asks what the classification of automorphism groups in
\cite{Silverman09}*{\S III.10} predicts.

\begin{wstheorem}[{Morphisms fixing $O$ \cite{Silverman09}*{Theorem~III.4.8}}]
\label{selmer:thm:hom}
A morphism of elliptic curves taking $O$ to $O$ is a group homomorphism.
\end{wstheorem}

\begin{wsproposition}[{Endomorphism rings \cite{Silverman09}*{\S III.4}}]
\label{selmer:prop:end}
The endomorphism ring of an elliptic curve has no zero divisors.
\end{wsproposition}

\begin{problem}[Complex multiplication, $\star$]
\label{selmer:prob:cm}
Let $\alpha(X:Y:Z) = (X:Y:\zeta Z)$, a map defined over $k(\zeta)$.
\begin{enumerate}
  \item\label{selmer:it:cm-alpha} Show that $\alpha$ maps $\Cd$ to itself and fixes $O$.
\begin{solution}
% Write your solution here.
\end{solution}
  \item\label{selmer:it:cm-automorphism} Conclude that $\alpha$ is a
        group homomorphism, hence an automorphism of the elliptic curve
        $(\Cd, O)$ over $k(\zeta)$.
\begin{solution}
% Write your solution here.
\end{solution}
  \item\label{selmer:it:cm-relation} Show that $\alpha^2 + \alpha + 1 = 0$ in
        $\End(\Cd \times_k k(\zeta))$, the ring of endomorphisms of $\Cd$ defined
        over $k(\zeta)$.\hint{$\alpha^3 = 1$, and the endomorphism ring of
        an elliptic curve has no zero divisors.}
\begin{solution}
% Write your solution here.
\end{solution}
  \item\label{selmer:it:cm-predict} Deduce that $\Z[\zeta] \subseteq \End(\Cd \times_k k(\zeta))$,
        so every Selmer curve has complex multiplication. What does
        the classification of automorphism groups then predict for $j(\Cd)$ and for
        $\#\Aut(\Cd)$ and $\#\Aut(\Cd \times_k k(\zeta))$?
\begin{solution}
% Write your solution here.
\end{solution}
\end{enumerate}
\end{problem}

\section{From Riemann--Roch to a Weierstrass equation}
\label{selmer:sec:rr}

\noindent
This section reproduces, for our family, the proof of
\cite{Silverman09}*{Proposition~III.3.1}: every elliptic curve is isomorphic to a
curve in Weierstrass form.

\Cref{selmer:prob:rr} is the longest problem here. Parts \ref{selmer:it:rr-z}--\ref{selmer:it:rr-bases} go quickly once
you see that \Cref{selmer:prop:bezout} converts ``pole order at $O$'' into
``multiplicity of $O$ in $h \cap \Cd$''. Part \ref{selmer:it:rr-relation} is a two-line computation with
a surprising answer, and \ref{selmer:it:rr-automorphisms} is worth coming back to.
Weierstrass equations and their coefficients $a_1, \dots, a_6$ are as in
\cite{Silverman09}*{\S III.1}.

\begin{wsdefinition}[Riemann--Roch spaces]
\label{selmer:def:rr}
For $D \in \Div \Cd$,
\[
  \mcL(D) = \brk{f \in k(\Cd)^\times : \div(f) \geq -D} \cup \brk{0},
  \qquad \div(f) \in \Div \Cd .
\]
This is a space of functions defined \emph{over $k$}.
\end{wsdefinition}

You may use the following two results without proof.

\begin{wsproposition}[B\'ezout]
\label{selmer:prop:bezout}
For a nonzero linear form $h$, the scheme-theoretic intersection $h \cap \Cd$
is an effective divisor of degree $3$ in $\Div \Cd$. Over $\kbar$ it becomes a
sum of three points of $\Cd(\kbar)$ counted with multiplicity. Consequently,
since $X+Y$ cuts out $3O$ by \Cref{selmer:prob:flex}\ref{selmer:it:flex-tangent},
\[
  \div\!\paren{\frac{h}{X+Y}} = (h \cap \Cd) - 3O .
\]
\end{wsproposition}

\begin{wsproposition}[{Riemann--Roch for genus one \cite{Silverman09}*{\S II.5 and proof of Proposition~III.3.1}}]
\label{selmer:prop:rr}
For every $n \geq 1$ we have $\dim_k \mcL(nO) = n$. Since $\Cd$ is
geometrically integral, $\mcL(D) \tensor_k \kbar = \mcL(D_{\kbar})$, so the
dimension is the same whether computed over $k$ or over $\kbar$.
\end{wsproposition}

\begin{problem}[The Weierstrass model of $\Cd$]
\label{selmer:prob:rr}
Consider the rational functions in $k(\Cd)$
\[
  w \coloneqq \frac{Z}{X+Y}, \quad v \coloneqq \frac{X-Y}{X+Y}.
\]
For parts \ref{selmer:it:rr-t}--\ref{selmer:it:rr-weierstrass} and \ref{selmer:it:rr-automorphisms} assume
$\Char k \neq 2$.
Part~\ref{selmer:it:rr-char2} replaces $v$ by $X/(X+Y) = (1+v)/2$, which also makes sense in
characteristic $2$, and reaches a model valid whenever $\Char k \neq 3$.
\begin{enumerate}
  \item\label{selmer:it:rr-z} Using \Cref{selmer:prob:flex}\ref{selmer:it:flex-infinity} and \Cref{selmer:prop:bezout}, compute
        $\div(w)$. Deduce that $w$ has a pole of order exactly $2$ at $O$ and
        no other pole, so $w \in \mcL(2O)$. Where are its zeros, and what are
        they in the language of \Cref{selmer:prob:flexes}?
\begin{solution}
% Write your solution here.
\end{solution}
  \item\label{selmer:it:rr-t} Show that the line $X - Y = 0$ does not pass through $O$. Deduce that
        $v$ has a pole of order exactly $3$ at $O$, so $v \in \mcL(3O)$. Where
        are its zeros?
\begin{solution}
% Write your solution here.
\end{solution}
  \item\label{selmer:it:rr-bases} Using \Cref{selmer:prop:rr}, explain why $\brk{1, w}$ is a basis of
        $\mcL(2O)$ and $\brk{1,w,v}$ is a basis of $\mcL(3O)$.
\begin{solution}
% Write your solution here.
\end{solution}
  \item\label{selmer:it:rr-seven} The seven functions $1, w, v, w^2, wv, w^3, v^2$ all lie in
        $\mcL(6O)$. Why must they satisfy a nontrivial linear relation over
        $k$? List their pole orders at $O$. Which two collide, and why does
        that force the relation to involve both?
\begin{solution}
% Write your solution here.
\end{solution}
  \item\label{selmer:it:rr-relation} Find the relation.\hint{work on the affine chart $X + Y = 1$,
        where $X = (1+v)/2$ and $Y = (1-v)/2$, and expand $X^3 + Y^3$.}
\begin{solution}
% Write your solution here.
\end{solution}
  \item\label{selmer:it:rr-weierstrass} Let $E_1\colon y_1^2 = x_1^3 - 432d^2$. Show
        that the relation of \ref{selmer:it:rr-relation} says that there is a morphism
        $\varphi_1\colon \Cd \to E_1$ with $\varphi_1^*x_1 = 12dw$ and
        $\varphi_1^*y_1 = 36dv$. Write down $\varphi_1$ and its inverse explicitly.
\begin{solution}
% Write your solution here.
\end{solution}
  \item\label{selmer:it:rr-char2} Where exactly did you use $\Char k \neq 2$? Show that there is instead a
        morphism $\varphi_2$ with $\varphi_2^*x_2 = 3dw$ and $\varphi_2^*y_2 = 9dX/(X+Y)$ to
        \begin{equation}
        \label{selmer:eq:char2-model}
          E_2\colon y_2^2 - 9d\,y_2 = x_2^3 - 27d^2 ,
        \end{equation}
        which is a Weierstrass equation valid whenever $\Char k \neq 3$. Which of
        $a_1, \dots, a_6$ are nonzero?
\begin{solution}
% Write your solution here.
\end{solution}
  \item\label{selmer:it:rr-automorphisms} ($\star$) The relation in \ref{selmer:it:rr-relation} has no $w^2$, $wv$, $w$ or $v$ term,
        although a general Weierstrass equation allows all four. Explain
        why, using the automorphisms $\sigma = \ominus$ of \Cref{selmer:prob:neg} and
        $\alpha$ of \Cref{selmer:prob:cm}.\hint{compute $\sigma^*w$,
        $\sigma^*v$, $\alpha^*w$, $\alpha^*v$. Then use \ref{selmer:it:rr-seven} to argue that the
        relation spans a one-dimensional space, so it is a simultaneous
        eigenvector.}
\begin{solution}
% Write your solution here.
\end{solution}
\end{enumerate}
\end{problem}

\section{Worked examples in \Magma}
\label{selmer:sec:magma}

\noindent
Most of you have never opened \Magma, so this section is written as worked
examples rather than exercises: read them, then re-run them yourself and change
things. Every line of output below is reproduced from \Magma\ V2.28-8
\cite{Magma}, with whitespace tidied and one long type name elided. Three short
variations to try on your own are collected at the end.

\begin{wswarning}[Closed points in \Magma]
\label{selmer:warn:magma}
A \defi{closed point} $P$ of $\Cd$ is a point of the scheme $\Cd$. Its residue
field $k(P)$ is a finite extension of $k$, and $\deg P \coloneqq [k(P):k]$. The
closed points of degree one are the rational points, and a closed point $P$ with
$k(P)/k$ separable corresponds to an orbit of $\deg P$ points of $\Cd(\kbar)$ under
$\Gal(\ksep/k)$. The divisor group $\Div \Cd$ is the free abelian group on the
closed points \cite{Poonen06}*{\S 1.9}, whereas \cite{Silverman09}*{\S II.3}
works with points of $\Cd(\kbar)$ and Galois-invariant divisors on $\Cdbar$.
\Magma\ follows the first convention: \mg{Places} returns closed points,
\mg{Divisor} builds elements of $\Div \Cd$, and \mg{RiemannRochSpace} returns
$\mcL(D)$ as a $k$-vector space.
\end{wswarning}

Each example continues the previous one. Begin with
\begin{magma}
QQ := Rationals();
P2<X,Y,Z> := ProjectiveSpace(QQ,2);
d := 7;
C := Curve(P2, X^3 + Y^3 - d*Z^3);
O := C![1,-1,0];
\end{magma}
The notation \mg{P2<X,Y,Z>} names the coordinates so that we can write
equations in them, and \mg{C!} coerces a list of three numbers into a rational point of
\mg{C}. Typing \mg{O := (1:-1:0);} does not work.

\begin{wsexample}[Setting up, and a trap]
\label{selmer:ex:magma-setup}
First, \Cref{selmer:prob:nice}:
\begin{magma}
IsNonSingular(C);  // true
Genus(C);          // 1
O;                 // (-1 : 1 : 0)
\end{magma}
The last line is worth a pause. We asked for $(1:-1:0)$ and got $(-1:1:0)$: a
point of $\PP^2(\Q)$ is defined only up to scaling, and \Magma\ normalizes the
representative it stores, here by $-1$.

Now ask for the Weierstrass model. The second return value \mg{mp} is the
isomorphism from \mg{C} to \mg{E}.
\begin{magma}
E, mp := EllipticCurve(C, O);
E;
// Elliptic Curve defined by y^2 - 9/49*y = x^3 - 27/2401 over Rational Field
jInvariant(E);                            // 0
E eq EllipticCurve([0,0,0,0,-432*d^2]);   // false
\end{magma}
That \mg{false} looks alarming, since \Cref{selmer:prob:rr}\ref{selmer:it:rr-weierstrass} says the answer should
be $E_1\colon y_1^2 = x_1^3 - 432d^2$. It is not a contradiction: \mg{eq} compares
Weierstrass \emph{coefficients}, and by \cite{Silverman09}*{Proposition~III.3.1(b)} a
Weierstrass model is unique only up to the changes of variables
$x = u^2x' + r$, $y = u^3y' + su^2x' + t$. The right question is
whether the curves are isomorphic, and against which model:
\begin{magma}
Eours := EllipticCurve([0,0,-9*d,0,-27*d^2]);
Eours;                            // y^2 - 63*y = x^3 - 1323
bool, iso := IsIsomorphic(Eours, E);
bool;                             // true
IsomorphismData(iso);             // [ 0, 0, 0, 1/7 ]
\end{magma}
So $r = s = t = 0$ and $u = 1/7 = 1/d$: \Magma\ found exactly the model $E_2$ of
\Cref{selmer:prob:rr}\ref{selmer:it:rr-char2}, with $x_2$ and $y_2$ divided by $d^2$ and $d^3$. The two maps agree as well,
\begin{magma}
DefiningEquations(mp);   // [ 3/7*Z, 9/49*X, X + Y ]
\end{magma}
against our $(3dZ : 9dX : X+Y)$, the ratios being $d^2$ and $d^3$ exactly as
$u = 1/d$ demands. The moral: for elliptic curves, \mg{eq} is almost never the
question you mean to ask.

We can even watch this happen for all $d$ at once, by working over the function
field $\Q(d)$ instead of over $\Q$:
\begin{magma}
Fd<D> := FunctionField(Rationals());
P2d<Xd,Yd,Zd> := ProjectiveSpace(Fd,2);
Cd := Curve(P2d, Xd^3 + Yd^3 - D*Zd^3);
Ed, mpd := EllipticCurve(Cd, Cd![1,-1,0]);
Ed;
// Elliptic Curve defined by y^2 - 9/D^2*y = x^3 - 27/D^4
DefiningEquations(mpd);
// [ 3/D*Zd, 9/D^2*Xd, Xd + Yd ]
\end{magma}
which is \eqref{selmer:eq:char2-model} rescaled by $u = d$, as an identity in $d$.

Finally, the predictions of \Cref{selmer:prob:cm}\ref{selmer:it:cm-predict}:
\begin{magma}
#AutomorphismGroup(E);                     // 2
K<zeta> := CyclotomicField(3);
#AutomorphismGroup(BaseChange(E,K));       // 6
\end{magma}
\end{wsexample}

\begin{wsexample}[The addition formulas, checked]
\label{selmer:ex:magma-group}
A \Magma\ \mg{function} takes arguments and returns a value. Here both are
sequences of three coordinates, so that \eqref{selmer:eq:chord} and
\eqref{selmer:eq:doubling} can be typed almost verbatim.
\begin{magma}
dbl := function(P)
  x := P[1]; y := P[2]; z := P[3];
  return [-y*(2*x^3+y^3), x*(x^3+2*y^3), z*(x^3-y^3)];
end function;

add := function(P1,P2)
  A := P1[1]^2*P2[1] + P1[2]^2*P2[2] - d*P1[3]^2*P2[3];
  B := P1[1]*P2[1]^2 + P1[2]*P2[2]^2 - d*P1[3]*P2[3]^2;
  return [B*P1[2]-A*P2[2], B*P1[1]-A*P2[1], B*P1[3]-A*P2[3]];
end function;

P := [2,-1,1];
Q := dbl(P);      C!Q;   //  (5/3 : 4/3 : 1)
R := dbl(Q);      C!R;   //  (-1256/183 : 1265/183 : 1)
S := add(P,Q);    C!S;   //  (-17/38 : 73/38 : 1)
\end{magma}
These are the rational points of \Cref{selmer:prob:d7}. Coercing through \mg{C!} normalizes
them, which is why $(15:12:9)$ prints as $(5/3 : 4/3 : 1)$. Now transport
everything to $E$ through \mg{mp} and compare with \Magma's own group law:
\begin{magma}
EP := mp(C!P);
mp(C!Q) eq 2*EP;         // true
mp(C!S) eq 3*EP;         // true
mp(C!R) eq 4*EP;         // true
\end{magma}
Three \mg{true}s: formulas derived from nothing but \eqref{selmer:eq:tri} and a
$3 \times 3$ determinant agree with a general-purpose implementation.
\end{wsexample}

\begin{wsexample}[Riemann--Roch, mechanically]
\label{selmer:ex:magma-rr}
\Magma\ can carry out \Cref{selmer:prob:rr} without being told the answer. The
generators \mg{x} and \mg{y} of \mg{FunctionField(C)} are the affine coordinates
$X/Z$ and $Y/Z$ of $\Cd$ (not Weierstrass coordinates), so our two functions are
built from those.
\begin{magma}
FF<x,y> := FunctionField(C);   // x = X/Z,  y = Y/Z
w := 1/(x+y);                  // Z/(X+Y)
v := (x-y)/(x+y);              // (X-Y)/(X+Y)
plO := Place(O);
Valuation(w, plO);             // -2
Valuation(v, plO);             // -3
\end{magma}
Those are the pole orders of \Cref{selmer:prob:rr}\ref{selmer:it:rr-z}, \ref{selmer:it:rr-t}. A valuation of $-2$ is a
pole of order $2$. As in \Cref{selmer:warn:magma}, \mg{Place(O)} is the closed point $O$ and
\mg{Divisor} builds an element of $\Div \Cd$, so everything here happens over
$\Q$. Next,
\Cref{selmer:prop:rr} itself:
\begin{magma}
D := Divisor(plO);
for n in [1..6] do print n, Dimension(RiemannRochSpace(n*D)); end for;
// 1 1
// 2 2
// 3 3
// 4 4
// 5 5
// 6 6
\end{magma}
Now ask for a basis of $\mcL(2O)$ without suggesting one:
\begin{magma}
V2, m2 := RiemannRochSpace(2*D);
[m2(b) : b in Basis(V2)];      // [ x^2 - y*x + y^2, 1 ]
\end{magma}
which looks nothing like our $w$. But $X^3+Y^3 = (X+Y)(X^2-XY+Y^2)$, and on
$\Cd$ the left side is $dZ^3$, so
\[
  \frac{X^2-XY+Y^2}{Z^2} = \frac{dZ^3}{Z^2(X+Y)} = d\,\frac{Z}{X+Y} = d\,w ,
\]
and indeed
\begin{magma}
(x^2 - x*y + y^2) eq d*w;      // true
\end{magma}
Finally the relation. We list the seven functions of \Cref{selmer:prob:rr}\ref{selmer:it:rr-seven}, express
each in a basis of $\mcL(6O)$ (that is what \mg{@@} does, applying the
inverse of the map \mg{mV}), and take the kernel of the resulting matrix.
\begin{magma}
V, mV := RiemannRochSpace(6*D);
fns := [FF!1, w, v, w^2, w*v, w^3, v^2];
[-Valuation(f, plO) : f in fns];   // [ 0, 2, 3, 4, 5, 6, 6 ]
M := Matrix([Eltseq(f @@ mV) : f in fns]);
K := Kernel(M);
Dimension(K);                      // 1
Basis(K);                          // [ ( 1  0  0  0  0  -28  3) ]
\end{magma}
The pole orders are $0,2,3,4,5,6,6$: seven functions, six available orders, the
collision at $6$ forcing the relation exactly as in \Cref{selmer:prob:rr}\ref{selmer:it:rr-seven}. And the
kernel is one-dimensional, spanned by the vector reading
\[
  1 - 28w^3 + 3v^2 = 0,
\]
which is $3v^2 + 1 = 4dw^3$ at $d = 7$. We got it by hand from a factorization.
\Magma\ got it from linear algebra in a Riemann--Roch space.
\end{wsexample}

\subsection*{Now try these}

\noindent
Each one is a small edit to \Cref{selmer:ex:magma-setup} or \Cref{selmer:ex:magma-group}, and
in each case the worksheet already tells you what the answer must be.

\begin{enumerate}[label=(\arabic*), ref=(\arabic*)]
  \item\label{selmer:it:try-d2} Set \mg{d := 2}. \Cref{selmer:prob:neg}\ref{selmer:it:neg-rational-2torsion} predicts a rational point of order
        two. Find it, verify with the group law that it has order two, and
        compute its image under $\varphi_1$ by hand.
\begin{solution}
% Write your solution here.
\end{solution}
  \item\label{selmer:it:try-d8} Set \mg{d := 8}. \Cref{selmer:prob:flexes}\ref{selmer:it:flexes-rational-3torsion} predicts rational $3$-torsion.
        Find the two rational $3$-torsion points other than $O$, use your \mg{dbl} function to check
        that $[2]P = \ominus P$, and compute $\varphi_1(P)$ by hand.
\begin{solution}
% Write your solution here.
\end{solution}
  \item\label{selmer:it:try-isomorphic} Use \mg{IsIsomorphic} to decide whether $C_1 \isom C_8$ and whether
        $C_1 \isom C_2$ over $\Q$. Pass it the \emph{Weierstrass} models: on
        the plane curves themselves \Magma\ answers \mg{For genus 0 or 1
        curves the basefield must be finite}. Reconcile the answers with the
        cubic-twist criterion of \Cref{selmer:sec:name}.
\begin{solution}
% Write your solution here.
\end{solution}
\end{enumerate}

\section{Where the name comes from}
\label{selmer:sec:name}

\noindent
Selmer \cite{Selmer-51} studied the general diagonal cubic
$aX^3 + bY^3 + cZ^3 = 0$, and produced the standard counterexample to the Hasse
principle: the curve $3X^3 + 4Y^3 + 5Z^3 = 0$ has real points and
$\Qp$-points for every prime $p$, but no rational point. Such a curve is a torsor under its Jacobian
\cite{Poonen2017}*{\S 5.12}, and that Jacobian is a Selmer curve in our sense:
\[
  \Jac(aX^3 + bY^3 + cZ^3) \isom \bigl(y^2 = x^3 - 432(abc)^2\bigr) .
\]
Two ways to test this in \Magma: the curve $X^3 + 2Y^3 + 3Z^3 = 0$ \emph{does}
have the rational point $(1:1:-1)$, so it is its own Jacobian and should be
isomorphic to $C_6$, that is to $y^2 = x^3 - 432\cdot 6^2$; and for Selmer's own
curve one can call \mg{Jacobian(GenusOneModel(3*X^3 + 4*Y^3 + 5*Z^3))} and
compare with $y^2 = x^3 - 432\cdot 60^2$. Both come out \mg{true}.

Finally, by \Cref{selmer:prob:rr}\ref{selmer:it:rr-weierstrass} together with the classification of
twists of $y^2 = x^3 + A$ \cite{Silverman09}*{\S X.5}, for $\Char k \neq 2, 3$ the
curves $\Cd$ and $C_{d'}$ are isomorphic over $k$ exactly when
$d/d' \in (k^\times)^3$, so the family is parametrized by
$\umtp{k}$. These are the \defi{cubic twists} of $X^3 + Y^3 = Z^3$, and they
are the natural first examples for the descent machinery of
\cite{Silverman09}*{\S X} and \cite{Poonen2017}*{Ch.~8}.

\bigskip

\part[The Weil pairing, endomorphisms and automorphisms]{The Weil pairing, endomorphisms and automorphisms}
\label{weil-end-aut:part}
\noindent
\Cref{weil-end-aut:sec:weil} shows that the determinant of the mod~$N$ Galois representation
is the cyclotomic character. \Cref{weil-end-aut:sec:end} exhibits an endomorphism ring that
is an order in a definite quaternion algebra, the third case of the
classification \cite{Silverman09}*{Theorem~III.9.3}, on a single curve over
$\F_9$. \Cref{weil-end-aut:sec:aut} computes the automorphism groups of order $12$ and $24$
in characteristic $3$ and $2$, the two cases of
\cite{Silverman09}*{Theorem~III.10.1} in which the group is not cyclic.

\subsection*{Reading}
\cite{Silverman09}*{\S\S III.8--III.10}, with \cite{Silverman09}*{Appendix~A.1}
for the automorphism groups in characteristic $2$ and $3$. For the quaternion
background: \cite{Voight21}*{Definition~2.2.1} for the algebra
$(a, b \mid \Q)$, \cite{Voight21}*{\S\S 3.2--3.3} for the standard involution, the
reduced trace and the reduced norm, \cite{Voight21}*{Main Theorem~5.4.4 and \S
  14.1} for split and definite algebras, and
\cite{Voight21}*{Definitions~10.2.1 and 10.4.1} for orders and maximal orders.
\Cref{weil-end-aut:prop:ram}, \cref{weil-end-aut:thm:maximal} and \cref{weil-end-aut:prop:index} are the results we actually use.

\subsection*{\texorpdfstring{\Magma}{Magma}}
If you have no local installation, every block below runs in the free
\href{http://magma.maths.usyd.edu.au/calc/}{online calculator}, which caps how long a
single job may take.

\subsection*{Standing notation}
\begin{itemize}
  \item A field $k$, with an algebraic closure $\kbar$ and the separable closure
        $\ksep \subseteq \kbar$.
  \item Multiplication by $m$, written $[m]$, and the dual $\hat\varphi$ of an isogeny
        $\varphi$.
  \item The $q$-power Frobenius endomorphism $\Frob \in \End(E)$ of a curve over $\F_q$.
        Beware that $\varphi_\ell$ already means the map induced by $\varphi$ on
        $T_\ell(E)$, so the Frobenius \emph{endomorphism} is $\Frob$ and never
        $\varphi_q$.
\end{itemize}

\section{The Weil pairing and the cyclotomic character}
\label{weil-end-aut:sec:weil}

Fix $N \geq 2$ with $\Char k \nmid N$, so that $E[N](\ksep) \cong (\Z/N\Z)^2$.

\begin{wsdefinition}[The Weil pairing]
\label{weil-end-aut:def:weil}
We write
\[
  e_N\colon E[N](\ksep) \times E[N](\ksep) \to \mu_N(\ksep)
\]
for the \defi{Weil pairing}.
\end{wsdefinition}

\begin{wsdefinition}[Galois representations]
\label{weil-end-aut:def:rep}
We write
\[
  \bar\rho_{E,N}\colon \Gal_k \to \Aut\bigl(E[N](\ksep)\bigr) \cong \GL_2(\Z/N\Z)
\]
for the \defi{mod~$N$ Galois representation}, the isomorphism depending on a choice of
basis, and $\rho_{E,\ell}$ for the \defi{$\ell$-adic representation} on $T_\ell(E)$.
\end{wsdefinition}

\begin{wsproposition}[{Properties of the Weil pairing \cite{Silverman09}*{Propositions~III.8.1 and III.8.2}}]
\label{weil-end-aut:prop:weil}
The Weil pairing $e_N$ is
\begin{enumerate}
  \item\label{weil-end-aut:it:weil-bilinear} bilinear;
  \item\label{weil-end-aut:it:weil-alt} alternating: $e_N(T, T) = 1$, so $e_N(S, T) = e_N(T, S)^{-1}$;
  \item\label{weil-end-aut:it:weil-nondeg} non-degenerate: if $e_N(S, T) = 1$ for all $S \in E[N](\ksep)$, then $T = O$;
  \item\label{weil-end-aut:it:weil-galois} Galois equivariant: $\sigma\bigl(e_N(S, T)\bigr) = e_N(\sigma S, \sigma T)$ for $\sigma \in \Gal_k$;
  \item\label{weil-end-aut:it:weil-compat} compatible: $e_{NN'}(S, T) = e_N([N']S, T)$;
  \item\label{weil-end-aut:it:weil-adjoint} adjoint for isogenies: $e_N(\varphi S, T) = e_N(S, \hat\varphi T)$.
\end{enumerate}
These properties are.
\end{wsproposition}

\begin{wsproposition}[{The determinant of $\varphi_\ell$ \cite{Silverman09}*{Proposition~III.8.6}}]
\label{weil-end-aut:prop:det}
Let $\varphi \in \End(E)$ and let $\ell \neq \Char k$ be a prime. Then
$\det\varphi_\ell = \deg\varphi$.
\end{wsproposition}

\begin{problem}[The determinant is the cyclotomic character]
\label{weil-end-aut:prob:weil-det}
\begin{enumerate}
  \item\label{weil-end-aut:it:weil-chi} Let $\zeta \in \mu_N(\ksep)$ be primitive. Show that for each
        $\sigma \in \Gal_k$ there is a unique $\chi_N(\sigma) \in \Z/N\Z$ with
        $\sigma(\zeta) = \zeta^{\chi_N(\sigma)}$, that $\chi_N(\sigma)$ is a unit,
        and that the resulting map
        $\chi_N\colon \Gal_k \to (\Z/N\Z)^\times$ is a homomorphism independent
        of the choice of $\zeta$. This is the \defi{mod $N$ cyclotomic character}.
\begin{solution}
% Write your solution here.
\end{solution}
  \item\label{weil-end-aut:it:weil-primitive} Let $S, T$ be a basis of $E[N](\ksep)$ and put
        $\zeta \coloneqq e_N(S, T)$. Show that $\zeta$ is a \emph{primitive}
        $N$-th root of unity, using only
        parts \ref{weil-end-aut:it:weil-bilinear}--\ref{weil-end-aut:it:weil-nondeg} of \cref{weil-end-aut:prop:weil}.
\begin{solution}
% Write your solution here.
\end{solution}
  \item\label{weil-end-aut:it:weil-det} Write $\sigma S = aS + cT$ and $\sigma T = bS + dT$, so that
        $\bar\rho_{E,N}(\sigma) = \smallmat{a}{b}{c}{d}$ in the basis $S, T$. Expand
        $e_N(\sigma S, \sigma T)$ and deduce
        \[
          \det \bar\rho_{E,N} = \chi_N .
        \]
\begin{solution}
% Write your solution here.
\end{solution}
  \item\label{weil-end-aut:it:weil-odd} Take $k = \Q$. Show that $\det\bar\rho_{E,N}$ is surjective for
        \emph{every} elliptic curve $E$ over $\Q$, that
        $\Q(\zeta_N) \subseteq \Q\bigl(E[N]\bigr)$, and that
        $\det\bar\rho_{E,N}(c) = -1$ for $c$ a complex conjugation: the
        representation is \defi{odd}.
\begin{solution}
% Write your solution here.
\end{solution}
  \item\label{weil-end-aut:it:weil-frob} Take $k = \F_q$ with $q$ prime to $N$. Show that
        $\det\bar\rho_{E,N}(\Frob) = q$, and check this against the identity
        $\det\varphi_\ell = \deg\varphi$.
\begin{solution}
% Write your solution here.
\end{solution}
  \item\label{weil-end-aut:it:weil-magma} Check \ref{weil-end-aut:it:weil-frob} numerically in \Magma, taking
        $N = 3$, $q = 5$ and $E \colon y^2 = x^3 + 1$ over $\F_5$.
\begin{magma}
q := 5;  N := 3;
E := EllipticCurve([GF(q) | 0, 1]);      // y^2 = x^3 + 1
#E;                                      // 6
Ek := ChangeRing(E, GF(q^2));
#Ek, AbelianGroup(Ek);                   // 36, Z/6 + Z/6
T3 := [P : P in Points(Ek) | N*P eq Ek!0];
#T3;                                     // 9
S := [P : P in T3 | P ne Ek!0][1];
T := [P : P in T3 | P ne Ek!0 and P ne S and P ne -S][1];
zeta := WeilPairing(S, T, N);  Order(zeta);    // 3: a primitive cube root
Frob := func<P | P eq Ek!0 select Ek!0 else Ek![P[1]^q, P[2]^q]>;
WeilPairing(Frob(S), Frob(T), N) eq zeta^q;    // true
\end{magma}
        Two lines carry the mathematics. That \mg{Order(zeta)} is $3$ is
        \ref{weil-end-aut:it:weil-primitive}. That the last line is \mg{true} is
        $\det\bar\rho_{E,3}(\Frob) = q$, which is \ref{weil-end-aut:it:weil-frob} on this curve. Note
        that \mg{Frob} raises both coordinates to the $q$-th power: that is the $q$-power
        Frobenius, an endomorphism because $E$ is defined over $\F_q$. Now write the
        matrix of $\Frob$ in the basis
        $S, T$ and read its determinant off:
\begin{magma}
[ <a, b> : a, b in [0..2] | a*S + b*T eq Frob(S) ];   // [ <1, 0> ]
[ <a, b> : a, b in [0..2] | a*S + b*T eq Frob(T) ];   // [ <0, 2> ]
GF(3)!q;                                              // 2
\end{magma}
        Check that a different choice of basis moves the matrix but not the
        determinant.
\begin{solution}
% Write your solution here.
\end{solution}
\end{enumerate}
\end{problem}

\begin{wsremark}
\label{weil-end-aut:rem:weil-reproof}
Part~\ref{weil-end-aut:it:weil-det} of \cref{weil-end-aut:prob:weil-det} gives back
\cite{Silverman09}*{Corollary~III.8.1.1}:
if $E[N](\ksep) \subseteq E(k)$, then $\bar\rho_{E,N}$ is trivial, so $\chi_N$ is
trivial, so every $\sigma \in \Gal_k$ fixes a primitive $N$-th root of unity, and
$\mu_N(\ksep) \subseteq k$.
\end{wsremark}

A word, for the curious, on the example in \cref{weil-end-aut:prob:weil-det}\ref{weil-end-aut:it:weil-magma}:
choosing it was half of
that part. We want $q \not\equiv 1 \pmod N$, so that $\chi_N$ is \emph{not} trivial and
the check has content. By \cref{weil-end-aut:rem:weil-reproof} that already forces
$E[N](\ksep) \not\subseteq E(\F_q)$, so the $N$-torsion becomes rational only over an
extension. The curve $E \colon y^2 = x^3 + 1$ over $\F_5$ has $\#E(\F_5) = 6$ and
$E(\F_{25}) \cong (\Z/6\Z)^2$, so all of $E[3](\Fbar_5)$ is already rational over the
quadratic extension: small enough to list, large enough to have something to say.

\section{A quaternionic endomorphism ring}
\label{weil-end-aut:sec:end}

The classification of \cite{Silverman09}*{Theorem~III.9.3} puts the
endomorphism ring of an elliptic curve in one of three cases: $\Z$, an order in
an imaginary quadratic field, or an order in a definite quaternion algebra over
$\Q$. The first two are easy to exhibit: a generic curve has endomorphism ring
$\Z$, and $y^2 = x^3 + x$ has $\Z[i]$ inside its endomorphism ring. Here is an
example of the third, on the smallest curve that has one.

\begin{wsdefinition}[{The quaternion algebra $D$ \cite{Voight21}*{Definition~2.2.1}}]
\label{weil-end-aut:def:quaternion}
Let $D \coloneqq (-1, -3 \mid \Q)$, the quaternion algebra with $i^2 = -1$, $j^2 = -3$
and $ij = -ji$. Its \defi{reduced discriminant} is the
product of the primes $p$ at which $D \otimes_\Q \Q_p$ is not split.
\end{wsdefinition}

\begin{wsdefinition}[{Definite quaternion algebras \cite{Voight21}*{Main Theorem~5.4.4 and \S 14.1}}]
\label{weil-end-aut:def:definite}
A quaternion algebra $B$ over $\Q$ is \defi{definite} when $B \otimes_\Q \R$ is
Hamilton's $\mathbb{H}$ rather than $M_2(\R)$, equivalently when its reduced norm is a
positive definite quadratic form in the four coordinates. Otherwise $B$ is \defi{indefinite}.
\end{wsdefinition}

\begin{wsdefinition}[{Orders and the reduced discriminant \cite{Voight21}*{\S 15.2}}]
\label{weil-end-aut:def:discrd}
Orders are as in \cite{Voight21}*{Definition~10.2.1}. The \defi{reduced discriminant}
$\discrd(\mathcal{O})$ of an order $\mathcal{O} \subseteq D$ is the positive integer with
$\discrd(\mathcal{O})^2 = \lvert\det\bigl(\trd(\alpha_k\alpha_l)\bigr)_{k,l}\rvert$
for a $\Z$-basis $\alpha_1, \dots, \alpha_4$ of $\mathcal{O}$.
\end{wsdefinition}

\begin{wsproposition}[{The characteristic polynomial of Frobenius \cite{Silverman09}*{\S V.2}}]
\label{weil-end-aut:prop:charpoly}
Let $E$ be an elliptic curve over $\F_q$ with $\Frob$ its $q$-power Frobenius,
$a \coloneqq q + 1 - \#E(\F_q)$, and $\ell \nmid q$ a prime. Then $\Frob$ satisfies
$\Frob^2 - [a]\Frob + [q] = 0$ in $\End(E)$: by
\cite{Silverman09}*{Proposition~III.8.6} the characteristic polynomial of
$\Frob_\ell$ is $T^2 - aT + q$, and every $\varphi \in \End(E)$ satisfies
$\varphi^2 - [\tr\varphi_\ell]\varphi + [\deg\varphi] = 0$ in $\End(E)$, not
merely on $T_\ell(E)$.
\end{wsproposition}

\begin{wscorollary}[{Factoring through a separable isogeny \cite{Silverman09}*{Corollary~III.4.11}}]
\label{weil-end-aut:cor:factor}
If $\varphi\colon E_1 \to E_2$ is separable and
$\psi\colon E_1 \to E_3$ is an isogeny with
$\ker\varphi \subseteq \ker\psi$, then $\psi = \lambda \circ \varphi$ for a
unique isogeny $\lambda\colon E_2 \to E_3$.
\end{wscorollary}

\begin{wsproposition}[{Where $D$ is split \cite{Voight21}*{Exercise~11.11 and \S 14.1}}]
\label{weil-end-aut:prop:ram}
The algebra $D \otimes_\Q \Q_p$ is split for every prime $p \neq 3$, and
$D \otimes_\Q \Q_3$ is not. So $D$ is a division algebra, and its reduced
discriminant is $3$.
\end{wsproposition}

\begin{wstheorem}[{Maximal orders \cite{Voight21}*{Theorem~15.5.5}}]
\label{weil-end-aut:thm:maximal}
An order $\mathcal{O} \subseteq D$ is maximal if and only if
$\discrd(\mathcal{O}) = 3$.
\end{wstheorem}

\begin{wsproposition}[{The index and the discriminant \cite{Voight21}*{\S 15.2}}]
\label{weil-end-aut:prop:index}
If $\mathcal{O} \subseteq \mathcal{O}'$ are orders of $D$, then
$\discrd(\mathcal{O}) = [\mathcal{O}' : \mathcal{O}]\discrd(\mathcal{O}')$.
\end{wsproposition}

For the rest of this section
\[
  E\colon y^2 = x^3 - x \quad\text{over } \F_3 ,
\]
so that the generic $E$ of \cref{weil-end-aut:prop:charpoly} is this curve, with $q = 3$.

\begin{problem}[$\End(E_{\F_9})$ is a maximal order in a quaternion algebra]
\label{weil-end-aut:prob:end}
Every part below has a hand computation and a check to run in \Magma. The curve is small
enough that $[i]$ and $\Frob$ can be verified on all sixteen points of $E(\F_9)$.
\begin{enumerate}
  \item\label{weil-end-aut:it:end-warmup} Check that $E$ is an elliptic curve over $\F_3$, list
        $E(\F_3)$, and deduce that $a = 0$ and that
        \[
          \Frob^2 = [-3] .
        \]
        Record also that every point of
        $E[2](\Fbar_3)$ lies in $E(\F_3)$: this is what makes
        part~\ref{weil-end-aut:it:end-maximal} work.
\begin{magma}
E3 := EllipticCurve([GF(3) | -1, 0]);    // prints as y^2 = x^3 + 2*x
Points(E3);
#E3;                                // 4
3 + 1 - #E3;                        // 0        (so a = 0)
jInvariant(E3);                     // 0
\end{magma}
        \Magma\ hands you the point count and the $j$-invariant. The identity
        $\Frob^2 = [-3]$ is yours to deduce, from
        \cref{weil-end-aut:prop:charpoly}. Read $E[2](\Fbar_3) \subseteq E(\F_3)$
        off the list of rational points --- which points of $y^2 = f(x)$ have order $2$?
\begin{solution}
% Write your solution here.
\end{solution}
  \item\label{weil-end-aut:it:end-i} Show that $x^2 + 1$ is irreducible over $\F_3$, so that
        $\F_9 = \F_3(i)$ with $i^2 = -1$. Show that
        \[
          [i](x, y) \coloneqq (-x, iy)
        \]
        defines an endomorphism of $E_{\F_9}$, and that $[i]^2 = [-1]$. Why is a
        morphism $E \to E$ fixing $O$ automatically a group
        homomorphism? Mind \cref{weil-end-aut:warn:gf9} before you start typing.
\begin{magma}
F9 := GF(9);  i := Sqrt(F9!(-1));
i^2, MinimalPolynomial(i);               // 2 ($.1^2 + 1), and 2 = -1
E9 := ChangeRing(E3, F9);  #E9;          // 16
imap := func<P | P eq E9!0 select E9!0 else E9![-P[1], i*P[2]]>;
forall{P : P in Points(E9) | imap(P) in Points(E9)};        // [i] is well defined
forall{<P,Q> : P, Q in Points(E9) | imap(P+Q) eq imap(P)+imap(Q)};  // additive
forall{P : P in Points(E9) | imap(imap(P)) eq -P};          // [i]^2 = [-1]
// all true
\end{magma}
        The second \mg{forall} is additivity, checked on all $16^2$ pairs. That is a
        proof only because the curve is finite, and it is no answer to the question
        above: additivity is not a coincidence to be confirmed point by point but a
        consequence of fixing $O$.
\begin{solution}
% Write your solution here.
\end{solution}
  \item\label{weil-end-aut:it:end-anti} Show that $\Frob[i] = -[i]\Frob$. Conclude that
        $\End(E_{\F_9})$ is not commutative, so that $E_{\F_9}$ falls into the
        third case of the classification.
\begin{magma}
Frob := func<P | P eq E9!0 select E9!0 else E9![P[1]^3, P[2]^3]>;
forall{P : P in Points(E9) | Frob(Frob(P)) eq -3*P};             // Frob^2 = [-3]
forall{P : P in Points(E9) | Frob(imap(P)) eq -imap(Frob(P))};   // Frob[i] = -[i]Frob
// both true
\end{magma}
        The first line verifies the $\Frob^2 = [-3]$ you deduced in \ref{weil-end-aut:it:end-warmup} from
        the characteristic polynomial, which is worth having independently. For the
        second, do the computation by hand as well: it comes down to a single fact about
        $i^3$.
\begin{solution}
% Write your solution here.
\end{solution}
  \item\label{weil-end-aut:it:end-algebra} Deduce that
        \[
          \End(E_{\F_9}) \otimes_\Z \Q \cong D = (-1, -3 \mid \Q).
        \]
        Write down the reduced norm of $D$ and check that $D$ is definite. How
        does the reduced norm compare with the degree of an endomorphism?
\begin{magma}
D<ii,jj> := QuaternionAlgebra<Rationals() | -1, -3>;
RamifiedPrimes(D), Discriminant(D);      // [ 3 ] 3
IsDefinite(D);                           // true
Trace(jj), Norm(jj);                     // 0 3
\end{magma}
        The two arguments are the squares of the generators, so the algebra \Magma\
        builds is the $D$ of \cref{weil-end-aut:def:quaternion}. Here \mg{Trace} and \mg{Norm} on a quaternion
        algebra are the reduced trace and the reduced norm, and the pair they return for
        \mg{jj} should be recognizable as $(\tr \Frob, \deg \Frob)$. All of this is about $D$
        alone. That $D$ is the \emph{whole} of $\End(E_{\F_9}) \otimes \Q$ is the
        one claim here \Magma\ cannot be asked, and it needs a rank count.
\begin{solution}
% Write your solution here.
\end{solution}
  \item\label{weil-end-aut:it:end-maximal} Show that $\frac{1 + \Frob}{2}$ and $\frac{i + i\Frob}{2}$ lie in
        $\End(E_{\F_9})$.\hint{use \cref{weil-end-aut:cor:factor} with
        $\varphi = [2]$.} Then show that
        \[
          \mathcal{O} \coloneqq \Z + \Z i + \Z\,\frac{1 + \Frob}{2} + \Z\,\frac{i + i\Frob}{2}
        \]
        is an order of $D$,
        compute $\discrd(\mathcal{O})$, and conclude with \cref{weil-end-aut:thm:maximal} that
        $\End(E_{\F_9}) = \mathcal{O}$ is a maximal order.
\begin{magma}
O0 := QuaternionOrder([D!1, ii, jj, ii*jj]);
Discriminant(O0);                        // 12
O  := QuaternionOrder([D!1, ii, (1+jj)/2, (ii+ii*jj)/2]);
Discriminant(O), IsMaximal(O);           // 3 true
\end{magma}
        \mg{QuaternionOrder} takes a $\Z$-basis and \mg{Discriminant} of an order is its
        reduced discriminant, so the two lists above are the $\Z$-bases of
        $\mathcal{O}_0$ and $\mathcal{O}$. Do not reach for \mg{Index} to
        compare the two orders: \Magma\ has no index for quaternion orders, and answers
        \mg{Runtime error in 'Index': Bad argument types}. Recover
        $[\mathcal{O} : \mathcal{O}_0]$ from the two discriminants instead, by
        \cref{weil-end-aut:prop:index}, and check it against the determinant of
        the change of basis.

        The first sentence of this part is where the real work is, and none of the code
        touches it: $\frac{1+\Frob}{2}$ is a quotient in a ring with no division, and only
        the kernel argument puts it in $\End(E_{\F_9})$.
\begin{solution}
% Write your solution here.
\end{solution}
\end{enumerate}
\end{problem}

\begin{wswarning}[$\F_9$ does not come with $i$]
\label{weil-end-aut:warn:gf9}
\mg{GF(9)} does \emph{not} come with $i$. \Magma's default generator of $\F_9$ is a root
of the Conway polynomial $x^2 + 2x + 2$, so writing \mg{F9<i> := GF(9)} gives an $i$
with $i^2 = i + 1$, not $i^2 = -1$, and every check in \cref{weil-end-aut:prob:end} would silently
fail. Construct the square root of $-1$ explicitly, and confirm it: \mg{MinimalPolynomial}
of the element you build must be $x^2 + 1$, and \mg{i^2} must print as $2$, which is
$-1$ in $\F_9$.
\end{wswarning}

\section{Automorphisms in characteristic \texorpdfstring{$2$ and $3$}{2 and 3}}
\label{weil-end-aut:sec:aut}

For $\Char k = 2, 3$ the geometric automorphism group of an elliptic curve may be
non-abelian, which is the subject of \cref{weil-end-aut:prob:aut}. In the following statements, $E$ is
an elliptic curve over $k$, given by a Weierstrass equation with coefficients
$a_1, a_2, a_3, a_4, a_6$.

\begin{wsdefinition}[Changes of variables]
\label{weil-end-aut:def:cov}
A \defi{change of variables} preserving a Weierstrass equation is a substitution
\begin{equation}
\label{weil-end-aut:eq:cov}
  x = u^2x' + r, \qquad y = u^3y' + u^2sx' + t, \qquad u \in \kbar^\times,\ r, s, t \in \kbar.
\end{equation}
For such $u, r, s, t$ we write
\[
  \Psi_{u,r,s,t}(x, y) \coloneqq (u^2x + r,\ u^3y + u^2sx + t).
\]
\end{wsdefinition}

\begin{wstheorem}[{The automorphism group \cite{Silverman09}*{Theorem~III.10.1}}]
\label{weil-end-aut:thm:aut}
The group $\Aut(E_{\kbar})$ is finite, of order
\[
  \#\Aut(E_{\kbar}) =
  \begin{cases}
    2 & \text{if } j(E) \neq 0, 1728, \\
    4 & \text{if } j(E) = 1728 \text{ and } \Char k \neq 2, 3, \\
    6 & \text{if } j(E) = 0 \text{ and } \Char k \neq 2, 3, \\
    12 & \text{if } j(E) = 0 = 1728 \text{ and } \Char k = 3, \\
    24 & \text{if } j(E) = 0 = 1728 \text{ and } \Char k = 2.
  \end{cases}
\]
\end{wstheorem}

\begin{wscorollary}[{The cyclic cases \cite{Silverman09}*{Corollary~III.10.2}}]
\label{weil-end-aut:cor:aut-cyclic}
If $\Char k \neq 2, 3$, then $\Aut(E_{\ksep}) \cong \mu_n(\ksep)$ with
$n = \#\Aut(E_{\kbar})$, via $[\zeta](x,y) \coloneqq (\zeta^2x, \zeta^3y)$. In particular,
the geometric automorphism group is cyclic in that case.
\end{wscorollary}

\begin{wsproposition}[{The action on the coefficients \cite{Silverman09}*{Proposition~III.3.1(b) and Table~3.1}}]
\label{weil-end-aut:prop:cov}
Under a change of variables \eqref{weil-end-aut:eq:cov} the coefficients transform by
\begin{equation}
\label{weil-end-aut:eq:cov-coeffs}
\begin{aligned}
  u a_1' &= a_1 + 2s, \\
  u^2 a_2' &= a_2 - sa_1 + 3r - s^2, \\
  u^3 a_3' &= a_3 + ra_1 + 2t, \\
  u^4 a_4' &= a_4 - sa_3 + 2ra_2 - (t + rs)a_1 + 3r^2 - 2st, \\
  u^6 a_6' &= a_6 + ra_4 + r^2a_2 + r^3 - ta_3 - t^2 - rta_1 .
\end{aligned}
\end{equation}
A substitution \eqref{weil-end-aut:eq:cov} carries $E$ to the curve with coefficients $a_i'$,
so it is a change of coordinates from $E$ to itself exactly when $a_i' = a_i$
for every $i$, and then $\Psi_{u,r,s,t}$ is the automorphism of $E$ that it defines.
Setting $a_i' = a_i$ in \eqref{weil-end-aut:eq:cov-coeffs} therefore turns those five lines into five
equations in the four unknowns $u, r, s, t$, whose solutions are exactly the elements of
$\Aut(E_{\kbar})$.
\end{wsproposition}

The whole point of this section is visible in \eqref{weil-end-aut:eq:cov-coeffs}: observe
the $2$'s and $3$'s in the coefficients. The characteristic enters through the terms
$2s$, $2t$ and $3r$. Away from $2$
and $3$ they are invertible multiples of $s$, $t$ and $r$, which determines
those three and leaves $u$ as the only free parameter: that is why the group is
cyclic there. In characteristic $3$ the term $3r$ vanishes, in characteristic
$2$ the terms $2s$ and $2t$ do, and the parameters they no longer pin down are
what make the two groups of \cref{weil-end-aut:prob:aut} large, and non-cyclic.

\begin{problem}[The two non-cyclic cases]
\label{weil-end-aut:prob:aut}
As in \cref{weil-end-aut:prob:end}, each part pairs a hand computation with a check in
\Magma. Here the machine is best at naming the group, which is the trickiest
thing to do by hand.
\begin{enumerate}
  \item\label{weil-end-aut:it:aut-j} Let $\Char k = 3$. Show that every short Weierstrass equation
        $y^2 = x^3 + Ax + B$ has $j = 0$, and explain why a curve with $j \neq 0$ in
        characteristic $3$ needs an $a_2x^2$ term. So the case
        ``$j = 0 = 1728$, $\Char k = 3$'' of \cref{weil-end-aut:thm:aut} is exactly the
        short-Weierstrass one. The first claim is a one-line sweep over the whole family:
\begin{magma}
{ jInvariant(EllipticCurve([GF(3) | A, B])) : A, B in GF(3) | A ne 0 };  // { 0 }
jInvariant(EllipticCurve([GF(3) | 0, 1, 0, 0, -1]));   // 1
\end{magma}
        The condition \mg{A ne 0} is what keeps the equation non-singular: work out what
        $\Delta \neq 0$ becomes for $y^2 = x^3 + Ax + B$ once $27 = 0$, and you will find
        that $B$ drops out. The set that comes back has one element, which is the first
        claim. The second line is a curve with an $a_2x^2$ term, and its $j$-invariant is
        not in that set.
\begin{solution}
% Write your solution here.
\end{solution}
  \item\label{weil-end-aut:it:aut-three-count} Let $\Char k = 3$ and $E\colon y^2 = x^3 - x$. Using
        \eqref{weil-end-aut:eq:cov-coeffs}, show first that $s = t = 0$, and then that
        \eqref{weil-end-aut:eq:cov} is an automorphism of $E$ if and only if
        \[
          u^4 = 1 \qquad\text{and}\qquad r^3 - r = 0 ,
        \]
        i.e.~$u \in \mu_4(\Fbar_3) \subseteq \F_9^\times$ and $r \in \F_3$. Conclude
        $\#\Aut(E_{\Fbar_3}) = 12$, and that all $12$ are already defined over $\F_9$.
        Where did $r$ become free, and what would have forced $r = 0$ if
        $\Char k \neq 3$? Two counts settle both claims at once:
\begin{magma}
E3 := EllipticCurve([GF(3) | -1, 0]);
#AutomorphismGroup(E3), #AutomorphismGroup(ChangeRing(E3, GF(9)));   // 6 12
\end{magma}
        Predict both numbers from your conditions on $u$ and $r$ before running it. Over
        $\F_3$ the fourth roots of unity are only $\pm 1$, while $r$ already ranges over
        all of $\F_3$. Over $\F_9$ all four roots of unity appear. The second number
        being $12$, and being reached over $\F_9$ rather than over some larger field, is
        exactly the claim that all twelve automorphisms are defined there.
\begin{solution}
% Write your solution here.
\end{solution}
  \item\label{weil-end-aut:it:aut-three-group} Write $\varphi_{u,r} \coloneqq \Psi_{u,r,0,0}$, so that
        $\varphi_{u,r}(x, y) = (u^2x + r, u^3y)$.
        Compute $\varphi_{u',r'} \circ \varphi_{u,r}$ and deduce
        \[
          \Aut(E_{\Fbar_3}) \cong \Z/3\Z \rtimes \Z/4\Z ,
        \]
        where $\Z/3\Z = \F_3$ is the group of translations $r$, where
        $\Z/4\Z = \mu_4(\Fbar_3)$ is the group of the $u$, and where $u$ acts on $r$ by
        $r \mapsto u^2r$. Since $u^4 = 1$ forces $u^2 = \pm 1$, that action is not
        faithful: it factors through the squaring map
        $\mu_4(\Fbar_3) \twoheadrightarrow \brk{\pm 1} = \Aut(\Z/3\Z)$, so $\pm 1$ act trivially
        and $\pm i$ act by $r \mapsto -r$. Show this group is non-abelian with a
        \emph{unique} element of order $2$, and deduce that it is neither cyclic, nor
        dihedral, nor $A_4$: it is the dicyclic group of order $12$. Which
        automorphism is the unique involution? \mg{AutomorphismGroup} returns a polycyclic
        group and \mg{IdentifyGroup} gives its label in the small-groups database.
\begin{magma}
G3 := AutomorphismGroup(ChangeRing(E3, GF(9)));
IdentifyGroup(G3), GroupName(G3);        // <12, 1>  C3:C4
IsAbelian(G3), IsCyclic(G3);             // false false
{* Order(g) : g in G3 *};                // {* 1, 2, 3^^2, 4^^6, 6^^2 *}
\end{magma}
        The last line is the multiset of element orders. Say how it rules out the
        dihedral group and $A_4$ --- how many involutions does each of those have? ---
        and match its single element of order $2$ against the automorphism you named.
\begin{solution}
% Write your solution here.
\end{solution}
  \item\label{weil-end-aut:it:aut-two-count} Let $\Char k = 2$ and $E\colon y^2 + y = x^3$. Using
        \eqref{weil-end-aut:eq:cov-coeffs}, show that $a_1' = a_1$ automatically, that $a_2' = 0$
        forces $r = s^2$, and then that \eqref{weil-end-aut:eq:cov} is an automorphism if and only if
        \[
          u^3 = 1, \qquad s^4 = s, \qquad t^2 + t = s^6 .
        \]
        Deduce that $s \in \F_4$, count the solutions, and conclude
        $\#\Aut(E_{\Fbar_2}) = 24$, all defined over $\F_4$. The same two counts as in
        \ref{weil-end-aut:it:aut-three-count}:
\begin{magma}
E2 := EllipticCurve([GF(2) | 0, 0, 1, 0, 0]);          // y^2 + y = x^3
#AutomorphismGroup(E2), #AutomorphismGroup(ChangeRing(E2, GF(4)));   // 2 24
\end{magma}
        Predict both first. Over $\F_2$ the three conditions bite hard: $u^3 = 1$ leaves
        only $u = 1$, and of the values $s \in \F_2$ one of them makes $t^2 + t = s^6$
        unsolvable in $\F_2$. The gap between the two numbers here is much wider than in
        characteristic $3$, and that gap is the content of ``all defined over $\F_4$''.
\begin{solution}
% Write your solution here.
\end{solution}
  \item\label{weil-end-aut:it:aut-two-group} Write $\psi_{s,t} \coloneqq \Psi_{1,s^2,s,t}$ for the
        automorphisms with $u = 1$, so that $\psi_{s,t}(x,y) = (x + s^2, y + sx + t)$.
        Show that
        \[
          \psi_{s',t'} \circ \psi_{s,t} = \psi_{s + s',\, t + t' + s's^2},
          \qquad \psi_{s,t}^2 = \psi_{0, s^3},
        \]
        and deduce that these $8$ automorphisms form a copy of the quaternion group
        $Q_8$. For the remaining factor, let $\omega \in \F_4$ be a primitive cube root
        of unity and $\varphi_\omega \coloneqq \Psi_{\omega,0,0,0}$, which has order $3$.
        Check that conjugation by it is
        \[
          \varphi_\omega \circ \psi_{s,t} \circ \varphi_\omega^{-1} = \psi_{\omega s,\, t} ,
        \]
        so $\varphi_\omega$ normalizes $Q_8$, fixes its center
        $\langle \psi_{0,1} \rangle$, and cycles the three pairs $\pm i, \pm j, \pm k$
        --- the order-$3$ automorphism of $Q_8$. Since
        $\langle \varphi_\omega \rangle \cap Q_8 = 1$, conclude that
        \[
          \Aut(E_{\Fbar_2}) \cong Q_8 \rtimes \Z/3\Z \cong \SL_2(\F_3) .
        \]
        As before:
\begin{magma}
G2 := AutomorphismGroup(ChangeRing(E2, GF(4)));
IdentifyGroup(G2), GroupName(G2);        // <24, 3>  SL(2,3)
{* Order(g) : g in G2 *};                // {* 1, 2, 3^^8, 4^^6, 6^^8 *}
\end{magma}
        Exactly one element of order $2$ again. Explain why that alone pins down $Q_8$
        among the groups of order $8$, and why it does \emph{not} by itself pin down
        $\SL_2(\F_3)$ among the groups of order $24$ --- what extra fact did you use?
\begin{solution}
% Write your solution here.
\end{solution}
\end{enumerate}
\end{problem}

\begin{wsremark}[Both groups are unit groups]
\label{weil-end-aut:rem:aut-units}
Automorphisms are the units of the endomorphism ring, $\Aut(E_{\kbar}) =
\End(E_{\kbar})^\times$. In characteristic $2$ and $3$ the curve with $j = 0$ has
$\End(E_{\kbar})$ a maximal order in the quaternion algebra over
$\Q$ that is not split exactly at $p$ and $\infty$ \cite{Silverman09}*{Theorem~V.3.1},
\cite{Voight21}*{Theorem~42.1.9}. For $p = 3$ that order is the one computed
in \cref{weil-end-aut:prob:end}\ref{weil-end-aut:it:end-maximal}, and its unit group is the dicyclic group of
order $12$ of \cref{weil-end-aut:prob:aut}\ref{weil-end-aut:it:aut-three-group}. For $p = 2$ it is the Hurwitz order
$\Z + \Z i + \Z j + \Z\frac{1 + i + j + ij}{2}$ in $(-1,-1 \mid \Q)$
\cite{Voight21}*{\S 11.1}, whose unit group is the binary tetrahedral group of order
$24$ of \cref{weil-end-aut:prob:aut}\ref{weil-end-aut:it:aut-two-group}. This explains the shape of \cref{weil-end-aut:thm:aut}: when
$\End(E_{\kbar})$ is commutative it is $\Z$ or an imaginary quadratic order, whose unit
groups are the cyclic groups of orders $2$, $4$ and $6$, giving its first three lines, while a non-commutative $\End(E_{\kbar})$ is a quaternion order, and
its unit group need not be, and here is not, cyclic.

Both counts can be recovered from the orders alone, with no elliptic curve anywhere in
the computation:
\begin{magma}
for d in [-1, -3] do
  Dd := QuaternionAlgebra<Rationals() | -1, d>;
  Od := MaximalOrder(Dd);
  print RamifiedPrimes(Dd), #UnitGroup(Od);
end for;
// [ 2 ] 24     <- Hurwitz order, characteristic 2
// [ 3 ] 12     <- the endomorphism ring of (*\cref{weil-end-aut:prob:end}*), characteristic 3
\end{magma}
So the $24$ and the $12$ of \cref{weil-end-aut:thm:aut} are the unit counts of the two maximal
quaternion orders.
\end{wsremark}

\subsection*{Looking ahead}
Curves whose Frobenius has $a = 0$, like the one of \cref{weil-end-aut:sec:end}, have a name and a
theory of their own. They come later in the course, with Hasse's bound and the structure
of $\End$ over $\F_q$. In characteristic $0$ the third case of the classification never occurs, and
$\End(E_{\kbar}) \neq \Z$ becomes complex multiplication. The oddness of
$\rho_{E,\ell}$ from \cref{weil-end-aut:prob:weil-det}\ref{weil-end-aut:it:weil-odd} is the first hint of
modularity.

\bigskip

\part[Differentials, dual isogenies, and the Tate module]{Differentials, dual isogenies, and the Tate module}
\label{diffs-dual-tate:part}
\noindent
Here are three computations by hand, on the invariant differential, the dual isogeny and
the Tate module. Each compares the multiplicative group $\Gm$ with the elliptic curve
\[
  E\colon y^2 = x^3 - x, \qquad \omega_E \coloneqq \frac{\dd x}{2y}.
\]
Throughout, $k$ is a field with $\Char k \neq 2$, $\kbar$ an algebraic closure and
$\ksep \subseteq \kbar$ the separable closure. Here $\Gm = \PP^1_k \smallsetminus \brk{0, \infty}$
has coordinate $t$ and invariant differential $\omega_{\Gm} \coloneqq \dd t/t$.

\subsection*{Reading}
\cite{Silverman09}*{\S\S III.4--III.7}, and \cite{Silverman09}*{\S III.8} for the Weil
pairing of \cref{diffs-dual-tate:prop:weil}.

\subsection*{A formula}
We use the following duplication formula on $E$, derived in
\cref{diffs-dual-tate:prob:omega}\ref{diffs-dual-tate:it:omega-dup}: if $x \in k(E)$ is the $x$-coordinate and
$[2]\colon E \to E$ is multiplication by $2$, then
\begin{equation}
\label{diffs-dual-tate:eq:dup}
  [2]^*x \coloneqq x \circ [2] = \frac{(x^2 + 1)^2}{4(x^3 - x)} .
\end{equation}

\section{Invariant differentials}

\begin{problem}[Poles and pullbacks]
\label{diffs-dual-tate:prob:omega}
\begin{enumerate}
  \item\label{diffs-dual-tate:it:omega-gm} Rewrite $\omega_{\Gm}$ in the coordinate $s \coloneqq 1/t$ near $\infty$, and
        compute $\div(\omega_{\Gm})$ on $\PP^1_k$. Compute $[m]^*\omega_{\Gm}$ for
        $[m](t) = t^m$, $m \in \Z$.
\begin{solution}
% Write your solution here.
\end{solution}
  \item\label{diffs-dual-tate:it:omega-no-zeros} Show that $\omega_E = \dd y/(3x^2 - 1)$, and that $y$ and $3x^2 - 1$ have no
        common zero on $E$. Deduce that $\omega_E$ has no zeros or poles on
        $E \smallsetminus \brk{O}$.
\begin{solution}
% Write your solution here.
\end{solution}
  \item\label{diffs-dual-tate:it:omega-at-O} Using $\ord_O x = -2$ and $\ord_O y = -3$, show that $\ord_O \omega_E = 0$.
\begin{solution}
% Write your solution here.
\end{solution}
  \item\label{diffs-dual-tate:it:omega-minus-one} Compute $[-1]^*\omega_E$ from $[-1](x, y) = (x, -y)$, and compare with
        $[-1]^*\omega_{\Gm}$.
\begin{solution}
% Write your solution here.
\end{solution}
  \item\label{diffs-dual-tate:it:omega-dup} We derive \eqref{diffs-dual-tate:eq:dup} by hand. Here $x, y \in k(E)$ are
        functions. For a point $P \in E(\kbar)$ write $x_P \coloneqq x(P)$ and
        $y_P \coloneqq y(P)$ for its coordinates.
        \begin{enumerate}
          \item\label{diffs-dual-tate:it:omega-dup-point} Let $P \in E(\kbar)$ with $y_P \neq 0$. Show that the tangent line to
                $E$ at $P$ has slope $\lambda_P = (3x_P^2 - 1)/(2y_P)$, and use that its
                third point of intersection with $E$ is $-[2]P$ to show that
                $x([2]P) = \lambda_P^2 - 2x_P$.
\begin{solution}
% Write your solution here.
\end{solution}
          \item\label{diffs-dual-tate:it:omega-dup-function} Let $\lambda \coloneqq (3x^2 - 1)/(2y) \in k(E)$, so that
                $\lambda(P) = \lambda_P$. Explain why \ref{diffs-dual-tate:it:omega-dup-point} gives the identity
                $[2]^*x = \lambda^2 - 2x$ of functions in $k(E)$, and simplify it using
                $y^2 = x^3 - x$.
\begin{solution}
% Write your solution here.
\end{solution}
        \end{enumerate}
  \item\label{diffs-dual-tate:it:omega-magma} Solve this part in \Magma, working in the function field
        $\Q(E) = \Q(x)[y]/(y^2 - x^3 + x)$. Compute $[2]^*x$ and $[2]^*y$ from the
        function $\lambda$ of~\ref{diffs-dual-tate:it:omega-dup}, check \eqref{diffs-dual-tate:eq:dup}, and show that
        \[
          [2]^*x - e = \frac{g_e(x)^2}{4(x^3 - x)}, \qquad
          g_0 = x^2 + 1,\quad g_1 = x^2 - 2x - 1,\quad g_{-1} = x^2 + 2x - 1,
        \]
        for $e \in \brk{0, 1, -1}$. Show that
        $[2]^*y = \pm (x^2 + 1)(x^4 - 6x^2 + 1)/(8y^3)$ and
        $[2]^*\omega_E = \pm 2\omega_E$. Which sign is the correct one?
\begin{solution}
% Write your solution here.
\end{solution}
\end{enumerate}
\end{problem}

\section{Dual isogenies}

Consider the elliptic curves
\[
  E_1\colon y_1^2 = x_1^3 - x_1, \qquad
  E_2\colon y_2^2 = x_2^3 + 4x_2, \qquad
  E_3\colon y_3^2 = x_3^3 - 16x_3,
\]
so that $E_1 = E$, and write $\omega_i \coloneqq \dd x_i/(2y_i)$ for the invariant
differential of $E_i$. Consider the maps
\[
  \begin{aligned}
    \varphi\colon E_1 &\to E_2, &
    (x_1, y_1) &\mapsto \paren{\frac{x_1^2 - 1}{x_1},\ \frac{y_1(x_1^2 + 1)}{x_1^2}},\\
    \psi\colon E_2 &\to E_3, &
    (x_2, y_2) &\mapsto \paren{\frac{y_2^2}{x_2^2},\ \frac{y_2(x_2^2 - 4)}{x_2^2}},\\
    \iota\colon E_3 &\to E_1, &
    (x_3, y_3) &\mapsto \paren{\frac{x_3}{4},\ \frac{y_3}{8}}.
  \end{aligned}
\]
The maps $\varphi$ and $\psi$ are special cases of
\cite{Silverman09}*{Example~III.4.5}, which treats $y^2 = x^3 + ax^2 + bx$ in general,
so they are isogenies. Also $\iota$ is an isomorphism. That example computes the kernels
and the dual, but not the constants $c_\alpha$ of \cref{diffs-dual-tate:def:c}: that is
part~\ref{diffs-dual-tate:it:dual-c-phi} of \cref{diffs-dual-tate:prob:dual}.

\begin{wsdefinition}[The constant $c_\alpha$]
\label{diffs-dual-tate:def:c}
For an isogeny $\alpha\colon E_i \to E_j$, define $c_\alpha$ by
$\alpha^*\omega_j = c_\alpha\,\omega_i$.
\end{wsdefinition}

\begin{problem}[A $2$-isogeny and its dual]
\label{diffs-dual-tate:prob:dual}
\begin{enumerate}
  \item\label{diffs-dual-tate:it:dual-phi} Check that $\varphi$ lands on $E_2$. Show that $\ker\varphi = \brk{O_1, (0,0)}$,
        that $\varphi(\pm 1, 0) = (0, 0)$, and that $\deg\varphi = 2$.\hint{write
        $y_1$ in terms of $\varphi^*y_2$ and $x_1$, and find a quadratic equation for
        $x_1$ over $\varphi^*k(E_2)$.}
\begin{solution}
% Write your solution here.
\end{solution}
  \item\label{diffs-dual-tate:it:dual-c-phi} Compute $c_\varphi$.
\begin{solution}
% Write your solution here.
\end{solution}
  \item\label{diffs-dual-tate:it:dual-psi-iota} Check that $\psi$ lands on $E_3$ and that $\iota$ lands on $E_1$, and compute
        $c_\psi$ and $c_\iota$.
\begin{solution}
% Write your solution here.
\end{solution}
  \item\label{diffs-dual-tate:it:dual-sign} Show that $(\iota \circ \psi \circ \varphi)^*x_1 = [2]^*x_1$, so that
        $\iota\circ\psi\circ\varphi = \pm[2]$. Use \ref{diffs-dual-tate:it:dual-c-phi} and \ref{diffs-dual-tate:it:dual-psi-iota} to decide the sign, and
        conclude that $\hat\varphi = \iota \circ \psi$, that
        $c_\varphi c_{\hat\varphi} = \deg\varphi$, and that
        $\ker\hat\varphi = \varphi(E_1[2](\kbar))$.
\begin{solution}
% Write your solution here.
\end{solution}
  \item\label{diffs-dual-tate:it:dual-gm} Is there an isogeny $\Gm \to \Gm$ of degree $2$? If so,
        does it have a dual?
\begin{solution}
% Write your solution here.
\end{solution}
\end{enumerate}
\end{problem}

\section{The Tate module}

We first describe the Tate module of $\Gm$, the familiar case, and then repeat each step
for $E$. What links the two is the Weil pairing, which you may use without proof.

\begin{wsproposition}[{The Weil pairing \cite{Silverman09}*{Proposition~III.8.1}}]
\label{diffs-dual-tate:prop:weil}
Let $m \geq 1$ with $\Char k \nmid m$. There is a pairing
\[
  e_m\colon E[m](\ksep) \times E[m](\ksep) \to \mu_m(\ksep)
\]
that is bilinear, alternating ($e_m(P, P) = 1$), non-degenerate, and Galois equivariant:
$\sigma\bigl(e_m(P, Q)\bigr) = e_m(\sigma P, \sigma Q)$ for every
$\sigma \in \Gal_k$.
Here $E[m](\ksep) = E[m](\kbar)$, since $\Char k \nmid m$. It follows that if $(P, Q)$ is
a basis of $E[m](\ksep)$ over $\Z/m\Z$, then $e_m(P, Q)$ is a primitive $m$-th root of
unity, and if $\sigma$ acts on this basis by the matrix $A \in \GL_2(\Z/m\Z)$, then
$\sigma\bigl(e_m(P, Q)\bigr) = e_m(P, Q)^{\det A}$.
\end{wsproposition}

\begin{problem}[The $2$-power torsion of $y^2 = x^3 - x$]
\label{diffs-dual-tate:prob:tate}
Regard $E$ as a curve over $\Q$, fix $\Qbar \subset \C$, and recall that $E[m](\Qbar)$ is the kernel
of $[m]$ on $E(\Qbar)$. Let $\sigma$ be complex conjugation, acting on $\Qbar$ and on
$E(\Qbar)$ coordinatewise.
\begin{enumerate}
  \item\label{diffs-dual-tate:it:tate-gm} Recall that $T_2(\Gm) \coloneqq \varprojlim_n \mu_{2^n}(\Qbar)$, with
        transition maps $\zeta \mapsto \zeta^2$. Find a $\Z_2$-basis of $T_2(\Gm)$, and write
        down its first three terms. How does $\sigma$ act on $\mu_4(\Qbar)$, and on
        $T_2(\Gm)$?
\begin{solution}
% Write your solution here.
\end{solution}
  \item\label{diffs-dual-tate:it:tate-e4} List $E[2](\Qbar)$. Using \cref{diffs-dual-tate:prob:omega}\ref{diffs-dual-tate:it:omega-magma}, find the $x$-coordinates of all
        $P \in E(\Qbar)$ with $[2]P \in E[2](\Qbar) \smallsetminus \brk{O}$, and conclude that
        $\#E[4](\Qbar) = 16$.
\begin{solution}
% Write your solution here.
\end{solution}
  \item\label{diffs-dual-tate:it:tate-basis} Show that $P_2 \coloneqq (i, 1 - i)$ and
        $Q_2 \coloneqq (1 + \sqrt2, 2 + \sqrt2)$ are points of $E(\Qbar)$, with $[2]P_2 = (0, 0)$ and
        $[2]Q_2 = (1, 0)$. Show that $(P_2, Q_2)$ is a basis of $E[4](\Qbar)$ over $\Z/4\Z$, and
        that $E[4](\Qbar) \subseteq E(\Q(\zeta_8))$. Compare with $\mu_8(\Qbar)$ and
        \cref{diffs-dual-tate:prop:weil}.
\begin{solution}
% Write your solution here.
\end{solution}
  \item\label{diffs-dual-tate:it:tate-conj} Show that $\sigma$ restricts to a group automorphism of
        $E[4](\Qbar)$. Compute $P_2 + (1, 0)$, and write $\sigma(P_2)$ and $\sigma(Q_2)$ in
        the basis $(P_2, Q_2)$. Compare the determinant of the resulting matrix in
        $\GL_2(\Z/4\Z)$ with the action of $\sigma$ on $\mu_4(\Qbar)$ found in
        \ref{diffs-dual-tate:it:tate-gm}, and explain the agreement using \cref{diffs-dual-tate:prop:weil}.
\begin{solution}
% Write your solution here.
\end{solution}
  \item\label{diffs-dual-tate:it:tate-module} Let $(P_n)_n$ and $(Q_n)_n$ be compatible sequences, $[2]P_{n+1} = P_n$ and
        $[2]Q_{n+1} = Q_n$, starting with $P_1 = (0,0)$, $Q_1 = (1,0)$ and $P_2$, $Q_2$ as
        above. Explain why they form a $\Z_2$-basis of $T_2(E)$, and compare with
        \ref{diffs-dual-tate:it:tate-gm}.
\begin{solution}
% Write your solution here.
\end{solution}
\end{enumerate}
\end{problem}

\section{\texorpdfstring{\Magma}{Magma} examples}

\begin{wsexample}[The isogenies of \cref{diffs-dual-tate:prob:dual}]
\label{diffs-dual-tate:ex:magma-dual}
We check in the function field of $E_1$ that $\varphi$ and $\psi$ land where they should,
that $\iota\circ\psi\circ\varphi = [2]$, and compare with \Magma's own $2$-isogeny.
\begin{magma}
E := EllipticCurve([-1,0]);                  // y^2 = x^3 - x
Kx<x> := FunctionField(Rationals());
R<t> := PolynomialRing(Kx);
F<y> := ext<Kx | t^2 - (x^3 - x)>;                      // x = x_1, y = y_1
x2 := F!((x^2 - 1)/x);  y2 := y*(x^2 + 1)/x^2;          // phi^*x_2, phi^*y_2
y2^2 eq x2^3 + 4*x2;                                    // true
x3 := y2^2/x2^2;  y3 := y2*(x2^2 - 4)/x2^2;             // (psi o phi)^*x_3, ^*y_3
y3^2 eq x3^3 - 16*x3;                                   // true
lam := (3*x^2 - 1)/(2*y);  xd := lam^2 - 2*x;  yd := -(y + lam*(xd - x));
x3/4 eq xd and y3/8 eq yd;                              // true: [2]
phi := TwoIsogeny(E![0,0]);
aInvariants(Codomain(phi));                             // [ 0, 0, 0, 4, 0 ]
phi(E![1,0]);                                           // (0 : 0 : 1)
DualIsogeny(phi);
// ... taking (x : y : 1) to ((1/4*x^2 + 1) / x : (1/8*x^2*y - 1/2*y) / x^2 : 1)
\end{magma}
\Magma's dual is $\iota\circ\psi$ written out in the coordinates $x_2, y_2$ of $E_2$:
$\paren{(x_2^2 + 4)/(4x_2),\ y_2(x_2^2 - 4)/(8x_2^2)}$.
\end{wsexample}

\begin{wsexample}[Division polynomials and {$E[4](\Qbar)$}]
\label{diffs-dual-tate:ex:magma-e4}
For $m \geq 2$, \mg{DivisionPolynomial(E,m)} vanishes exactly at the $x$-coordinates of
$E[m](\Qbar) \smallsetminus \brk{O}$. Its degree is $(m^2 - 1)/2$ for odd $m$ and
$(m^2 + 2)/2$ for even $m$, and its factors for $m = 4$ are $x$, $x \pm 1$ and the
$g_e$ of \cref{diffs-dual-tate:prob:omega}\ref{diffs-dual-tate:it:omega-magma}.
\begin{magma}
[ Degree(DivisionPolynomial(E,m)) : m in [2..8] ];
// [ 3, 4, 9, 12, 19, 24, 33 ]
Factorization(DivisionPolynomial(E,4));
// x - 1, x, x + 1, x^2 - 2*x - 1, x^2 + 1, x^2 + 2*x - 1
Q8<z> := CyclotomicField(8);  i := z^2;  s := z + z^7;  // s^2 = 2
E8 := BaseChange(E, Q8);
P2 := E8![i, 1-i];  Q2 := E8![1+s, 2+s];
2*P2, 2*Q2;                        // (0 : 0 : 1) (1 : 0 : 1)
#DivisionPoints(E8!0, 4);          // 16
E8![-i, 1+i] eq 3*P2 + 2*Q2;       // true: complex conjugation
\end{magma}
\end{wsexample}

\begin{wsexample}[Characteristic $p$]
\label{diffs-dual-tate:ex:magma-char-p}
Reducing modulo $p$, the $p$-division polynomial collapses, because
$\#E[p](\Fbar_p) \in \brk{1, p}$.
\begin{magma}
for p in [3,5,7,13] do
  Ep  := ChangeRing(E, GF(p));
  psi := DivisionPolynomial(Ep, p);
  print p, IsSupersingular(Ep), Degree(psi), Factorization(psi);
end for;
//  3 true   0  []
//  5 false 10  [ <x^2 + 2, 5> ]
//  7 true   0  []
// 13 false 78  [ <x^6 + 11*x^4 + 5*x^2 + 11, 13> ]
\end{magma}
For $p = 3, 7$ the curve is supersingular and $E[p](\Fbar_p) = 0$. In general,
$y^2 = x^3 - x$ is supersingular exactly when $p \equiv 3 \pmod 4$. For $p = 5, 13$
the division polynomial is the $p$-th power of a polynomial of degree $(p - 1)/2$, whose
roots are the $x$-coordinates of the $p - 1$ non-zero points of
$E[p](\Fbar_p) \cong \Z/p\Z$. The exponent $p$ is the inseparable degree of $[p]$.
The same computation for $\Gm$ is $t^p - 1 = (t - 1)^p$.
\end{wsexample}

\subsection*{Looking ahead}
The Weil pairing of \cref{diffs-dual-tate:prop:weil} also gives $\det\varphi_\ell = \deg\varphi$ for the map
$\varphi_\ell$ induced by an isogeny on Tate modules
(\cite{Silverman09}*{Proposition~III.8.6}). Formal groups make $c_\varphi$ an honest
derivative. And Frobenius factors $[p]$ in characteristic $p$.

\bigskip

\part[The zeta function of an elliptic curve]{The zeta function of an elliptic curve}
\label{zeta:part}
\noindent
\Cref{zeta:sec:euler} defines the zeta function of a curve over $\F_q$ as a generating
function for its numbers of rational points, and shows that it is also an Euler product
over the closed points of the curve. \Cref{zeta:sec:zeta} computes the zeta function of an
elliptic curve $E$ from the action of Frobenius on the Tate module, proves the Riemann
hypothesis for it, and deduces Hasse's bound
$\lvert \#E(\F_q) - q - 1\rvert \leq 2\sqrt q$. \Cref{zeta:sec:isogeny} shows that the zeta
function, equivalently the single number $\#E(\F_q)$, is a complete isogeny invariant.

\subsection*{Reading}
\cite{Silverman09}*{\S V.2}, with \cite{Silverman09}*{\S\S III.7--III.8} for the Tate
module, and \cite{Poonen06}*{\S\S 2.2.1 and 3.4--3.5} for closed points and the Euler
product. Three of the problems are exercises from \cite{Silverman09}*{Chapter~V}:
\cref{zeta:prob:tracedet} is Exercise~5.3, \cref{zeta:prob:rec} is Exercise~5.13 and
\cref{zeta:prob:isog} is Exercise~5.4.

\subsection*{\texorpdfstring{\Magma}{Magma}}
If you have no local installation, every block below runs in the free
\href{http://magma.maths.usyd.edu.au/calc/}{online calculator}.

\subsection*{Standing notation}
\begin{itemize}
  \item A power $q = p^r$ of a prime $p$, and the Galois group
        $\Gal_{\F_q} = \Gal(\Fbar_q/\F_q)$.
  \item The automorphism $\sigma \in \Gal_{\F_q}$, $t \mapsto t^q$, of $\Fbar_q$, which
        generates $\Gal_{\F_q}$ topologically. Its fixed field in $\Fbar_q$ is $\F_q$,
        and that of $\sigma^n$ is $\F_{q^n}$.
  \item In \cref{zeta:sec:zeta,zeta:sec:isogeny}, an elliptic curve $E$ over $\F_q$, given by a
        Weierstrass equation with coefficients in $\F_q$.
  \item The ring $\End(E)$ of the endomorphisms of $E$ \emph{defined over $\F_q$}.
  \item The $q$-power Frobenius endomorphism
        \[
          \Frob\colon E \to E, \qquad (x, y) \mapsto (x^q, y^q).
        \]
  \item The integers
        \[
          a \coloneqq q + 1 - \#E(\F_q), \qquad a_n \coloneqq q^n + 1 - \#E(\F_{q^n})
          \quad (n \geq 1), \qquad a_0 \coloneqq 2,
        \]
        so that $a_1 = a$.
  \item For a prime $\ell \neq p$, the Tate module
        $T_\ell(E) = \varprojlim_n E[\ell^n](\Fbar_q)$, a free $\Z_\ell$-module of
        rank~$2$, and the map $\psi_\ell \in \End\bigl(T_\ell(E)\bigr)$ induced by
        $\psi \in \End(E)$.
\end{itemize}

\section{Zeta functions and closed points}
\label{zeta:sec:euler}

Let $C$ be a smooth projective curve over $\F_q$. The zeta function of $C$
(\cref{zeta:def:zeta}) packages the numbers of rational points over all the extensions of
$\F_q$ into one power series. The exponential in it looks odd: the more obvious series is
$\sum_n \#C(\F_{q^n})T^n$. \Cref{zeta:prob:euler} explains the choice. It is
\cite{Poonen06}*{Proposition~3.4.3 and Exercise~3.7}. In \cite{Poonen06} the ground field
is perfect throughout, as $\F_q$ is.

\begin{wsdefinition}[{Closed points \cite{Poonen06}*{Definition~2.2.1}}]
\label{zeta:def:closed}
A \defi{closed point} $P$ of $C$ is an integral subvariety of dimension $0$. Its
\defi{residue field} $\F_q(P)$ is a finite extension of $\F_q$, and its \defi{degree} is
$\deg P \coloneqq [\F_q(P) : \F_q]$.
\end{wsdefinition}

\begin{wsdefinition}[The zeta function]
\label{zeta:def:zeta}
The \defi{zeta function} of $C$ is
\[
  Z(C/\F_q; T) \coloneqq \exp\Bigl( \sum_{n \geq 1} \#C(\F_{q^n}) \frac{T^n}{n} \Bigr)
  \in \Q[[T]],
\]
so that $\#C(\F_{q^n})$ is the coefficient of $T^{n-1}$ in
$\frac{d}{dT}\log Z(C/\F_q; T)$.
\end{wsdefinition}

\begin{wsproposition}[{Closed points are Galois orbits \cite{Poonen06}*{Remark~2.2.2}}]
\label{zeta:prop:closed}
Sending a point of $C(\Fbar_q)$ to its image is a surjection from $C(\Fbar_q)$ onto the
set of closed points of $C$, whose fibers are the $\Gal_{\F_q}$-orbits. The orbit over a
closed point $P$ has $\deg P$ elements: a point of $C(\Fbar_q)$ over $P$ is an
$\F_q$-embedding $\F_q(P) \hookrightarrow \Fbar_q$, and there are $\deg P$ of these.
\end{wsproposition}

\begin{problem}[The Euler product]
\label{zeta:prob:euler}
For $d \geq 1$, let $N_d$
be the number of closed points of $C$ of degree $d$.
\begin{enumerate}
  \item\label{zeta:it:euler-orbit} Show that a point of $C(\Fbar_q)$ lies in $C(\F_{q^n})$ if
        and only if $\sigma^n$ fixes it. Deduce that a $\Gal_{\F_q}$-orbit of size $d$ is
        contained in $C(\F_{q^n})$ if $d \mid n$, and is disjoint from it otherwise.
\begin{solution}
% Write your solution here.
\end{solution}
  \item\label{zeta:it:euler-count} Show that $N_d$ is finite, and that
        \[
          \#C(\F_{q^n}) = \sum_{d \mid n} d N_d \qquad \text{for every } n \geq 1 .
        \]
\begin{solution}
% Write your solution here.
\end{solution}
  \item\label{zeta:it:euler-product} Deduce the identity of formal power series
        \[
          Z(C/\F_q; T) = \prod_{P} \bigl(1 - T^{\deg P}\bigr)^{-1}
          = \prod_{d \geq 1} (1 - T^d)^{-N_d},
        \]
        the first product running over the closed points $P$ of $C$. Why is the product
        well defined? Conclude that $Z(C/\F_q; T) \in 1 + T\Z[[T]]$, which is not
        obvious from the exponential.
\begin{solution}
% Write your solution here.
\end{solution}
  \item\label{zeta:it:euler-mobius} Invert \ref{zeta:it:euler-count}: show that
        $N_d = \frac1d \sum_{e \mid d} \mu(d/e)\,\#C(\F_{q^e})$, where $\mu$ is the Möbius
        function. So the numbers of rational points over the extensions of $\F_q$ and the
        numbers of closed points of each degree determine each other.
\begin{solution}
% Write your solution here.
\end{solution}
  \item\label{zeta:it:euler-p1} Take $C = \PP^1$. Show that
        $Z(\PP^1/\F_q; T) = 1/\bigl((1 - T)(1 - qT)\bigr)$. The closed points of
        $\PP^1$ are the point at infinity, of degree $1$, and the closed points of
        $\AA^1$, which are the monic irreducible polynomials $f \in \F_q[t]$, with
        $\deg P = \deg f$. Deduce that the number of monic irreducible polynomials of
        degree $d$ over $\F_q$ is $\frac1d\sum_{e \mid d}\mu(d/e)q^e$. How many are there
        of degree $2$ over $\F_2$, and which are they?
\begin{solution}
% Write your solution here.
\end{solution}
  \item\label{zeta:it:euler-s} Put $\zeta_C(s) \coloneqq Z(C/\F_q; q^{-s})$ and
        $NP \coloneqq \#\F_q(P)$. Show that
        $\zeta_C(s) = \prod_P \bigl(1 - NP^{-s}\bigr)^{-1}$, and compare with the Euler
        product $\zeta(s) = \prod_p (1 - p^{-s})^{-1}$ of the Riemann zeta function.
\begin{solution}
% Write your solution here.
\end{solution}
\end{enumerate}
\end{problem}

\begin{wsremark}[The Weil conjectures for curves]
\label{zeta:rem:weil}
Let $C$ be a smooth projective curve of genus $g$ over $\F_q$. The Weil conjectures
(proved for curves by Weil) say:
\begin{enumerate}
  \item\label{zeta:it:weil-rational} (Rationality) $Z(C/\F_q; T) = P(T)/\bigl((1-T)(1-qT)\bigr)$
        with $P \in \Z[T]$ of degree $2g$ and $P(0) = 1$.
  \item\label{zeta:it:weil-fe} (Functional equation)
        $Z\bigl(C/\F_q; 1/(qT)\bigr) = q^{1-g}T^{2-2g}Z(C/\F_q; T)$.
  \item\label{zeta:it:weil-rh} (Riemann hypothesis) $P(T) = \prod_{i=1}^{2g}(1 - \alpha_i T)$
        with $\lvert\alpha_i\rvert = \sqrt q$ for every $i$.
\end{enumerate}
For $C = \PP^1$, of genus $0$, \cref{zeta:prob:euler}\ref{zeta:it:euler-p1} gives all three with
$P = 1$. We prove them for $g = 1$ in \cref{zeta:sec:zeta}.
\end{wsremark}

\section{The zeta function of an elliptic curve}
\label{zeta:sec:zeta}

\Cref{zeta:prob:tracedet} is \cite{Silverman09}*{Exercise~5.3}.

\begin{problem}[Traces and determinants]
\label{zeta:prob:tracedet}
Let $A$ be a square matrix over a field $K$
of characteristic $0$. We prove, in two ways, that
\begin{equation}
\label{zeta:eq:tracedet}
  \exp\Bigl(\sum_{n \geq 1} \tr(A^n)\frac{T^n}{n}\Bigr) = \frac{1}{\det(1 - AT)}
  \qquad \text{in } K[[T]].
\end{equation}
\begin{enumerate}
  \item\label{zeta:it:tracedet-lemma} Let $f, g \in 1 + TK[[T]]$. Show that $f$ and $g$ are
        units of $K[[T]]$, so that their \defi{logarithmic derivatives} $f'/f$ and
        $g'/g$ are defined, and that if $f'/f = g'/g$, then $f = g$.\hint{compute the
        derivative of $f/g$.} Where do you use that $K$ has characteristic $0$? Give a
        counterexample in characteristic $p$.
\begin{solution}
% Write your solution here.
\end{solution}
  \item\label{zeta:it:tracedet-logder} Show that both sides of \eqref{zeta:eq:tracedet} lie in
        $1 + TK[[T]]$, compute their logarithmic derivatives, and conclude with
        \ref{zeta:it:tracedet-lemma}.\hint{for $M(T)$ an invertible matrix over
        $K[[T]]$, $\frac{d}{dT}\det M(T) = \det M(T)\,\tr\bigl(M(T)^{-1}M'(T)\bigr)$.}
\begin{solution}
% Write your solution here.
\end{solution}
  \item\label{zeta:it:tracedet-triangular} Give a second proof: triangularize $A$ over an
        algebraic closure $\bar K$ of $K$, reduce \eqref{zeta:eq:tracedet} to the case of a
        $1 \times 1$ matrix $A = (\lambda)$, and prove that case with
        \ref{zeta:it:tracedet-lemma}, applied over $\bar K$.
\begin{solution}
% Write your solution here.
\end{solution}
\end{enumerate}
\end{problem}

\Cref{zeta:prob:zeta} is the proof of \cite{Silverman09}*{Theorems~V.2.3.1 and V.2.4},
arranged around \cref{zeta:prob:tracedet}.

In the following results, $E$ is an elliptic curve over $\F_q$ and $\ell \neq p$ is a
prime.

\begin{wsproposition}[{Frobenius is purely inseparable \cite{Silverman09}*{Proposition~II.2.11}}]
\label{zeta:prop:deg}
The Frobenius endomorphism $\Frob$ is purely inseparable of degree $q$.
\end{wsproposition}

\begin{wscorollary}[{$1 - \Frob$ is separable \cite{Silverman09}*{Corollary~III.5.5}}]
\label{zeta:cor:sep}
The endomorphism $1 - \Frob$ is separable.
\end{wscorollary}

\begin{wstheorem}[{Kernels of separable isogenies \cite{Silverman09}*{Theorem~III.4.10(c)}}]
\label{zeta:thm:ker}
A nonzero separable isogeny $\psi$ satisfies
$\#\ker\psi = \deg\psi$, where $\ker \psi$ is taken in $E(\Fbar_q)$.
\end{wstheorem}

\begin{wsproposition}[{Determinant and trace on $T_\ell(E)$ \cite{Silverman09}*{Proposition~III.8.6}}]
\label{zeta:prop:trace}
For $\psi \in \End(E)$,
$\det \psi_\ell = \deg\psi$ and $\tr\psi_\ell = 1 + \deg\psi - \deg(1 - \psi)$.
\end{wsproposition}

\begin{wstheorem}[{Endomorphisms act faithfully on $T_\ell(E)$ \cite{Silverman09}*{Theorem~III.7.4}}]
\label{zeta:thm:inj}
The map $\psi \mapsto \psi_\ell$ is an injective ring
homomorphism $\End(E) \to \End\bigl(T_\ell(E)\bigr)$.
\end{wstheorem}

\begin{problem}[The zeta function of $E$]
\label{zeta:prob:zeta}
\begin{enumerate}
  \item\label{zeta:it:zeta-frob} Show that $\sigma(Q) = \Frob(Q)$ for every
        $Q \in E(\Fbar_q)$. Deduce that the closed points of $E$ of degree $d$ are the
        $\Frob$-orbits of size $d$ in $E(\Fbar_q)$, and that
        $E(\F_{q^n}) = \ker(1 - \Frob^n)$.
\begin{solution}
% Write your solution here.
\end{solution}
  \item\label{zeta:it:zeta-count} Show that $\Frob^n$ is the $q^n$-power Frobenius
        endomorphism of $E$ over $\F_{q^n}$, and deduce that
        \[
          \#E(\F_{q^n}) = \deg(1 - \Frob^n) = \det(1 - \Frob_\ell^n)
          = 1 - \tr(\Frob_\ell^n) + q^n .
        \]
        Where does the argument use that $1 - \Frob^n$ is separable?
\begin{solution}
% Write your solution here.
\end{solution}
  \item\label{zeta:it:zeta-charpoly} Taking $n = 1$, show that the characteristic polynomial
        of $\Frob_\ell$ is $T^2 - aT + q$. In particular it has integer coefficients and
        does not depend on $\ell$.
\begin{solution}
% Write your solution here.
\end{solution}
  \item\label{zeta:it:zeta-rational} Use \cref{zeta:prob:tracedet} to show that
        \[
          Z(E/\F_q; T) = \frac{\det(1 - \Frob_\ell T)}{(1 - T)(1 - qT)}
          = \frac{1 - aT + qT^2}{(1 - T)(1 - qT)} .
        \]
        By \cref{zeta:prob:euler}, this is also $\prod_P(1 - T^{\deg P})^{-1}$ over the
        closed points $P$ of $E$.
\begin{solution}
% Write your solution here.
\end{solution}
  \item\label{zeta:it:zeta-fe} Verify the functional equation
        $Z\bigl(E/\F_q; 1/(qT)\bigr) = Z(E/\F_q; T)$ of \cref{zeta:rem:weil}.
\begin{solution}
% Write your solution here.
\end{solution}
\end{enumerate}
\end{problem}

\begin{problem}[The Riemann hypothesis and Hasse's bound]
\label{zeta:prob:rh}
Let $\alpha, \beta \in \C$ be the roots of $T^2 - aT + q$, so that
$1 - aT + qT^2 = (1 - \alpha T)(1 - \beta T)$.
\begin{enumerate}
  \item\label{zeta:it:rh-rh} For $r, s \in \Z$ with $s \neq 0$, show that
        $\det(r/s - \Frob_\ell) = \deg(r - s\Frob)/s^2 \geq 0$. Deduce that
        $T^2 - aT + q$ is nonnegative on $\R$, and hence that $\beta = \bar\alpha$ and
        $\lvert\alpha\rvert = \lvert\beta\rvert = \sqrt q$. This is the Riemann
        hypothesis of \cref{zeta:rem:weil}. With $\zeta_E$ as in
        \cref{zeta:prob:euler}\ref{zeta:it:euler-s}, show that every zero of $\zeta_E(s)$ satisfies
        $\operatorname{Re}(s) = 1/2$, and that $\zeta_E(s) = \zeta_E(1 - s)$.
\begin{solution}
% Write your solution here.
\end{solution}
  \item\label{zeta:it:rh-hasse} Show that $a_n = \alpha^n + \beta^n$ for every $n \geq 0$.
        Deduce Hasse's bound
        \[
          \lvert \#E(\F_q) - q - 1\rvert \leq 2\sqrt q ,
        \]
        and more generally $\lvert a_n\rvert \leq 2q^{n/2}$ for every $n \geq 1$.
\begin{solution}
% Write your solution here.
\end{solution}
  \item\label{zeta:it:rh-converse} Conversely, suppose only that there is a constant $c$ with
        $\lvert a_n\rvert \leq c\,q^{n/2}$ for all $n \geq 1$. Show that the power series
        $\sum_{n\geq1}a_nT^n = \frac{\alpha T}{1 - \alpha T} + \frac{\beta T}{1 - \beta T}$
        converges for $\lvert T\rvert < q^{-1/2}$, and deduce that
        $\lvert\alpha\rvert = \lvert\beta\rvert = \sqrt q$. So the Riemann hypothesis is
        equivalent to a Hasse bound over every extension of $\F_q$, even with a worse
        constant.
\begin{solution}
% Write your solution here.
\end{solution}
\end{enumerate}
\end{problem}

\Cref{zeta:prob:rec} is \cite{Silverman09}*{Exercise~5.13}.

\begin{problem}[A recurrence]
\label{zeta:prob:rec}
\begin{enumerate}
  \item\label{zeta:it:rec-one} Show that the single number $\#E(\F_q)$ determines
        $\#E(\F_{q^n})$ for every $n \geq 1$.
\begin{solution}
% Write your solution here.
\end{solution}
  \item\label{zeta:it:rec-rec} Prove that
        \[
          a_{n+2} = a_1 a_{n+1} - q a_n \qquad \text{for all } n \geq 0 .
        \]
        (Recall the convention $a_0 = 2$.)
\begin{solution}
% Write your solution here.
\end{solution}
  \item\label{zeta:it:rec-magma} Here is the recurrence on $y^2 = x^3 + x + 1$ over $\F_{101}$:
\begin{magma}
q := 101;  E := EllipticCurve([GF(q) | 1, 1]);          // y^2 = x^3 + x + 1
a := [TraceOfFrobenius(ChangeRing(E, GF(q^n))) : n in [1..6]];
a;                                // [ -3, -193, 882, 16847, -139623, -1282678 ]
b := [2, a[1]];
for n in [1..5] do Append(~b, a[1]*b[n+1] - q*b[n]); end for;
b[2..7] eq a;                     // true
\end{magma}
        Here \mg{TraceOfFrobenius(E)} is $a$. Compute $a_2$ by hand from $a_1 = -3$.
        Check Hasse's bound for $a_6$ over $\F_{101^6}$.
\begin{solution}
% Write your solution here.
\end{solution}
\end{enumerate}
\end{problem}

\begin{problem}[A zeta function by hand]
\label{zeta:prob:byhand}
Let $E$ be the elliptic curve $y^2 + xy = x^3 + 1$ over $\F_2$ (its discriminant
is $1$).
\begin{enumerate}
  \item\label{zeta:it:byhand-count} List the rational points of $E(\F_2)$, and find $a$.
\begin{solution}
% Write your solution here.
\end{solution}
  \item\label{zeta:it:byhand-zeta} Write down $Z(E/\F_2; T)$, and check the Riemann hypothesis
        for it directly.
\begin{solution}
% Write your solution here.
\end{solution}
  \item\label{zeta:it:byhand-ext} Use \cref{zeta:prob:rec} to predict $\#E(\F_4)$ and $\#E(\F_8)$.
        Check the first by listing the rational points of $E(\F_4)$, where
        $\F_4 = \brk{0, 1, \omega, \omega^2}$ with $\omega^2 = \omega + 1$.\hint{for $x \neq 0$, substitute $y = xz$.}
\begin{solution}
% Write your solution here.
\end{solution}
  \item\label{zeta:it:byhand-closed} Use \cref{zeta:prob:euler}\ref{zeta:it:euler-mobius} to find the
        numbers $N_1$, $N_2$, $N_3$ of closed points of $E$ of degree $1$, $2$ and $3$.
        Deduce that every $\F_8$-rational point of $E$ is $\F_2$-rational. Is
        $\#E(\F_8)$ close to the edge of the Hasse interval?
\begin{solution}
% Write your solution here.
\end{solution}
  \item\label{zeta:it:byhand-euler} Check the Euler product of
        \cref{zeta:prob:euler}\ref{zeta:it:euler-product} against \ref{zeta:it:byhand-zeta} up to the
        term in $T^2$.
\begin{solution}
% Write your solution here.
\end{solution}
  \item\label{zeta:it:byhand-magma} Check, where \mg{Places(E, d)} lists the closed points of
        degree $d$:
\begin{magma}
E := EllipticCurve([GF(2) | 1, 0, 0, 0, 1]);            // y^2 + xy = x^3 + 1
[#ChangeRing(E, GF(2^n)) : n in [1..8]];     // [ 4, 8, 4, 16, 44, 56, 116, 288 ]
[#Places(E, d) : d in [1..6]];               // [ 4, 2, 0, 2, 8, 8 ]
ZetaFunction(E);              // (2*$.1^2 + $.1 + 1)/(2*$.1^2 - 3*$.1 + 1)
\end{magma}
\begin{solution}
% Write your solution here.
\end{solution}
\end{enumerate}
\end{problem}

\section{Isogenous curves}
\label{zeta:sec:isogeny}

\Cref{zeta:prob:isog} is \cite{Silverman09}*{Exercise~5.4}. Let $E_1$ and $E_2$ be elliptic
curves over $\F_q$, with $q$-power Frobenius endomorphisms $\Frob_1$ and $\Frob_2$ and
traces $a(E_1)$, $a(E_2)$.

\begin{wsdefinition}[Isogenies and Galois-equivariant maps]
\label{zeta:def:hom}
We write $\Hom(E_1, E_2)$ for the isogenies $E_1 \to E_2$ \emph{defined over $\F_q$}. The
Galois group $\Gal_{\F_q}$ acts on each Tate module $T_\ell(E_i)$, and we write
$\Hom_{\Gal_{\F_q}}\bigl(T_\ell(E_1), T_\ell(E_2)\bigr)$ for the $\Z_\ell$-linear maps
that commute with this action.
\end{wsdefinition}

\begin{wstheorem}[{Tate's isogeny theorem \cite{Silverman09}*{Theorem~III.7.7(a)}}]
\label{zeta:thm:tate}
Let $\ell \neq p$ be prime. The natural map
\[
  \Hom(E_1, E_2) \otimes \Z_\ell \to \Hom_{\Gal_{\F_q}}\bigl(T_\ell(E_1), T_\ell(E_2)\bigr)
\]
is an isomorphism.
\end{wstheorem}

\begin{wscorollary}[{Isogenies form a free module \cite{Silverman09}*{Corollary~III.7.5}}]
\label{zeta:cor:hom-free}
The group $\Hom(E_1, E_2)$ is a free $\Z$-module of finite rank, being a subgroup of the
free module of.
\end{wscorollary}

\begin{problem}[Isogenous curves have the same number of rational points]
\label{zeta:prob:isog}
\begin{enumerate}
  \item\label{zeta:it:isog-count} Let $\varphi\colon E_1 \to E_2$ be a nonzero isogeny defined
        over $\F_q$. Show that $\varphi \circ \Frob_1 = \Frob_2 \circ \varphi$, and deduce
        that $\#E_1(\F_q) = \#E_2(\F_q)$.
\begin{solution}
% Write your solution here.
\end{solution}
  \item\label{zeta:it:isog-zeta} Deduce that $Z(E_1/\F_q; T) = Z(E_2/\F_q; T)$.
\begin{solution}
% Write your solution here.
\end{solution}
\end{enumerate}
For the converse, suppose that $\#E_1(\F_q) = \#E_2(\F_q)$, so that
$\Frob_{1,\ell}$ and $\Frob_{2,\ell}$ have the same characteristic polynomial
$f(T) = T^2 - aT + q$ (\cref{zeta:prob:zeta}\ref{zeta:it:zeta-charpoly}). Fix a prime
$\ell \neq p$ and put $V_i \coloneqq T_\ell(E_i) \otimes \Q_\ell$.
\begin{enumerate}[resume]
  \item\label{zeta:it:isog-galois} Show that a $\Z_\ell$-linear map
        $T_\ell(E_1) \to T_\ell(E_2)$ commutes with the action of $\Gal_{\F_q}$ if and only
        if it intertwines $\Frob_{1,\ell}$ and $\Frob_{2,\ell}$.
\begin{solution}
% Write your solution here.
\end{solution}
  \item\label{zeta:it:isog-cyclic} Suppose $a^2 \neq 4q$. Show that $\Frob_{1,\ell}$ is not a
        scalar, that $V_1$ has a basis $v, \Frob_{1,\ell}v$, and that $V_1 \cong
        \Q_\ell[T]/(f)$ as $\Q_\ell[T]$-modules, with $T$ acting by $\Frob_{1,\ell}$.
        Deduce that there is a nonzero $\Z_\ell$-linear map
        $T_\ell(E_1) \to T_\ell(E_2)$ intertwining the two Frobenius maps.
\begin{solution}
% Write your solution here.
\end{solution}
  \item\label{zeta:it:isog-scalar} Suppose $a^2 = 4q$. Show that $q$ is a square, that
        $a/2 \in \Z$, and that $\Frob_i = [a/2]$ in $\End(E_i)$ for $i = 1, 2$.\hint{compute $\deg(\Frob_i - [a/2])$ with \cref{zeta:prop:trace}.} Deduce that every
        $\Z_\ell$-linear map $T_\ell(E_1) \to T_\ell(E_2)$ intertwines the two Frobenius
        maps.
\begin{solution}
% Write your solution here.
\end{solution}
  \item\label{zeta:it:isog-tate} Use \cref{zeta:thm:tate} to conclude that $E_1$ and $E_2$ are
        isogenous over $\F_q$.
\begin{solution}
% Write your solution here.
\end{solution}
\end{enumerate}
\end{problem}

\begin{wsremark}
\label{zeta:rem:isog-classes}
Together with \cref{zeta:prob:zeta}\ref{zeta:it:zeta-rational}, \cref{zeta:prob:isog} says that for
elliptic curves $E_1$, $E_2$ over $\F_q$ the following are equivalent: $E_1$ and $E_2$ are
isogenous over $\F_q$; $\#E_1(\F_q) = \#E_2(\F_q)$; $Z(E_1/\F_q; T) = Z(E_2/\F_q; T)$;
and $E_1$ and $E_2$ have the same number of closed points of each degree
(\cref{zeta:prob:euler}\ref{zeta:it:euler-mobius}).
\end{wsremark}

\subsection*{Looking ahead}
By \cref{zeta:prob:rh}\ref{zeta:it:rh-rh}, $\alpha = \sqrt q\,e^{i\theta}$ for a unique
$\theta \in [0, \pi]$, and $a = 2\sqrt q\cos\theta$. Fix an elliptic curve over $\Q$,
reduce it modulo every prime $p$ of good reduction, and let $\theta_p$ be the angle of
the reduction over $\F_p$. How are the $\theta_p$ distributed in $[0, \pi]$ as $p$
varies? That is the Sato--Tate conjecture. The product over the primes of good
reduction of the factors $(1 - a_pp^{-s} + p^{1-2s})^{-1}$, the reciprocals of the
numerators of the $\zeta$-functions of the reductions, is, up to finitely many factors, the
$L$-function of the curve, the subject of the Birch and Swinnerton-Dyer conjecture.

\bigskip

\part[Supersingular curves and the trace of Frobenius]{Supersingular curves and the trace of Frobenius}
\label{supersingular:part}
\noindent
\Cref{supersingular:sec:trace} proves that an elliptic curve over $\F_q$ is supersingular if and only
if $p$ divides its trace of Frobenius, and works out what this says over $\F_p$, over
$\F_2$ and $\F_3$, and over $\F_8$. \Cref{supersingular:sec:group} shows that
$E(\F_q) \cong \Z/m\Z \times \Z/mn\Z$ with $m \mid q - 1$, and pins down $m$ for a
supersingular curve over $\F_p$. \Cref{supersingular:sec:ordinary} shows that an elliptic curve over
$\Q$ is ordinary modulo infinitely many primes.

\subsection*{Reading}
\cite{Silverman09}*{\S\S V.3--V.4}, with \cite{Silverman09}*{\S V.2} for the trace of
Frobenius, \cite{Silverman09}*{\S III.6} for the torsion subgroups and
\cite{Silverman09}*{\S III.8} for the Weil pairing. The three problems are exercises
from \cite{Silverman09}*{Chapter~V}: \cref{supersingular:prob:trace} is Exercise~5.10,
\cref{supersingular:prob:group} is Exercise~5.6 and \cref{supersingular:prob:ordinary} is Exercise~5.11.

\subsection*{\texorpdfstring{\Magma}{Magma}}
If you have no local installation, every block below runs in the free
\href{http://magma.maths.usyd.edu.au/calc/}{online calculator}.

\subsection*{Standing notation}
\begin{itemize}
  \item A power $q = p^e$ of a prime $p$.
  \item The $p^r$-power Frobenius morphism $\Frob^r\colon C \to C^{(p^r)}$ of a curve $C$
        over a field of characteristic $p$, for $r \geq 1$, which raises coordinates to
        the $p^r$-th power. For an elliptic curve $E$, its dual isogeny
        $\Versch^r\colon E^{(p^r)} \to E$ is the Verschiebung.
  \item The $q$-power Frobenius \emph{endomorphism} of an elliptic curve $E$ over $\F_q$.
        Since $E^{(q)} = E$, the morphism
        \[
          \Frob^e\colon E \to E, \qquad (x, y) \mapsto (x^q, y^q),
        \]
        is an endomorphism of $E$, with dual $\Versch^e \in \End(E)$. Here $\End(E)$
        consists of the endomorphisms defined over $\F_q$.
  \item The group $E[m](\Fbar_q)$ of $m$-torsion points of $E$ with coordinates in
        $\Fbar_q$, for $m \geq 1$.
  \item The \defi{supersingular} elliptic curves $E$ over $\F_q$: those with
        $E[p^r](\Fbar_q) = 0$ for one, equivalently every, $r \geq 1$
        \cite{Silverman09}*{Theorem~V.3.1 and the definition after it}. The others are
        \defi{ordinary}. Silverman also says that a supersingular curve has ``Hasse
        invariant $0$'' and an ordinary one ``Hasse invariant $1$''. That is the phrasing
        of Exercise~5.11.
\end{itemize}

\section{Supersingularity and the trace of Frobenius}
\label{supersingular:sec:trace}

This section is \cite{Silverman09}*{Exercise~5.10}. Part~\ref{supersingular:it:trace-square} is a
little stronger than \cite{Silverman09}*{Exercise~5.10(e)}, and
part~\ref{supersingular:it:trace-eight} needs the stronger form. By Deuring's theorem
\cite{Silverman09}*{Theorem~V.3.1}, a supersingular curve is one whose Verschiebung is
purely inseparable, whose endomorphism ring over $\Fbar_q$ is an order in a quaternion
algebra, and whose formal group has height $2$. None of these is easy to read off from a
Weierstrass equation. Over a finite field there is a criterion that is: count the
rational points.

\begin{wsdefinition}[Trace of Frobenius]
\label{supersingular:def:trace}
Let $E$ be an elliptic curve over $\F_q$. The \defi{trace of Frobenius} of $E$ is
\[
  a \coloneqq q + 1 - \#E(\F_q).
\]
\end{wsdefinition}

\begin{wsdefinition}[Separable degree]
\label{supersingular:def:sepdeg}
For an isogeny $\psi$, $\deg_s\psi$ denotes its \defi{separable degree}.
\end{wsdefinition}

In the following results, $E$ is an elliptic curve over $\F_q$, with $q = p^e$, and $a$ is
its trace of Frobenius.

\begin{wsproposition}[{Frobenius is purely inseparable \cite{Silverman09}*{Proposition~II.2.11}}]
\label{supersingular:prop:frob}
For $r \geq 1$, the Frobenius morphism $\Frob^r$ is purely inseparable of degree $p^r$.
\end{wsproposition}

\begin{wscorollary}[{Purely inseparable isogenies \cite{Silverman09}*{Corollary~II.2.12}}]
\label{supersingular:cor:factor}
A purely inseparable isogeny $\psi\colon E \to E'$ of
degree $p^r$ factors as $\psi = \lambda \circ \Frob^r$ with
$\lambda\colon E^{(p^r)} \to E'$ an isomorphism.
\end{wscorollary}

\begin{wstheorem}[{Kernels and separable degree \cite{Silverman09}*{Theorem~III.4.10(a)}}]
\label{supersingular:thm:kernel}
A nonzero isogeny $\psi$ satisfies
$\#\ker\psi = \deg_s\psi$, the kernel taken on points with coordinates in
$\Fbar_q$.
\end{wstheorem}

\begin{wscorollary}[{Separability of $[m] + [n]\Frob^e$ \cite{Silverman09}*{Corollary~III.5.5}}]
\label{supersingular:cor:sep}
For $m, n \in \Z$, the endomorphism $[m] + [n]\Frob^e$ is
separable if and only if $p \nmid m$.
\end{wscorollary}

\begin{wstheorem}[{Dual isogenies \cite{Silverman09}*{Theorem~III.6.2(a)}}]
\label{supersingular:thm:dual}
Every isogeny $\psi$ satisfies $\hat\psi \circ \psi = [\deg\psi]$ and
$\psi \circ \hat\psi = [\deg\psi]$. So
$\Versch^e \circ \Frob^e = \Frob^e \circ \Versch^e = [q]$.
\end{wstheorem}

\begin{wscorollary}[{The degree of $[m]$ \cite{Silverman09}*{Corollary~III.6.4(a)}}]
\label{supersingular:cor:degree}
For $m \in \Z$, the multiplication-by-$m$ map has degree $\deg[m] = m^2$.
\end{wscorollary}

\begin{wscorollary}[{The $p$-power torsion \cite{Silverman09}*{Corollary~III.6.4(c)}}]
\label{supersingular:cor:ptorsion}
Either $E[p^r](\Fbar_q) = 0$ for every $r \geq 1$, or
$E[p^r](\Fbar_q) \cong \Z/p^r\Z$ for every $r \geq 1$.
\end{wscorollary}

\begin{wstheorem}[{The characteristic polynomial of Frobenius \cite{Silverman09}*{Theorem~V.2.3.1(b)}}]
\label{supersingular:thm:charpoly}
The Frobenius endomorphism satisfies $(\Frob^e)^2 - [a]\Frob^e + [q] = 0$ in $\End(E)$.
\end{wstheorem}

\begin{wstheorem}[{Hasse's bound \cite{Silverman09}*{Theorem~V.1.1}}]
\label{supersingular:thm:hasse}
The trace of Frobenius satisfies $\lvert a\rvert \leq 2\sqrt q$.
\end{wstheorem}

\begin{problem}[Supersingularity and the trace of Frobenius]
\label{supersingular:prob:trace}
Let $E$ be an elliptic curve over $\F_q$.
\begin{enumerate}
  \item\label{supersingular:it:trace-crit} Show that $\Versch^e = [a] - \Frob^e$ in $\End(E)$, and
        deduce that $\Versch^e$ is separable if and only if $p \nmid a$. Show that
        $\ker\Versch^e$ has $\#E[q](\Fbar_q)$ elements, so that
        $\deg_s\Versch^e \in \brk{1, q}$. Conclude that
        \[
          E \text{ is supersingular} \iff a \equiv 0 \pmod p .
        \]
\begin{solution}
% Write your solution here.
\end{solution}
  \item\label{supersingular:it:trace-prime} Suppose $q = p \geq 5$. Show that $E$ is supersingular if
        and only if $\#E(\F_p) = p + 1$. Where is $p \geq 5$ used?
\begin{solution}
% Write your solution here.
\end{solution}
  \item\label{supersingular:it:trace-three} Every elliptic curve over $\F_3$ is isomorphic to one
        given by $y^2 = x^3 + a_2x^2 + a_4x + a_6$. \Magma\ lists the isomorphism
        classes, with $\#E(\F_3)$, $a$ and whether $E$ is supersingular:
\begin{magma}
K := GF(3);  reps := [];
for a2, a4, a6 in K do
  ok, E := IsEllipticCurve([K | 0, a2, 0, a4, a6]);
  if ok and not exists{F : F in reps | IsIsomorphic(E, F)} then
    Append(~reps, E);
  end if;
end for;
#reps;                                                            // 8
for E in reps do
  printf "%o  %o  %o  %o\n", aInvariants(E), #E,
         TraceOfFrobenius(E), IsSupersingular(E);
end for;
// [ 0, 0, 0, 1, 0 ]  4  0  true
// [ 0, 0, 0, 2, 0 ]  4  0  true
// [ 0, 0, 0, 2, 1 ]  7  -3  true
// [ 0, 0, 0, 2, 2 ]  1  3  true
// [ 0, 1, 0, 0, 1 ]  6  -2  false
// [ 0, 1, 0, 0, 2 ]  3  1  false
// [ 0, 2, 0, 0, 1 ]  5  -1  false
// [ 0, 2, 0, 0, 2 ]  2  2  false
\end{magma}
        Which column decides supersingularity, by part~\ref{supersingular:it:trace-crit}? Count the
        rational points of $y^2 = x^3 - x + 1$ and of $y^2 = x^3 - x - 1$ by hand, and
        explain why part~\ref{supersingular:it:trace-prime} fails for $p = 3$.
\begin{solution}
% Write your solution here.
\end{solution}
  \item\label{supersingular:it:trace-two} The same over $\F_2$, where a Weierstrass equation needs
        all five coefficients:
\begin{magma}
K := GF(2);  reps := [];
for a1, a2, a3, a4, a6 in K do
  ok, E := IsEllipticCurve([K | a1, a2, a3, a4, a6]);
  if ok and not exists{F : F in reps | IsIsomorphic(E, F)} then
    Append(~reps, E);
  end if;
end for;
#reps;                                                            // 5
for E in reps do
  printf "%o  %o  %o  %o\n", aInvariants(E), #E,
         TraceOfFrobenius(E), IsSupersingular(E);
end for;
// [ 0, 0, 1, 0, 0 ]  3  0  true
// [ 0, 0, 1, 1, 0 ]  5  -2  true
// [ 0, 0, 1, 1, 1 ]  1  2  true
// [ 1, 0, 0, 0, 1 ]  4  -1  false
// [ 1, 0, 1, 0, 1 ]  2  1  false
\end{magma}
        Count the rational points of the three supersingular curves by hand. Which
        coefficient decides supersingularity over $\F_2$, and what does the analogue of
        part~\ref{supersingular:it:trace-prime} say for $p = 2$?
\begin{solution}
% Write your solution here.
\end{solution}
  \item\label{supersingular:it:trace-square} Suppose $p \mid a$. Use \cref{supersingular:cor:factor} to write
        $\Versch^e = \lambda \circ \Frob^e$ with $\lambda \in \Aut(E)$, deduce that
        $[a] = ([1] + \lambda) \circ \Frob^e$, and conclude that $q \mid a^2$. So
        \[
          a \equiv 0 \pmod p \iff q \mid a^2 \iff p^{\lceil e/2 \rceil} \mid a .
        \]
\begin{solution}
% Write your solution here.
\end{solution}
  \item\label{supersingular:it:trace-eight} Show that no elliptic curve over $\F_8$ has exactly $7$
        or exactly $11$ rational points. \Magma\ collects the values of $\#E(\F_8)$ over
        all Weierstrass equations, and the corresponding traces:
\begin{magma}
K := GF(8);  counts := {};
for a1, a2, a3, a4, a6 in K do
  ok, E := IsEllipticCurve([a1, a2, a3, a4, a6]);
  if ok then Include(~counts, #E); end if;
end for;
counts;                           // { 4, 5, 6, 8, 9, 10, 12, 13, 14 }
{ 9 - n : n in counts };          // { -5, -4, -3, -1, 0, 1, 3, 4, 5 }
\end{magma}
        Which integers in the Hasse interval are missing, and which traces belong to
        supersingular curves? Predict both from parts \ref{supersingular:it:trace-crit} and
        \ref{supersingular:it:trace-square} before you look.
\begin{solution}
% Write your solution here.
\end{solution}
\end{enumerate}
\end{problem}

\begin{wsremark}[The traces of supersingular curves]
\label{supersingular:rem:waterhouse}
The identity $a^2 = \deg([1] + \lambda)\cdot q$ of \cref{supersingular:prob:trace}\ref{supersingular:it:trace-square}
says more than $q \mid a^2$. Since $a^2 \leq 4q$, the integer $d \coloneqq \deg([1] + \lambda)$
is at most $4$, so a supersingular curve over $\F_q$ has
\[
  a^2 \in \brk{0,\ q,\ 2q,\ 3q,\ 4q},
\]
and only the values with $dq$ a perfect square can occur. Over $\F_8$ these are
$a^2 \in \brk{0, 16}$, as found above; over $\F_p$ with $p \geq 5$ only $a = 0$; over
$\F_3$ also $a^2 = 9 = 3q$, and over $\F_2$ also $a^2 = 4 = 2q$, the cases of parts
\ref{supersingular:it:trace-three} and \ref{supersingular:it:trace-two}. Which of the five values occur, over which
$q$, is a theorem of Waterhouse. We do not need it.
\end{wsremark}

\section{The group structure of \texorpdfstring{$E(\F_q)$}{E(F_q)}}
\label{supersingular:sec:group}

This section is \cite{Silverman09}*{Exercise~5.6}.

In the following results, $E$ is an elliptic curve over a field $K$ and $m \geq 1$ is an
integer with $\Char K \nmid m$.

\begin{wscorollary}[{The $m$-torsion points \cite{Silverman09}*{Corollary~III.6.4(b)}}]
\label{supersingular:cor:torsion}
The group $E[m](\Kbar)$ is isomorphic to $\Z/m\Z \times \Z/m\Z$.
\end{wscorollary}

\begin{wscorollary}[{Rational torsion and roots of unity \cite{Silverman09}*{Corollary~III.8.1.1}}]
\label{supersingular:cor:mu}
If every $m$-torsion point of $E(\Kbar)$ is $K$-rational,
then $K$ contains the $m$-th roots of unity.
\end{wscorollary}

\begin{problem}[The group structure of $E(\F_q)$]
\label{supersingular:prob:group}
Let $E$ be an elliptic curve over $\F_q$.
\begin{enumerate}
  \item\label{supersingular:it:group-structure} Show that there are integers $m, n \geq 1$ with
        $p \nmid m$ and $E(\F_q) \cong \Z/m\Z \times \Z/mn\Z$.\hint{decompose the finite abelian group $E(\F_q)$ into its $\ell$-primary parts, and
        count the cyclic factors of each with \cref{supersingular:cor:torsion} for $\ell \neq p$ and
        \cref{supersingular:cor:ptorsion} for $\ell = p$.}
\begin{solution}
% Write your solution here.
\end{solution}
  \item\label{supersingular:it:group-mu} With $m$ as in part~\ref{supersingular:it:group-structure}, show that
        $q \equiv 1 \pmod m$.
\begin{solution}
% Write your solution here.
\end{solution}
  \item\label{supersingular:it:group-ss} Suppose $q = p \geq 5$ and $E$ is supersingular. Show that
        $m = 1$ or $m = 2$, and that $m = 1$ if $p \equiv 1 \pmod 4$.
\begin{solution}
% Write your solution here.
\end{solution}
  \item\label{supersingular:it:group-magma} Both values of $m$ occur when $p \equiv 3 \pmod 4$. Over
        $\F_7$ the curves $y^2 = x^3 + x$ and $y^2 = x^3 - x$, both with $j = 1728$, are
        supersingular. Here \mg{Invariants(AbelianGroup(E))} returns $[m, mn]$, or $[mn]$
        when $m = 1$:
\begin{magma}
E1 := EllipticCurve([GF(7) | 1, 0]);  E2 := EllipticCurve([GF(7) | -1, 0]);
IsSupersingular(E1), IsSupersingular(E2);                        // true true
Invariants(AbelianGroup(E1));                                    // [ 8 ]
Invariants(AbelianGroup(E2));                                    // [ 2, 4 ]
[ j : j in GF(13) | IsSupersingular(EllipticCurveFromjInvariant(j)) ];  // [ 5 ]
E3 := EllipticCurveFromjInvariant(GF(13)!5);  aInvariants(E3);   // [ 1, 0, 0, 7, 2 ]
Invariants(AbelianGroup(E3));                                    // [ 14 ]
Invariants(AbelianGroup(QuadraticTwist(E3)));                    // [ 14 ]
\end{magma}
        Explain the two values of $m$ over $\F_7$ from the rational points of order
        $2$ of each curve, and the value over $\F_{13}$ from part~\ref{supersingular:it:group-ss}.
\begin{solution}
% Write your solution here.
\end{solution}
\end{enumerate}
\end{problem}

\section{Infinitely many ordinary primes}
\label{supersingular:sec:ordinary}

This section is \cite{Silverman09}*{Exercise~5.11}, with ``ordinary'' for ``Hasse
invariant $1$''. Silverman's exercise asks for infinitely many ordinary primes. The curve
of part~\ref{supersingular:it:ordinary-magma} and its supersingular primes are
\cite{Silverman09}*{Example~V.4.6}. The proof uses results that we have not proved.

\begin{wsdefinition}[Reduction modulo $p$]
\label{supersingular:def:reduction}
Let $E$ be an elliptic curve over $\Q$, given by a Weierstrass equation with
coefficients in $\Z$ and discriminant $\Delta \in \Z$. For a prime $p \nmid \Delta$,
reducing the coefficients modulo $p$ gives a Weierstrass equation over $\F_p$ with
nonzero discriminant, hence an elliptic curve $E_p$ over $\F_p$.
\end{wsdefinition}

\begin{wsdefinition}[Supersingular and ordinary primes]
\label{supersingular:def:ordinary-prime}
A prime $p \nmid \Delta$ is \defi{supersingular} or \defi{ordinary} for $E$ according as
$E_p$ is.
\end{wsdefinition}

\begin{wsdefinition}[Primes that split completely]
\label{supersingular:def:split}
Let $L/\Q$ be a finite Galois extension. A prime $p$ \defi{splits completely} in $L$ if
$p\calO_L$ is a product of $[L : \Q]$ distinct prime ideals of the ring of integers
$\calO_L$.
\end{wsdefinition}

In the following results, $E$, $\Delta$ and $E_p$ are as in \cref{supersingular:def:reduction}, and
$L/\Q$ is a finite Galois extension.

\begin{wsproposition}[{Reduction modulo $p$ \cite{Silverman09}*{Proposition~VII.2.1}}]
\label{supersingular:prop:reduction-map}
Let $p \nmid \Delta$. Every point of $E(\Qp)$ can be written $[X : Y : Z]$ with
$X, Y, Z \in \Zp$ not all divisible by $p$, and reducing the coordinates modulo
$p$ gives a group homomorphism $E(\Qp) \to E_p(\F_p)$.
\end{wsproposition}

\begin{wsproposition}[{Reduction of torsion \cite{Silverman09}*{Proposition~VII.3.1(b)}}]
\label{supersingular:prop:reduction-torsion}
Let $p \nmid \Delta$ and let $m \geq 1$ with $p \nmid m$. The restriction of the
reduction homomorphism $E(\Qp) \to E_p(\F_p)$ to the $m$-torsion points of $E(\Qp)$ is
injective.
\end{wsproposition}

\begin{wsproposition}[{Completions at a completely split prime \cite{Neukirch}}]
\label{supersingular:prop:split-completion}
Let $p$ be a prime that splits completely in $L$. For each prime $\mathfrak p$ of $L$
above $p$ the completion $L_{\mathfrak p}$ is $\Qp$, so that $L$ embeds in $\Qp$.
\end{wsproposition}

\begin{wstheorem}[{Chebotarev density theorem \cite{Neukirch}}]
\label{supersingular:thm:chebotarev}
Infinitely many primes split completely in $L$: they have density $1/[L : \Q]$ among all
primes.
\end{wstheorem}

\begin{problem}[Infinitely many ordinary primes]
\label{supersingular:prob:ordinary}
Fix a prime $\ell \geq 3$, and let $L \coloneqq \Q(E[\ell])$ be the
field generated over $\Q$ by the affine coordinates of the nonzero $\ell$-torsion
points of $E(\Qbar)$.
\begin{enumerate}
  \item\label{supersingular:it:ordinary-galois} Show that $L$ is a finite Galois extension of $\Q$,
        and that $L$ contains the $\ell$-th roots of unity.
\begin{solution}
% Write your solution here.
\end{solution}
  \item\label{supersingular:it:ordinary-torsion} Let $p \nmid \ell\Delta$ be a prime that splits
        completely in $L$. Show that $\ell^2 \mid \#E_p(\F_p)$ and that
        $p \equiv 1 \pmod\ell$.
\begin{solution}
% Write your solution here.
\end{solution}
  \item\label{supersingular:it:ordinary-contradiction} Show that every prime $p \geq 5$ with
        $p \nmid \ell\Delta$ that splits completely in $L$ is ordinary for $E$, and
        conclude. Where did you use $\ell \geq 3$?
\begin{solution}
% Write your solution here.
\end{solution}
  \item\label{supersingular:it:ordinary-magma} Take the curve
        $E\colon y^2 + y = x^3 - x^2 - 10x - 20$, of discriminant $-11^5$, and
        $\ell = 3$. By part~\ref{supersingular:it:ordinary-torsion}, a prime $p \neq 3, 11$ that
        splits completely in $\Q(E[3])$ has all nine $3$-torsion points of
        $E_p(\Fbar_p)$ rational. \Magma\ lists the supersingular primes below $100$ and
        the primes below $3000$ with that property:
\begin{magma}
E := EllipticCurve([0, -1, 1, -10, -20]);          // y^2 + y = x^3 - x^2 - 10x - 20
Discriminant(E);                                    // -161051
good := [p : p in PrimesUpTo(3000) | p ne 11];
[p : p in good | p lt 100 and IsSupersingular(ChangeRing(E, GF(p)))];
// [ 2, 19, 29 ]
full := function(p, l)       // are all the l-torsion points of E_p rational?
  I := Invariants(AbelianGroup(ChangeRing(E, GF(p))));
  return #I eq 2 and I[1] mod l eq 0;
end function;
split := [p : p in good | p ne 3 and full(p, 3)];  split;
// [ 337, 523, 1087, 2437, 2719, 2749 ]
[p mod 3 : p in split];                                    // [ 1, 1, 1, 1, 1, 1 ]
[IsSupersingular(ChangeRing(E, GF(p))) : p in split];
// [ false, false, false, false, false, false ]
[p : p in good | p lt 300 and p ne 5 and full(p, 5)];      // and the same for l = 5:
// [ 31, 41, 61, 71, 101, 131, 151, 181, 191, 211, 241, 251, 271, 281 ]
\end{magma}
        Check the output against parts \ref{supersingular:it:ordinary-torsion} and
        \ref{supersingular:it:ordinary-contradiction}, and compare the list of supersingular primes
        with the list in the example cited before this problem. The last line runs the
        same test for $\ell = 5$: which primes below $300$ appear, and what does that suggest about
        the field $\Q(E[5])$? (The rational point $(5, 5)$ of $E$ has order $5$.)
\begin{solution}
% Write your solution here.
\end{solution}
\end{enumerate}
\end{problem}

\begin{wsremark}[Supersingular primes]
\label{supersingular:rem:supersingular-primes}
The other direction is much harder. For the curve of \cref{supersingular:prob:ordinary}\ref{supersingular:it:ordinary-magma},
D.\,H.~Lehmer found exactly $27$ odd supersingular primes below $31500$
\cite{Silverman09}*{Example~V.4.6}. For an elliptic curve over $\Q$ without complex
multiplication, the supersingular primes have density $0$ (Serre, Elkies), and yet there
are infinitely many of them (Elkies) \cite{Silverman09}*{Theorems~V.4.7 and V.4.9}.
The Lang--Trotter conjecture predicts that about $c\sqrt x/\log x$ of them lie below
$x$ \cite{Silverman09}*{Conjecture~V.4.8}. For a curve with complex multiplication,
such as $y^2 = x^3 + x$, the supersingular primes have density $1/2$
\cite{Silverman09}*{Examples~V.4.4 and V.4.5}.
\end{wsremark}

\bigskip

\part[Bounding torsion by reduction]{Reduction modulo $p$ and a bound on the torsion}
\label{torsion:part}
\noindent
The aim is a \Magma\ script, written from scratch, that bounds the torsion subgroup of an
elliptic curve over $\Q$ using only the local theory of
\cite{Silverman09}*{\S\S VII.1--VII.3}. \Cref{torsion:sec:inject} shows that reduction modulo
an odd prime of good reduction is injective on the rational torsion points,
\cref{torsion:sec:bound} turns this into two bounds, and \cref{torsion:sec:lmfdb} tests them against the
LMFDB. Each problem prints a piece of the script with its output. You prove the
mathematics behind it.

\subsection*{Standing notation}
\begin{itemize}
  \item An elliptic curve $E$ over $\Q$, given by a Weierstrass equation with coefficients
        $a_1, \dots, a_6 \in \Z$ and discriminant $\Delta$.
  \item The torsion subgroup $T \coloneqq E(\Q)_{\tors}$.
  \item A prime $p$ is \defi{good} if $p \nmid \Delta$. Then reducing the equation modulo
        $p$ gives an elliptic curve $E_p$ over $\F_p$.
  \item For an abelian group $G$, $G[m]$ is the subgroup killed by $m$.
  \item A finite abelian group $G$ is isomorphic to $\Z/n_1\Z \times \Z/n_2\Z$ for at most
        one pair of integers $1 \leq n_1 \mid n_2$. When it is, we write $n_1(G)$ and
        $n_2(G)$ for them.
\end{itemize}

\section{Reduction is injective on torsion}
\label{torsion:sec:inject}

Everything in this section is in \cite{Silverman09}*{\S\S VII.1--VII.3}. The curve of
\cref{torsion:prob:inject} is from Application~VII.3.5 there.

In the following definition and results, $p$ is a good prime for $E$.

\begin{wsdefinition}[Reduction modulo $p$]
\label{torsion:def:reduction}
We regard $E(\Q)$ as a subgroup of $E(\Qp)$, and write $P \mapsto \bar P$ for the
reduction map $E(\Qp) \to E_p(\F_p)$ and $E_1(\Qp)$ for its kernel.
\end{wsdefinition}

\begin{wsproposition}[{Minimality \cite{Silverman09}*{Remark~VII.1.1}}]
\label{torsion:prop:minimal}
The equation is minimal at $p$, because $v_p(\Delta) = 0 < 12$. So $E_p$ is the
reduction of $E$ modulo $p$.
\end{wsproposition}

\begin{wsproposition}[{Reduction is a homomorphism \cite{Silverman09}*{Proposition~VII.2.1 and the discussion after it}}]
\label{torsion:prop:hom}
The reduction map $E(\Qp) \to E_p(\F_p)$ is a surjective group homomorphism. Its kernel
$E_1(\Qp)$ consists of $O$ and the $(x, y) \in E(\Qp)$ with $v_p(x) < 0$.
\end{wsproposition}

\begin{wsproposition}[{Torsion prime to $p$ \cite{Silverman09}*{Proposition~VII.3.1}}]
\label{torsion:prop:prime-to-p}
If $p \nmid m$, then $E_1(\Qp)[m] = \{O\}$, and the reduction map above restricts to
an injection $E(\Qp)[m] \hookrightarrow E_p(\F_p)$.
\end{wsproposition}

\begin{wstheorem}[{Torsion of $p$-power order \cite{Silverman09}*{Theorem~VII.3.4(b)}}]
\label{torsion:thm:p-power}
Let $P \in E(\Qp)$ have exact order $p^n$ with $n \geq 1$, and let
$r \coloneqq \lfloor 1/(p^n - p^{n-1}) \rfloor$. Then $p^{2r}x(P) \in \Zp$. In particular
$x(P) \in \Zp$ unless $p = 2$ and $n = 1$.
\end{wstheorem}

\begin{problem}[Injectivity of reduction on torsion]
\label{torsion:prob:inject}
\begin{enumerate}
  \item\label{torsion:it:inject-odd} Let $p \geq 3$ be a good prime. Show that the reduction
        map \[E(\Q) \hookrightarrow E(\Qp) \to E_p(\F_p)\] is injective on $T$. In
        particular $T$ is finite, and isomorphic to a subgroup of
        $E_p(\F_p)$.\hint{a nonzero torsion point in $E_1(\Qp)$ has
        a multiple of prime order in $E_1(\Qp)$.}
\begin{solution}
% Write your solution here.
\end{solution}
  \item\label{torsion:it:inject-two} The prime $2$ is different: even if $p = 2$ is good,
        $T \to E_2(\F_2)$ need not be injective. Take for instance the curve
        $E\colon y^2 + xy = x^3 + 4x + 1$:
\begin{magma}
E := EllipticCurve([1, 0, 0, 4, 1]);            // y^2 + xy = x^3 + 4x + 1
Factorization(Integers()!Discriminant(E));      // [ <5, 2>, <13, 2> ]
P := E![-1/4, 1/8];  2*P eq E!0;                // true
\end{magma}
        So $2$ is a good prime and $P = (-1/4, 1/8)$ is a rational point of order $2$.
        Show that $\bar P = O$ in $E_2(\F_2)$. Which hypothesis of
        \cref{torsion:prop:prime-to-p,torsion:thm:p-power} fails, and what does
        \cref{torsion:prop:prime-to-p} still give at $p = 2$?
\begin{solution}
% Write your solution here.
\end{solution}
\end{enumerate}
\end{problem}

\section{Two bounds}
\label{torsion:sec:bound}

This section is \cite{Silverman09}*{Application~VII.3.2}, and its three curves are
\cite{Silverman09}*{Examples~VII.3.3.1--VII.3.3.3}.

\begin{problem}[The order bound]
\label{torsion:prob:order}
The first piece of the script computes the greatest common divisor $b$ of the orders
$\#E_p(\F_p)$ over the good odd primes $p \leq N$:
\begin{magma}
// b = gcd of #E_p(F_p) over the good odd primes p <= N (there must be one)
OrderBound := function(E, N)
  Delta := Integers()!Discriminant(E);
  good := [p : p in PrimesUpTo(N) | p gt 2 and Delta mod p ne 0];
  return GCD([#ChangeRing(E, GF(p)) : p in good]);
end function;
E1 := EllipticCurve([0, 0, 1, -1, 1]);           // y^2 + y = x^3 - x + 1
Discriminant(E1);                                // -611
[#ChangeRing(E1, GF(p)) : p in [3, 5, 7]];       // [ 1, 8, 6 ]
OrderBound(E1, 30);                              // 1
E2 := EllipticCurve([0, 3]);                     // y^2 = x^3 + 3
Discriminant(E2);                                // -3888
[#ChangeRing(E2, GF(p)) : p in [5, 7, 11, 13]];  // [ 6, 13, 12, 9 ]
OrderBound(E2, 30);                              // 1
\end{magma}
\begin{enumerate}
  \item\label{torsion:it:order-examples} Deduce that both curves
        have trivial torsion, and that $y^2 = x^3 + 3$ has infinitely many
        rational points.
\begin{solution}
% Write your solution here.
\end{solution}
  \item\label{torsion:it:order-model} \mg{OrderBound} skips the primes dividing the
        discriminant of the equation it is given. Compare two equations for the first
        curve, with $N = 7$:
\begin{magma}
E4 := EllipticCurve([0, 0, 27, -81, 729]);       // y^2 + 27y = x^3 - 81x + 729
Discriminant(E1), Discriminant(E4);              // -611 -324710451
IsIsomorphic(E1, E4);                            // true
OrderBound(E1, 7), OrderBound(E4, 7);            // 1 2
\end{magma}
        Here $-324710451 = -3^{12} \cdot 611$. Find a change of variables between the
        two equations. Explain why a non-minimal equation can make the script skip a prime
        it could have used, but never makes it use one it should not.
\begin{solution}
% Write your solution here.
\end{solution}
\end{enumerate}
\end{problem}

\begin{wscorollary}[{The $m$-torsion \cite{Silverman09}*{Corollary~III.6.4}}]
\label{torsion:cor:m-torsion}
Let $E$ be an elliptic curve over a field $K$, and $m \geq 1$. If $\Char K \nmid m$, then
$E[m](\Kbar) \cong (\Z/m\Z)^2$. If $\Char K = p > 0$, then $E[p^e](\Kbar)$ is cyclic for
every $e \geq 1$.
\end{wscorollary}

\begin{problem}[The structure bound]
\label{torsion:prob:structure}
The curve $y^2 = x^3 + x$ shows that the orders are not the whole story. \mg{Invariants}
returns $[n_1(G), n_2(G)]$ for $G = E_p(\F_p)$, omitting $n_1(G)$ when it is $1$, and
returns $[\,]$ when $G$ is trivial:
\begin{magma}
E3 := EllipticCurve([1, 0]);                     // y^2 = x^3 + x
Discriminant(E3);                                // -64
OrderBound(E3, 30);                              // 4
[Invariants(AbelianGroup(ChangeRing(E3, GF(p)))) : p in [3, 5, 7, 11, 13]];
// [ [ 4 ], [ 2, 2 ], [ 8 ], [ 12 ], [ 2, 10 ] ]
\end{magma}
\begin{enumerate}
  \item\label{torsion:it:structure-subgroup} Let $G$ be a finite subgroup of $E(\Kbar)$, for an
        elliptic curve $E$ over a field $K$. Show that $n_1(G)$ and $n_2(G)$ are
        defined, and that $n_2(G)$ is the exponent of $G$. Show that
        $\#E[n_1(G)](\Kbar) = n_1(G)^2$, and that an integer $m \geq 1$ with
        $\#E[m](\Kbar) = m^2$ satisfies $E[m](\Kbar) \subseteq G$ if and only if
        $m \mid n_1(G)$. Deduce that $n_1(H) \mid n_1(G)$ and $n_2(H) \mid n_2(G)$ for
        every subgroup $H \subseteq G$.\hint{for a prime $\ell$, the
        group $G[\ell]$ lies in $E[\ell](\Kbar)$. Use \cref{torsion:cor:m-torsion}.}
\begin{solution}
% Write your solution here.
\end{solution}
  \item\label{torsion:it:structure-magma} Let $b_1$ and $b_2$ be the greatest common divisors of
        $n_1(E_p(\F_p))$ and of $n_2(E_p(\F_p))$ over the good odd primes $p \leq N$. Show
        that $n_1(T) \mid b_1$, $n_2(T) \mid b_2$ and $b_1b_2 \mid b$. This is the second
        piece of the script:
\begin{magma}
// [b1, b2]: n1(T) | b1 and n2(T) | b2
StructureBound := function(E, N)
  Delta := Integers()!Discriminant(E);
  good := [p : p in PrimesUpTo(N) | p gt 2 and Delta mod p ne 0];
  b1 := 0;  b2 := 0;
  for p in good do
    I := Invariants(AbelianGroup(ChangeRing(E, GF(p))));
    b1 := GCD(b1, #I eq 2 select I[1] else 1);
    b2 := GCD(b2, #I eq 0 select 1 else I[#I]);
  end for;
  return [b1, b2];
end function;
StructureBound(E3, 30);                          // [ 1, 2 ]
P := E3![0, 0];  2*P eq E3!0;                    // true
\end{magma}
        Conclude that $T = \{O, (0, 0)\} \cong \Z/2\Z$ for $y^2 = x^3 + x$. Why does $b$
        stay at $4$ however large $N$ is?
\begin{solution}
% Write your solution here.
\end{solution}
  \item\label{torsion:it:structure-lower} Show that a rational point $P$ is a torsion point if
        and only if $b_2P = O$, and that rational torsion points generating a subgroup of
        order $b_1b_2$ generate $T$.
\begin{solution}
% Write your solution here.
\end{solution}
\end{enumerate}
\end{problem}

\section{How good are the bounds? The LMFDB}
\label{torsion:sec:lmfdb}

The \href{https://www.lmfdb.org}{L-functions and modular forms database} \cite{lmfdb}
lists millions of elliptic curves over $\Q$ with their torsion subgroups, computed by
methods beyond \cite{Silverman09}*{Chapter~VII}. We use it as the answer key.

\begin{wsproposition}[{Reduction of isogenies \cite{Milne20}*{proof of Theorem~V.4.1}}]
\label{torsion:prop:isogeny-reduction}
An isogeny over $\Q$ reduces to an isogeny over $\F_p$ at every good prime $p$ of both
equations.
\end{wsproposition}

\begin{wsproposition}[{Isogenous curves over $\F_p$ \cite{Silverman09}*{Exercise~5.4(a)}}]
\label{torsion:prop:isogeny-count}
Isogenous curves over $\F_p$ have the same number of rational points.
\end{wsproposition}

\begin{problem}[Testing the bounds]
\label{torsion:prob:lmfdb}
\begin{enumerate}
  \item\label{torsion:it:lmfdb-magma} Here are the curves of smallest conductor with each
        torsion structure that occurs over $\Q$, with their LMFDB labels,
        $a$-invariants and torsion structures ($[1]$ for the trivial group):
\begin{magma}
curves := [<"11.a1", [0, -1, 1, -7820, -263580], [1]>,
  <"14.a2", [1, 0, 1, -171, -874], [2]>,
  <"19.a2", [0, 1, 1, -9, -15], [3]>, <"15.a4", [1, 1, 1, -80, 242], [4]>,
  <"11.a2", [0, -1, 1, -10, -20], [5]>, <"14.a6", [1, 0, 1, 4, -6], [6]>,
  <"26.b2", [1, -1, 1, -3, 3], [7]>, <"15.a8", [1, 1, 1, 35, -28], [8]>,
  <"54.b2", [1, -1, 1, -14, 29], [9]>, <"66.c3", [1, 0, 0, -45, 81], [10]>,
  <"90.c7", [1, -1, 1, -122, 1721], [12]>, <"15.a2", [1, 1, 1, -135, -660], [2, 2]>,
  <"15.a5", [1, 1, 1, -10, -10], [2, 4]>, <"30.a6", [1, 0, 1, -19, 26], [2, 6]>,
  <"210.e6", [1, 0, 0, -1070, 7812], [2, 8]>];
for c in curves do
  E := EllipticCurve(c[2]);
  printf "%-7o LMFDB %-9o order bound %-3o structure bound %o\n",
    c[1], c[3], OrderBound(E, 50), StructureBound(E, 50);
end for;
// 11.a1   LMFDB [ 1 ]     order bound 5   structure bound [ 1, 5 ]
// 14.a2   LMFDB [ 2 ]     order bound 6   structure bound [ 1, 6 ]
// 19.a2   LMFDB [ 3 ]     order bound 3   structure bound [ 1, 3 ]
// 15.a4   LMFDB [ 4 ]     order bound 8   structure bound [ 1, 4 ]
// 11.a2   LMFDB [ 5 ]     order bound 5   structure bound [ 1, 5 ]
// 14.a6   LMFDB [ 6 ]     order bound 6   structure bound [ 1, 6 ]
// 26.b2   LMFDB [ 7 ]     order bound 7   structure bound [ 1, 7 ]
// 15.a8   LMFDB [ 8 ]     order bound 8   structure bound [ 1, 8 ]
// 54.b2   LMFDB [ 9 ]     order bound 9   structure bound [ 1, 9 ]
// 66.c3   LMFDB [ 10 ]    order bound 10  structure bound [ 1, 10 ]
// 90.c7   LMFDB [ 12 ]    order bound 12  structure bound [ 1, 12 ]
// 15.a2   LMFDB [ 2, 2 ]  order bound 8   structure bound [ 2, 4 ]
// 15.a5   LMFDB [ 2, 4 ]  order bound 8   structure bound [ 2, 4 ]
// 30.a6   LMFDB [ 2, 6 ]  order bound 12  structure bound [ 2, 6 ]
// 210.e6  LMFDB [ 2, 8 ]  order bound 16  structure bound [ 2, 8 ]
\end{magma}
        Check that no bound contradicts the LMFDB. Where is each bound attained?
\begin{solution}
% Write your solution here.
\end{solution}
  \item\label{torsion:it:lmfdb-isogeny} The curves 15.a4 and 15.a8 are isogenous:
\begin{magma}
E := EllipticCurve([1, 1, 1, -80, 242]);        // 15.a4
F := EllipticCurve([1, 1, 1, 35, -28]);          // 15.a8
good := [p : p in PrimesUpTo(50) | p ne 3 and p ne 5];
[#ChangeRing(E, GF(p)) : p in good] eq [#ChangeRing(F, GF(p)) : p in good];
// true
[<p, Invariants(AbelianGroup(ChangeRing(E, GF(p)))),
     Invariants(AbelianGroup(ChangeRing(F, GF(p))))> : p in [7, 11, 13]];
// [ <7, [ 2, 4 ], [ 8 ]>, <11, [ 2, 8 ], [ 16 ]>, <13, [ 16 ], [ 2, 8 ]> ]
\end{magma}
        Using \cref{torsion:prop:isogeny-reduction,torsion:prop:isogeny-count}, explain why the order
        bound cannot separate isogenous curves, and why the structure bound separates these two.
\begin{solution}
% Write your solution here.
\end{solution}
  \item\label{torsion:it:lmfdb-certify} The LMFDB lists $(5, -2)$ as a torsion generator of
        15.a4:
\begin{magma}
E := EllipticCurve([1, 1, 1, -80, 242]);  P := E![5, -2];        // 15.a4
[k*P eq E!0 : k in [1..4]];                      // [ false, false, false, true ]
\end{magma}
        Prove that $T \cong \Z/4\Z$ for 15.a4 without trusting the LMFDB. Can you do
        the same for 11.a1, 14.a2 or 15.a2?
\begin{solution}
% Write your solution here.
\end{solution}
  \item\label{torsion:it:lmfdb-browse} Find these curves yourself: on
        \url{https://www.lmfdb.org/EllipticCurve/Q/}, search by ``torsion structure'',
        such as \mg{[2,4]}, and sort by conductor. On a curve's page, find its equation,
        discriminant, torsion subgroup and bad primes. Pick a curve of your own and
        compute both bounds.
\begin{solution}
% Write your solution here.
\end{solution}
\end{enumerate}
\end{problem}

\begin{wsremark}[Beyond the local theory]
\label{torsion:rem:beyond}
Over $\Q$ the bounds complete to an algorithm. For $y^2 = x^3 + Ax + B$ with
$A, B \in \Z$, a rational torsion point $P \neq O$ has integer coordinates, and either
$2P = O$ or $y(P)^2 \mid 4A^3 + 27B^2$ (Nagell--Lutz,
\cite{Silverman09}*{Corollary~VIII.7.2}). This leaves finitely many candidates, and the
test $b_2P = O$ decides which are torsion points. Mazur proved that $E(\Q)_{\tors}$ is
one of the fifteen groups in the listing \cite{Silverman09}*{Theorem~VIII.7.5}.
\end{wsremark}

\bigskip

\part[Good and bad reduction]{Good and bad reduction: examples}
\label{reduction:part}
\noindent
An elliptic curve over $\Qp$ has good, multiplicative or additive reduction, and
multiplicative reduction is split or nonsplit
\cite{Silverman09}*{\S VII.5}. The problems below run through concrete curves of each kind:
\cref{reduction:prob:read} reads the type off an equation, \cref{reduction:prob:lines} counts the rational
points of the singular reduction, and \cref{reduction:prob:change} changes the type by twisting the
curve or extending the field. \Cref{reduction:prob:frey} ($\star$) goes further: the curves Frey
attached to solutions of $A + B = C$ have no additive reduction at all.

\subsection*{Standing notation}
\begin{itemize}
  \item A prime $p$, with $v = v_p$ its valuation on $\Qp$, extended to $\overline{\Q}_p$.
  \item An elliptic curve $E$ over $\Qp$, given by a Weierstrass equation with
        coefficients in $\Zp$, with discriminant $\Delta$ and invariants $c_4$, $c_6$
        \cite{Silverman09}*{\S III.1}.
  \item When the equation is minimal, $E_p$ is its reduction modulo $p$, a Weierstrass
        cubic over $\F_p$.
  \item We write
        \[
        a_p \coloneqq p + 1 - \#E_p(\F_p),
        \]
        counting the singular point too.
\end{itemize}

\section{Reading off the type}
\label{reduction:sec:read}

The following definition and results hold for every prime $p$, and are used throughout.
In them, $K$ is $\Qp$ or a finite extension of it, with normalized valuation $v_K$, ring
of integers $\calO_K$ and residue field $k$, and $E$ is an elliptic curve over $K$ given by
a Weierstrass equation with coefficients in $\calO_K$.

\begin{wsdefinition}[{Reduction types \cite{Silverman09}*{\S VII.5}}]
\label{reduction:def:types}
Suppose that the equation is minimal, and consider its reduction, a Weierstrass cubic
over $k$. Then $E$ has \defi{good reduction} if the reduction is nonsingular,
\defi{multiplicative reduction} if the reduction has a node, and \defi{additive
reduction} if it has a cusp. Multiplicative reduction is \defi{split} if the slopes of the
tangent lines at the node lie in $k$, and \defi{nonsplit} otherwise.
\end{wsdefinition}

\begin{wsproposition}[{Discriminant and $c_4$ \cite{Silverman09}*{\S III.1}}]
\label{reduction:prop:invariants}
For an equation $y^2 = x^3 + a_2x^2 + a_4x + a_6$ (that is, $a_1 = a_3 = 0$),
\begin{equation}
\label{reduction:eq:invariants}
\Delta = -16\bigl(4a_2^3a_6 - a_2^2a_4^2 - 18a_2a_4a_6 + 4a_4^3 + 27a_6^2\bigr),
\qquad c_4 = 16\bigl(a_2^2 - 3a_4\bigr),
\end{equation}
which follow from the formulas of.
\end{wsproposition}

\begin{wsproposition}[{Minimal equations \cite{Silverman09}*{Remark~VII.1.1}}]
\label{reduction:prop:minimal}
Any other Weierstrass equation for $E$ has discriminant $u^{-12}\Delta$ and invariant
$u^{-4}c_4$ for some $u \in K^\times$. So if $v_K(\Delta) < 12$ or $v_K(c_4) < 4$, the
equation is minimal.
\end{wsproposition}

\begin{wsproposition}[{Singular Weierstrass cubics \cite{Silverman09}*{Proposition~III.1.4(a)}}]
\label{reduction:prop:singular}
A Weierstrass cubic over a field with $\Delta = 0$ has exactly one singular point over
an algebraic closure. It is a node, with two distinct tangent lines, if $c_4 \neq 0$,
and a cusp, with one tangent line, if $c_4 = 0$.
\end{wsproposition}

\begin{wsproposition}[{Criteria for the reduction type \cite{Silverman09}*{Proposition~VII.5.1}}]
\label{reduction:prop:types}
Suppose that the equation is minimal. Then $E$ has good reduction if $v_K(\Delta) = 0$,
multiplicative reduction if $v_K(\Delta) > 0$ and $v_K(c_4) = 0$, and additive reduction
if $v_K(\Delta) > 0$ and $v_K(c_4) > 0$.
\end{wsproposition}

\begin{wsproposition}[{Extending the field \cite{Silverman09}*{Proposition~VII.5.4(b)}}]
\label{reduction:prop:extensions}
Let $L$ be a finite extension of $K$ with ramification index $e$, so that $v_L = e\,v_K$
on $K$. If $E$ has good, respectively multiplicative, reduction, then so does $E_L$.
\end{wsproposition}

\begin{wsproposition}[{Potential good reduction \cite{Silverman09}*{Proposition~VII.5.5}}]
\label{reduction:prop:potential}
There is a finite extension $L$ of $K$ such that $E_L$ has good reduction if and only if
$j(E) \in \calO_K$.
\end{wsproposition}

\begin{problem}[Reading off the type]
\label{reduction:prob:read}
Let $p \geq 5$, and let $\varepsilon \in \Z$ be a non-square modulo $p$. For each
equation, decide whether it is minimal at $p$, and find the reduction type of the curve
over $\Qp$: good, split multiplicative, nonsplit multiplicative, or additive.
\begin{solution}
% Write your solution here.
\end{solution}
\begin{enumerate}
  \item\label{reduction:it:read-split} $y^2 = x^3 + x^2 + p$.
\begin{solution}
% Write your solution here.
\end{solution}
  \item\label{reduction:it:read-nonsplit} $y^2 = x^3 + \varepsilon x^2 + p$.
\begin{solution}
% Write your solution here.
\end{solution}
  \item\label{reduction:it:read-additive} $y^2 = x^3 + p$.
\begin{solution}
% Write your solution here.
\end{solution}
  \item\label{reduction:it:read-scaled} $y^2 = x^3 + p^6$.
\begin{solution}
% Write your solution here.
\end{solution}
  \item\label{reduction:it:read-twelve} $y^2 = x^3 + x^2 + p^{12}$.
\begin{solution}
% Write your solution here.
\end{solution}
  \item\label{reduction:it:read-magma} Here is \Magma's opinion at $p = 7$ and $\varepsilon = 3$.
        For each curve it prints the valuation of the minimal discriminant, the Kodaira
        symbol of the reduction ($\mathrm{I}_0$ for good reduction, $\mathrm{I}_n$ with
        $n \geq 1$ for multiplicative, anything else for additive), and the polynomial
        $1 - a_7T$, or $1 - a_7T + 7T^2$ in the case of good reduction.
\begin{magma}
R<T> := PolynomialRing(Integers());
p := 7; eps := 3;                          // 3 is not a square modulo 7
Info := function(a)
  E := EllipticCurve(a); L := LocalInformation(E, p);
  return <L[2], L[5], R!EulerFactor(E, p)>;
end function;
Info([0, 1, 0, 0, p]);                     // <1, I1, -T + 1>
Info([0, eps, 0, 0, p]);                   // <1, I1, T + 1>
Info([0, 0, 0, 0, p]);                     // <2, II, 1>
Info([0, 0, 0, 0, p^6]);                   // <0, I0, 7*T^2 + 4*T + 1>
Info([0, 1, 0, 0, p^12]);                  // <12, I12, -T + 1>
\end{magma}
        Check that it agrees with your answers. Note the values of $a_7$ at the curves
        with bad reduction. \Cref{reduction:prob:lines} explains them.
\begin{solution}
% Write your solution here.
\end{solution}
\end{enumerate}
\end{problem}

\section{Counting rational points on the singular reduction}
\label{reduction:sec:lines}

\begin{wsdefinition}[The nonsingular locus]
\label{reduction:def:ns}
When the equation of $E$ is minimal, the \defi{nonsingular locus} $E_p^{\ns}$ is $E_p$
minus its singular point, if it has one.
\end{wsdefinition}

\begin{problem}[Counting with lines through the singular point]
\label{reduction:prob:lines}
Let $E$ be an elliptic curve over $\Qp$ with bad reduction, and let $S$ be the singular
point of its reduction $E_p$.
\begin{enumerate}
  \item\label{reduction:it:lines-rational} Show that $S \in E_p(\F_p)$, and that
        $S \neq O$.\hint{what does the Frobenius automorphism of
        $\overline{\F}_p$ do to a singular point?}
\begin{solution}
% Write your solution here.
\end{solution}
  \item\label{reduction:it:lines-normal} Show that a change of variables over $\F_p$ moves $S$ to
        $(0, 0)$ and puts $E_p$ in the form
        \[
        y^2 + a_1xy = x^3 + a_2x^2,
        \]
        and that the tangent lines at $S$ are the lines $y = \alpha x$ with
        $g(\alpha) = 0$, where $g(t) \coloneqq t^2 + a_1t - a_2$.
\begin{solution}
% Write your solution here.
\end{solution}
  \item\label{reduction:it:lines-param} Let $t \in \F_p$. Show that the line $y = tx$ meets
        $E_p$ at $S$ and, if $g(t) \neq 0$, at exactly one more $\F_p$-rational point
        $P_t = \bigl(g(t), t\,g(t)\bigr)$, and nowhere else in the affine plane. Show that
        $t \mapsto P_t$ is a bijection from $\{t \in \F_p : g(t) \neq 0\}$ onto
        $E_p^{\ns}(\F_p) - \{O\}$.
\begin{solution}
% Write your solution here.
\end{solution}
  \item\label{reduction:it:lines-count} Deduce that $\#E_p^{\ns}(\F_p) = p + 1 - r$, where $r$ is the
        number of roots of $g$ in $\F_p$, and hence that
        \[
        a_p =
        \begin{cases}
        1 & \text{for split multiplicative reduction,} \\
        -1 & \text{for nonsplit multiplicative reduction,} \\
        0 & \text{for additive reduction.}
        \end{cases}
        \]
        Compare with the output in \cref{reduction:prob:read}\ref{reduction:it:read-magma}.
\begin{solution}
% Write your solution here.
\end{solution}
  \item\label{reduction:it:lines-small} Check \ref{reduction:it:lines-count} by hand for
        $y^2 = x^3 + x^2$ and $y^2 = x^3 + 2x^2$ over $\F_5$.
\begin{solution}
% Write your solution here.
\end{solution}
\end{enumerate}
\end{problem}

\section{Changing the curve, changing the field}
\label{reduction:sec:change}

\begin{problem}[Twists and extensions]
\label{reduction:prob:change}
Let $p \geq 5$. For $d \in \Qp^\times$ and $E\colon y^2 = x^3 + a_2x^2 + a_4x + a_6$, the
\defi{quadratic twist} of $E$ by $d$ is
\[
E^d\colon y^2 = x^3 + da_2x^2 + d^2a_4x + d^3a_6.
\]
\begin{enumerate}
  \item\label{reduction:it:change-twist} Show that $(x, y) \mapsto (dx, d^{3/2}y)$ is an isomorphism
        $E \to E^d$ over $\Qp(\sqrt d)$, and that $\Delta(E^d) = d^6\Delta(E)$,
        $c_4(E^d) = d^2c_4(E)$ and $j(E^d) = j(E)$.
\begin{solution}
% Write your solution here.
\end{solution}
\end{enumerate}
From now on, $E\colon y^2 = x^3 + x^2 + p$ is the curve of
\cref{reduction:prob:read}\ref{reduction:it:read-split}, with split multiplicative reduction.
\begin{enumerate}[resume]
  \item\label{reduction:it:change-unramified} Let $\varepsilon \in \Z$ be a non-square modulo $p$.
        Show that $E^\varepsilon$ has nonsplit multiplicative reduction, and that
        $a_p(E^\varepsilon) = -a_p(E)$. Since $\Qp(\sqrt\varepsilon)$ is an unramified
        extension of $\Qp$, what does this say about which features of the reduction are
        preserved by unramified extensions?
\begin{solution}
% Write your solution here.
\end{solution}
  \item\label{reduction:it:change-ramified} Show that $E^p\colon y^2 = x^3 + px^2 + p^4$ has additive
        reduction. Over which quadratic extension of $\Qp$ does it acquire multiplicative
        reduction? Show that it has good reduction over no finite extension of $\Qp$.
\begin{solution}
% Write your solution here.
\end{solution}
  \item\label{reduction:it:change-six} Let $E_3\colon y^2 = x^3 + p$, and let $L$ be a finite
        extension of $\Qp$ with ramification index $e$. Show that $E_3$ has good
        reduction over $L$ if and only if $6 \mid e$.\hint{for ``if'',
        write $p = \varpi^e w$ with $\varpi$ a uniformizer of $L$ and $w$ a unit.}
\begin{solution}
% Write your solution here.
\end{solution}
  \item\label{reduction:it:change-four} Show that $y^2 = x^3 + px$ has additive reduction, and good
        reduction over $L$ as in \ref{reduction:it:change-six} if and only if $4 \mid e$.
\begin{solution}
% Write your solution here.
\end{solution}
\end{enumerate}
\end{problem}

\section{\texorpdfstring{$\star$}{*} Frey curves}
\label{reduction:sec:frey}

This section is \cite{Silverman09}*{Exercise~8.23}, with the conductor of
\cite{Milne20}*{Aside~II.3.5}. The \href{https://www.lmfdb.org}{LMFDB} \cite{lmfdb} is
the answer key.

\begin{problem}[Frey curves are semistable]
\label{reduction:prob:frey}
An elliptic curve over $\Q$ is \defi{semistable} if it has good or multiplicative
reduction at every prime. Let $A, B, C$ be nonzero integers with $A + B = C$ and
$\gcd(A, B, C) = 1$. Let
\[
E\colon y^2 = x(x + A)(x - B).
\]
\begin{enumerate}
  \item\label{reduction:it:frey-invariants} Show that $\Delta = 16A^2B^2C^2$ and
        $c_4 = 16(A^2 + AB + B^2)$.
\begin{solution}
% Write your solution here.
\end{solution}
  \item\label{reduction:it:frey-odd} Let $p$ be an odd prime dividing $ABC$. Show that the equation
        is minimal at $p$ and that $E$ has multiplicative reduction at $p$. Show that the
        reduction is split if and only if $-B$, $A$ or $B$ is a square modulo $p$,
        according as $p$ divides $A$, $B$ or $C$.
\begin{solution}
% Write your solution here.
\end{solution}
  \item\label{reduction:it:frey-two} Suppose that $A \equiv 0 \pmod{16}$ and $B \equiv 3 \pmod 4$.
        Show that the substitution $x = 4x'$, $y = 8y' + 4x'$ gives an equation with
        coefficients in $\Z$,
        \[
        y'^2 + x'y' = x'^3 + \frac{A - B - 1}{4}\,x'^2 - \frac{AB}{16}\,x',
        \]
        with discriminant $(ABC)^2/2^8$ and $c_4 = A^2 + AB + B^2$. Deduce that $E$ has
        good reduction at $2$ if $v_2(A) = 4$, and multiplicative reduction at $2$ if
        $v_2(A) \geq 5$. So $E$ is semistable.
\begin{solution}
% Write your solution here.
\end{solution}
  \item\label{reduction:it:frey-conductor} The \defi{conductor} of $E$ is $N = \prod_p p^{f_p}$, where
        $f_p = 0$ if $E$ has good reduction at $p$, $f_p = 1$ if it has multiplicative
        reduction, and $f_p \geq 2$ if it has additive reduction. Under the hypotheses of
        part~\ref{reduction:it:frey-two}, compute $N$ in terms of $A, B, C$. Find the curves for
        $(A, B, C) = (16, -1, 15)$ and $(32, -5, 27)$ in the LMFDB, and check their
        conductors and their split and nonsplit primes against your answers.
\begin{solution}
% Write your solution here.
\end{solution}
  \item\label{reduction:it:frey-fermat} Suppose that $a^\ell + b^\ell = c^\ell$ with $\ell \geq 5$
        prime and $a, b, c$ nonzero coprime integers. Show that, after reordering and
        changing signs, $(A, B, C) = (a^\ell, b^\ell, c^\ell)$ satisfies the hypotheses
        of part~\ref{reduction:it:frey-two}, and that the resulting curve has conductor
        $\prod_{p \mid abc} p$ and minimal discriminant $(abc)^{2\ell}/2^8$.
\begin{solution}
% Write your solution here.
\end{solution}
\end{enumerate}
\end{problem}

\bigskip

\part[Compactness and inertia]{Compactness and inertia}
\label{compactness-inertia:part}
\noindent
Let $E$ be an elliptic curve over a local field $K$. The rational points of $E$ with
nonsingular reduction form a subgroup $E_0(K)$ of finite index in $E(K)$
\cite{Silverman09}*{Theorem~VII.6.1 and Corollary~VII.6.2}, and $E$ has good reduction if
and only if the inertia group acts trivially on its torsion, by the criterion of
N\'eron--Ogg--Shafarevich \cite{Silverman09}*{Theorem~VII.7.1}. \Cref{compactness-inertia:prob:compact}
proves the finiteness of $E(K)/E_0(K)$ by a compactness argument. \Cref{compactness-inertia:prob:inertia}
studies how the inertia group acts on the torsion of a curve with potential good
reduction.

\subsection*{Standing notation}
\begin{itemize}
  \item A finite extension $K$ of $\Qp$, with normalized valuation $v$, ring of integers
        $\calO_K$, uniformizer $\pi$ and residue field $k = \calO_K/\pi\calO_K$, a finite
        field of characteristic $p$. (Silverman's Chapter~VII assumes that $K$ and $k$
        are perfect, and here both are.)
  \item A fixed algebraic closure $\overline{K}$ of $K$.
  \item A bar denotes reduction modulo $\pi$.
  \item An elliptic curve $E$ over $K$, given by a minimal Weierstrass equation
        $F(X, Y, Z) = 0$, with $F \in \calO_K[X, Y, Z]$ homogeneous of degree $3$. It
        embeds $E$ in $\PP^2$ as $V_+(F)$.
  \item The \defi{primitive coordinates} $[X : Y : Z]$ of $P \in \PP^2(K)$: coordinates
        in $\calO_K$, not all in $\pi\calO_K$. Every $P$ has them, and they are unique up
        to a factor in $\calO_K^\times$.
  \item The \defi{reduction} $\bar P \coloneqq [\bar X : \bar Y : \bar Z] \in \PP^2(k)$ of
        $P \in \PP^2(K)$, computed from primitive coordinates, hence well defined.
  \item The reduced curve $E_\pi$, defined by $\bar F = 0$. If $P \in E(K)$, then
        $\bar P \in E_\pi(k)$.
\end{itemize}

\section{Finiteness by compactness}
\label{compactness-inertia:sec:compact}

This section is \cite{Silverman09}*{Exercise~7.6}. The argument uses that $k$ is finite,
so it does not apply, for instance, to the completion of the maximal unramified extension
of $\Qp$: a field complete with respect to a discrete valuation, with residue field
$\overline{\F}_p$ (item~\ref{compactness-inertia:it:compact-residue}). Yet this is the case that the criterion
of N\'eron--Ogg--Shafarevich needs: Silverman's proof of it uses the finiteness of
$E(K^{\unr})/E_0(K^{\unr})$, where $K^{\unr}$ is the maximal unramified extension of $K$,
whose residue field is an algebraic closure of $k$
\cite{Silverman09}*{proof of Theorem~VII.7.1}.

\begin{wsdefinition}[The topology of $K$]
\label{compactness-inertia:def:topology}
We give $K$ the topology defined by $v$: the sets $a + \pi^n\calO_K$, for $a \in K$ and
$n \in \Z$, form a basis, and each of them is both open and closed. With it, $K$ is a
Hausdorff topological field.
\end{wsdefinition}

\begin{wsdefinition}[The subgroup $E_0(K)$]
\label{compactness-inertia:def:E0}
Let $E_\pi^{\ns}$ be $E_\pi$ minus its singular point, if it has one. Then
\[
E_0(K) \coloneqq \{P \in E(K) : \bar P \in E_\pi^{\ns}(k)\}.
\]
\end{wsdefinition}

\begin{wsdefinition}[{Regular functions \cite{Poonen06}*{Definition~1.8.3}}]
\label{compactness-inertia:def:regular}
A function $\psi \in K(E)$ is \defi{regular} at $Q \in E(\overline{K})$ if $\psi = g/h$
with $g, h \in K[X, Y, Z]$ homogeneous of the same degree and $h(Q) \neq 0$. Then
$\psi(Q) \coloneqq g(Q)/h(Q)$, an element of $K$ if $Q \in E(K)$. The functions regular
at $Q$ form a discrete valuation ring with fraction field $K(E)$.
\end{wsdefinition}

\begin{wsproposition}[{$E_0(K)$ is a subgroup \cite{Silverman09}*{Proposition~VII.2.1}}]
\label{compactness-inertia:prop:subgroup}
The set $E_0(K)$ is a subgroup of $E(K)$.
\end{wsproposition}

\begin{wsproposition}[{Compactness \cite{Poonen2017}*{\S 1.1.2}}]
\label{compactness-inertia:prop:compact}
The ring $\calO_K$ is compact: it is complete, and for every $n$ it is covered by the
$(\#k)^n$ translates of $\pi^n\calO_K$. Equivalently, $K$ is locally compact.
\end{wsproposition}

\begin{wsproposition}[Morphisms $E \to E$ over $K$]
\label{compactness-inertia:prop:morphism}
\leavevmode
\begin{enumerate}
  \item\label{compactness-inertia:it:morphism-points} The morphisms $\alpha\colon E \to E$ defined over $K$
        are exactly the points $\alpha = [\psi_0 : \psi_1 : \psi_2] \in E(K(E))$, that
        is, $\psi_0, \psi_1, \psi_2 \in K(E)$, not all zero, with
        $F(\psi_0, \psi_1, \psi_2) = 0$ \cite{Poonen2017}*{Proposition~3.6.5(b)}.
  \item\label{compactness-inertia:it:morphism-eval} If $Q \in E(\overline{K})$ and $\psi_0, \psi_1, \psi_2$
        are regular at $Q$ and not all zero there, then
        $\alpha(Q) = [\psi_0(Q) : \psi_1(Q) : \psi_2(Q)]$ \cite{Silverman09}*{\S I.3}.
  \item\label{compactness-inertia:it:morphism-rescale} Replacing each $\psi_i$ by $\phi\psi_i$, for
        $\phi \in K(E)^\times$, gives the same point of $E(K(E))$, and for every
        $Q \in E(\overline{K})$ some $\phi$ makes $\phi\psi_0, \phi\psi_1, \phi\psi_2$
        regular at $Q$ and not all zero there: a power of a uniformizer of the discrete
        valuation ring of \cref{compactness-inertia:def:regular} will do
        \cite{Silverman09}*{proof of Proposition~II.2.1}.
\end{enumerate}
\end{wsproposition}

\begin{wsproposition}[{Translations \cite{Silverman09}*{Theorem~III.3.6 and Example~III.4.7}}]
\label{compactness-inertia:prop:translation}
Let $P \in E(K)$. The translation $\tau_P\colon E \to E$, $Q \mapsto Q + P$, is an
isomorphism defined over $K$, with inverse $\tau_{-P}$.
\end{wsproposition}

\begin{problem}[$E(K)/E_0(K)$ is finite]
\label{compactness-inertia:prob:compact}
Let $q\colon K^3 - \{0\} \to \PP^2(K)$ be the quotient map. Give $K^3 - \{0\}$ the
product topology and $\PP^2(K)$ the quotient topology: $U \subseteq \PP^2(K)$ is open if
and only if $q^{-1}(U)$ is open. Let
\[
S \coloneqq \{\bfx = (x_0, x_1, x_2) \in \calO_K^3 : x_i \in \calO_K^\times
\text{ for some } i\}
\]
be the set of primitive coordinate vectors.
\begin{enumerate}
  \item\label{compactness-inertia:it:compact-proj} Show that $S$ is compact and open in $K^3 - \{0\}$, and
        that $q(S) = \PP^2(K)$. Deduce that $\PP^2(K)$ is compact. Show that
        $U \subseteq \PP^2(K)$ is open if and only if $q^{-1}(U) \cap S$ is open in
        $S$.\hint{$q^{-1}(U)$ is stable under multiplication by
        $K^\times$.}
\begin{solution}
% Write your solution here.
\end{solution}
  \item\label{compactness-inertia:it:compact-closed} Show that $E(K) = V_+(F)(K)$ is closed in $\PP^2(K)$,
        and deduce that it is compact.
\begin{solution}
% Write your solution here.
\end{solution}
  \item\label{compactness-inertia:it:compact-maps} Let $g, h \in K[X, Y, Z]$ be homogeneous of the same
        degree. Show that $W \coloneqq \{Q \in \PP^2(K) : h(Q) \neq 0\}$ is open, and that
        $W \to K$, $Q \mapsto g(Q)/h(Q)$, is continuous. Deduce from
        \cref{compactness-inertia:prop:morphism,compactness-inertia:prop:translation} that every translation $\tau_P\colon E(K) \to E(K)$ is a
        homeomorphism.
\begin{solution}
% Write your solution here.
\end{solution}
  \item\label{compactness-inertia:it:compact-open} Show that the reduction map $\PP^2(K) \to \PP^2(k)$,
        $P \mapsto \bar P$, is locally constant. Deduce that $E_0(K)$ is open in $E(K)$,
        and also closed.
\begin{solution}
% Write your solution here.
\end{solution}
  \item\label{compactness-inertia:it:compact-finite} Prove that $E(K)/E_0(K)$ is finite.
\begin{solution}
% Write your solution here.
\end{solution}
  \item\label{compactness-inertia:it:compact-residue} Let $K'$ be a field complete with respect to a discrete
        valuation whose residue field $k'$ is infinite, for instance algebraically closed.
        Show that its ring of integers is not compact. Which step of the argument above
        fails for an elliptic curve over $K'$?
\begin{solution}
% Write your solution here.
\end{solution}
\end{enumerate}
\end{problem}

\section{Inertia acting on torsion}
\label{compactness-inertia:sec:inertia}

This section is \cite{Silverman09}*{Exercise~7.9}, with the hypothesis $m \geq 3$ added to
its first two parts (item~\ref{compactness-inertia:it:inertia-two} explains why) and a \Magma\ check. We use
the notation of \cite{Silverman09}*{\S VII.4}.

\begin{wsdefinition}[Reduction types]
\label{compactness-inertia:def:reduction}
The curve $E$ has \defi{good reduction} if $E_\pi$ is nonsingular, that is, if
$v(\Delta) = 0$, and \defi{bad reduction} otherwise. It has \defi{potential good
reduction} if it has good reduction over some finite extension of $K$.
\end{wsdefinition}

\begin{wsdefinition}[Inertia]
\label{compactness-inertia:def:inertia}
Let $\Gal_K \coloneqq \Gal(\overline{K}/K)$, $K^{\unr} \subseteq \overline{K}$ the maximal
unramified extension of $K$, and $I_v \coloneqq \Gal(\overline{K}/K^{\unr})$ the
\defi{inertia group}. A set on which $\Gal_K$ acts is \defi{unramified} if $I_v$ acts
trivially on it.
\end{wsdefinition}

\begin{wsdefinition}[Ramification of finite extensions]
\label{compactness-inertia:def:ramification}
For an extension $L$ of $K$ inside $\overline{K}$, let $\Gal_L \coloneqq \Gal(\overline{K}/L)$.
If $L/K$ is finite and Galois, the \defi{inertia group} of $L/K$ is the image of $I_v$ in
$\Gal(L/K)$, isomorphic to $I_v/(I_v \cap \Gal_L)$. The extension $L/K$ is
\defi{unramified} if its inertia group is trivial, and \defi{tamely ramified} if its order
is prime to $p$.
\end{wsdefinition}

\begin{wsdefinition}[{Torsion fields and Tate modules \cite{Silverman09}*{\S III.7}}]
\label{compactness-inertia:def:tate}
For $m \geq 1$ with $p \nmid m$, $K(E[m])$ is the field generated over $K$ by the
coordinates of the points of $E[m](\overline{K})$, so $\Gal_{K(E[m])}$ is the kernel of the
action of $\Gal_K$ on $E[m](\overline{K})$. For a prime $\ell \neq p$, the \defi{Tate module}
$T_\ell(E)$ is the inverse limit of the $E[\ell^n](\overline{K})$ along the surjections $[\ell]$,
and $\rho_\ell\colon \Gal_K \to \Aut T_\ell(E)$ is the action
of $\Gal_K$ on it.
\end{wsdefinition}

In the following results, $\ell \neq p$ is a prime and $m \geq 1$ is an integer with
$p \nmid m$.

\begin{wsproposition}[{Minimal equations \cite{Silverman09}*{Remark~VII.1.1 and \S III.1}}]
\label{compactness-inertia:prop:minimal}
A Weierstrass equation with coefficients in $\calO_K$ and $v(\Delta) < 12$ is minimal. For $y^2 = x^3 + Ax + B$,
$\Delta = -16(4A^3 + 27B^2)$.
\end{wsproposition}

\begin{wscorollary}[{Torsion \cite{Silverman09}*{Corollary~III.6.4(b)}}]
\label{compactness-inertia:cor:torsion}
The group $E[m](\overline{K})$ is isomorphic to $(\Z/m\Z)^2$.
\end{wscorollary}

\begin{wsproposition}[{The Tate module \cite{Silverman09}*{Proposition~III.7.1(a)}}]
\label{compactness-inertia:prop:tate}
The Tate module $T_\ell(E)$ is isomorphic to $\Z_\ell^2$. Hence each
projection $T_\ell(E) \to E[\ell^n](\overline{K})$ is a surjective homomorphism of
$\Gal_K$-modules, and its kernel is $\ell^nT_\ell(E)$, since both have index $\ell^{2n}$.
\end{wsproposition}

\begin{wsproposition}[{Good reduction \cite{Silverman09}*{Proposition~VII.4.1}}]
\label{compactness-inertia:prop:good}
If $E$ has good reduction, then $E[m](\overline{K})$ and $T_\ell(E)$ are unramified.
\end{wsproposition}

\begin{wstheorem}[{N\'eron--Ogg--Shafarevich \cite{Silverman09}*{Theorem~VII.7.1}}]
\label{compactness-inertia:thm:nos}
If $T_\ell(E)$ is unramified, then $E$ has good reduction.
\end{wstheorem}

\begin{wscorollary}[{Potential good reduction \cite{Silverman09}*{Corollary~VII.7.3}}]
\label{compactness-inertia:cor:potential}
The curve $E$ has potential good reduction if and only if $\rho_\ell(I_v)$ is finite.
\end{wscorollary}

\begin{wsproposition}[{Inertia groups of finite extensions \cite{Silverman09}*{proof of Corollary~VII.7.3}}]
\label{compactness-inertia:prop:finite-index}
If $H \subseteq I_v$ is a closed normal subgroup of finite index, there is a finite
extension $K'$ of $K$ inside $\overline{K}$ whose inertia group
$\Gal(\overline{K}/K'K^{\unr}) = I_v \cap \Gal_{K'}$ is $H$.
\end{wsproposition}

\begin{problem}[The inertia group of $K(E{[m]})/K$]
\label{compactness-inertia:prob:inertia}
Suppose that $E$ has potential good reduction, and let $J$ be the subgroup of the elements
of $I_v$ that act trivially on $E[m](\overline{K})$ for every $m \geq 1$ with $p \nmid m$.
\begin{enumerate}
  \item\label{compactness-inertia:it:inertia-torsionfree} Let $\ell$ be a prime and $r \geq 1$ with
        $\ell^r \geq 3$, and let $g \in \GL_2(\Z_\ell)$ be an element of finite order with
        $g \equiv I \pmod{\ell^r}$. Show that $g = I$.\hint{reduce to $g$ of prime order
        $q$, write $g = I + \ell^s h$ with $h \in M_2(\Z_\ell)$ not divisible by $\ell$,
        and expand $g^q$.}
\begin{solution}
% Write your solution here.
\end{solution}
  \item\label{compactness-inertia:it:inertia-ell} Show that $J = I_v \cap \ker\rho_\ell$ for every prime
        $\ell \neq p$, and that $J$ has finite index in $I_v$.
\begin{solution}
% Write your solution here.
\end{solution}
  \item\label{compactness-inertia:it:inertia-m} Let $m \geq 3$ with $p \nmid m$. Show that
        $J = I_v \cap \Gal_{K(E[m])}$, and deduce that the inertia group of $K(E[m])/K$ is
        isomorphic to $I_v/J$, independently of $m$.
\begin{solution}
% Write your solution here.
\end{solution}
  \item\label{compactness-inertia:it:inertia-good} Let $m \geq 3$ with $p \nmid m$. Show that $K(E[m])/K$ is
        unramified if and only if $E$ has good reduction. Which direction holds for every
        $m \geq 1$ with $p \nmid m$, and for every elliptic curve over $K$?
\begin{solution}
% Write your solution here.
\end{solution}
  \item\label{compactness-inertia:it:inertia-two} Let $p \geq 3$ and $E\colon y^2 = x^3 - \pi^2x$.
        \begin{enumerate}
          \item\label{compactness-inertia:it:inertia-two-reduction} Show that $E$ has bad reduction, potential
                good reduction, and $K(E[2]) = K$.
\begin{solution}
% Write your solution here.
\end{solution}
          \item\label{compactness-inertia:it:inertia-two-twist} Show that $I_v/J$ has order $2$, and that its
                nontrivial element acts on every $T_\ell(E)$ as
                $-1$.\hint{over $K(\sqrt\pi)$, $E$ is isomorphic to
                $y^2 = x^3 - x$.}
\begin{solution}
% Write your solution here.
\end{solution}
          \item\label{compactness-inertia:it:inertia-two-where} So \ref{compactness-inertia:it:inertia-m} and \ref{compactness-inertia:it:inertia-good}
                fail for $m = 2$. Where did your proofs use $m \geq 3$?
\begin{solution}
% Write your solution here.
\end{solution}
        \end{enumerate}
  \item\label{compactness-inertia:it:inertia-tame} Suppose that $p \geq 5$. Show that $K(E[m])/K$ is tamely
        ramified for every $m \geq 1$ with $p \nmid m$.\hint{$\#\GL_2(\F_3) = 48$.}
\begin{solution}
% Write your solution here.
\end{solution}
  \item\label{compactness-inertia:it:inertia-magma} \Magma\ computes $\#(I_v/J)$ for elliptic curves over
        $\Q_7$:
\begin{magma}
p := 7;
Inertia := func< A, B |
  #InertiaGroup(GaloisRepresentation(EllipticCurve([A, B]), p)) >;
Inertia(-1, 0);                     // 1
Inertia(-p^2, 0);                   // 2
Inertia(0, p);                      // 6
\end{magma}
        Check the first two values against your answers, and the third against
        \ref{compactness-inertia:it:inertia-tame}.
\begin{solution}
% Write your solution here.
\end{solution}
\end{enumerate}
\end{problem}

\bigskip

\part[Hilbert 90 and Kummer theory]{Hilbert 90 and Kummer theory}
\label{kummer:part}
\noindent
Hilbert's Theorem 90 says that $H^1(k, \Gm) = 1$. Applied to the Kummer sequence
$1 \to \mu_m \to \Gm \xrightarrow{m} \Gm \to 1$, it gives an isomorphism
$k^\times/k^{\times m} \cong H^1(k, \mu_m)$ with an explicit formula, and, when $k$
contains the $m$-th roots of unity, a classification of the abelian extensions of $k$ of
exponent dividing $m$. The problems below prove all of this from the definitions in
\cite{Silverman09}*{Appendix~B}. \Cref{kummer:prob:cocycles} computes cohomology groups by
hand, \cref{kummer:prob:h90} proves Hilbert 90, \cref{kummer:prob:kseq} unwinds the Kummer sequence and
\cref{kummer:prob:kth} derives Kummer theory. \Cref{kummer:prob:local,kummer:prob:global} compute the groups
$k^\times/k^{\times m}$ for finite, local and global fields, and cut out a finite subgroup
by conditions at the primes. Later, the same sequence for an elliptic curve $E$ will give
an injection $E(k)/m \hookrightarrow H^1(k, E[m])$. Then there is no Hilbert 90, and its
failure is measured by $H^1(k, E)$.

\subsection*{Standing notation}
\begin{itemize}
  \item A field $k$, not assumed perfect, a separable closure $\ksep$ of $k$, and the
        absolute Galois group $\Gal_k \coloneqq \Gal(\ksep/k)$. For an extension $L$ of $k$
        inside $\ksep$, $\Gal_L \coloneqq \Gal(\ksep/L)$. (Silverman's Appendix~B assumes $k$
        perfect and works with an algebraic closure. Nothing in it changes with $\ksep$
        in its place.)
  \item A profinite group $G$, often finite, and a discrete $G$-module $A$: an abelian
        group with a left action of $G$, written $a \mapsto ga$, such that $G \times A \to A$
        is continuous for the discrete topology on $A$. (Silverman writes right actions,
        $a \mapsto a^\sigma$.)
  \item Elements $g, h$ of $G$. The letter $\sigma$ is kept for a generator of a cyclic
        group.
  \item An integer $m \geq 1$ with $\Char k \nmid m$, and the algebraic group $\mu_m$ of
        $m$-th roots of unity, the kernel of $x \mapsto x^m$ on $\Gm$. For a field $L$
        containing $k$, $\mu_m(L)$ is the group of $m$-th roots of unity in $L$.
  \item \textbf{An abuse of notation.} For a commutative algebraic group $A$ over $k$,
        such as $\Gm$, $\mu_m$ or an elliptic curve, we write
        $H^i(k, A) \coloneqq H^i(\Gal_k, A(\ksep))$, the cohomology of the
        $\Gal_k$-module of $\ksep$-points of $A$. For instance,
        \[
        H^1(k, \Gm) = H^1(\Gal_k, (\ksep)^\times), \qquad
        H^1(k, \mu_m) = H^1(\Gal_k, \mu_m(\ksep)).
        \]
        This is the only place where $A$ stands for $A(\ksep)$.
  \item The quotient $B/m \coloneqq B/mB$ of an abelian group $B$. For the
        multiplicative group of a field $L$ it is $L^\times/L^{\times m}$, where
        $L^{\times m}$ is the subgroup of $m$-th powers.
\end{itemize}

\section{Cocycles by hand}
\label{kummer:sec:cocycles}

This section uses the definitions of \cite{Silverman09}*{\S\S B.1--B.2} and
\cite{Poonen02}*{\S\S 2--3}. Its last item is \cite{Silverman09}*{Exercise~B.2(a)}.

\begin{wsdefinition}[{Cohomology in degrees $0$ and $1$ \cite{Silverman09}*{\S\S B.1--B.2}, \cite{Poonen02}*{\S 3}}]
\label{kummer:def:cohomology}
Let $A$ be a discrete $G$-module. Then $H^0(G, A) \coloneqq A^G$, the subgroup of the
$a \in A$ with $ga = a$ for every $g \in G$. A \defi{left $1$-cocycle}, or simply $1$-cocycle, is a continuous map
$\xi\colon G \to A$, $g \mapsto \xi_g$, such that
\[
\xi_{gh} = \xi_g + g\xi_h \qquad \text{for all } g, h \in G.
\]
A \defi{left $1$-coboundary}, or simply $1$-coboundary, is a map of the form
$g \mapsto ga - a$, for some $a \in A$. Both are \emph{left}: they are written for the
left action $a \mapsto ga$. Silverman's are the right ones, for $a \mapsto a^\sigma$:
$\xi_{\sigma\tau} = \xi_\sigma^\tau + \xi_\tau$ and $\xi_\sigma = a^\sigma - a$. The
$1$-cocycles form a group $Z^1(G, A)$ under addition of values, the $1$-coboundaries a
subgroup $B^1(G, A)$, and $H^1(G, A) \coloneqq Z^1(G, A)/B^1(G, A)$. If $G$ is finite,
every map $G \to A$ is continuous. When the group law of $A$ is written multiplicatively,
as for $L^\times$ or $\mu_m(L)$, the conditions read $\xi_{gh} = \xi_g \cdot g(\xi_h)$ and
$\xi_g = g(a)/a$.
\end{wsdefinition}

\begin{wsproposition}[{Trivial action \cite{Silverman09}*{Remarks~B.1.1 and B.2.2}, \cite{Poonen02}*{\S 3}}]
\label{kummer:prop:trivial}
If $G$ acts trivially on $A$, then $H^0(G, A) = A$ and
$H^1(G, A) = \Hom_{\mathrm{conts}}(G, A)$, the group of continuous homomorphisms.
\end{wsproposition}

\begin{problem}[Computing $H^1$ from the definition]
\label{kummer:prob:cocycles}
In this problem $G$ is finite.
\begin{enumerate}
  \item\label{kummer:it:coc-sign} Let $G = \{1, \sigma\}$ have order $2$, acting on $\Z$ by
        $\sigma a = -a$. Compute $H^0(G, \Z)$ and $H^1(G, \Z)$. Compare with the trivial
        action of $G$ on $\Z$.
\begin{solution}
% Write your solution here.
\end{solution}
  \item\label{kummer:it:coc-cyclic} Let $G = \langle \sigma \rangle$ be cyclic of order $n$, and
        $N \coloneqq 1 + \sigma + \dots + \sigma^{n - 1}$, acting on $A$. Show that
        $\xi \mapsto \xi_\sigma$ induces an
        isomorphism\hint{show that $\xi_{\sigma^i} = (1 + \sigma + \dots +
        \sigma^{i - 1})\xi_\sigma$ for $i \geq 1$.}
        \[
        H^1(G, A) \cong \ker(N\colon A \to A)/(\sigma - 1)A.
        \]
\begin{solution}
% Write your solution here.
\end{solution}
  \item\label{kummer:it:coc-real} Let $G = \Gal(\C/\R) = \{1, c\}$. Use \ref{kummer:it:coc-cyclic} to
        compute $H^1(G, \C^\times)$ and $H^1(G, \mu_m(\C))$.
        Compare the order of the second group with that of $\R^\times/\R^{\times m}$.
\begin{solution}
% Write your solution here.
\end{solution}
  \item\label{kummer:it:coc-order} Suppose that $G$ has order $n$. Show that $nH^1(G, A) = 0$.
        \hint{sum the cocycle condition $\xi_{gh} = \xi_g + g\xi_h$ over $h \in G$.}
\begin{solution}
% Write your solution here.
\end{solution}
\end{enumerate}
\end{problem}


\section{Hilbert 90}
\label{kummer:sec:h90}

This section is \cite{Poonen2017}*{Proposition~1.3.15(ii) and Exercises~1.10--1.11},
with the proof of \cite{Milne20}*{Proposition~IV.1.3} (also \cite{Milne22}*{Theorem~5.23}).
\cite{Silverman09}*{Proposition~B.2.5} states the result and refers to Serre for the
proof. Item~\ref{kummer:it:h90-cyclic} is the form in which Hilbert stated it.

\begin{wstheorem}[{Dedekind \cite{Milne22}*{Theorem~5.14}}]
\label{kummer:thm:dedekind}
Let $F$ be a field and $H$ a group. If $\chi_1, \dots, \chi_r$ are distinct
homomorphisms $H \to F^\times$, then they are linearly independent over $F$: if
$a_1\chi_1 + \dots + a_r\chi_r = 0$ as a function $H \to F$, with $a_i \in F$, then
$a_1 = \dots = a_r = 0$.
\end{wstheorem}

\begin{wsdefinition}[Norm and trace]
\label{kummer:def:norm}
Let $L/k$ be a finite Galois extension with group $G$. The \defi{norm} and the
\defi{trace} of $a \in L$ are $N_{L/k}(a) \coloneqq \prod_{g \in G} g(a)$ and
$\Tr_{L/k}(a) \coloneqq \sum_{g \in G} g(a)$. Both lie in $k$.
\end{wsdefinition}

\begin{wsproposition}[Infinite Galois theory]
\label{kummer:prop:galois}
\leavevmode
\begin{enumerate}
  \item\label{kummer:it:galois-basis} The subgroups $\Gal_L$, for $L$ running over the finite
        Galois extensions of $k$ inside $\ksep$, are open and normal in $\Gal_k$, and they
        form a basis of open neighborhoods of $1$ \cite{Silverman09}*{\S B.2}.
  \item\label{kummer:it:galois-corr} The maps $H \mapsto (\ksep)^H$ and $L \mapsto \Gal_L$ are
        inverse bijections between the closed subgroups of $\Gal_k$ and the extensions of $k$
        inside $\ksep$. The subgroup $H$ is open if and only if $(\ksep)^H$ is finite over $k$,
        and then $[\Gal_k : H] = [(\ksep)^H : k]$. It is normal if and only if $(\ksep)^H$ is
        Galois over $k$, and then $\Gal((\ksep)^H/k) \cong \Gal_k/H$ by restriction. In
        particular, $(\ksep)^{\Gal_k} = k$ \cite{Milne22}*{Theorem~7.13}.
\end{enumerate}
\end{wsproposition}

\begin{wsdefinition}[{Restriction and inflation \cite{Silverman09}*{\S B.2}, \cite{Poonen02}*{\S 4}}]
\label{kummer:def:res-inf}
Let $A$ be a discrete $G$-module and $H \subseteq G$ a closed subgroup. The
\defi{restriction} $\Res\colon H^1(G, A) \to H^1(H, A)$ maps the class of $\xi$ to the
class of $\xi|_H$. If $H$ is normal, the \defi{inflation}
$\Inf\colon H^1(G/H, A^H) \to H^1(G, A)$ maps the class of $\xi$ to the class of the
composition $G \to G/H \xrightarrow{\xi} A^H \subseteq A$. A $1$-cocycle of $G$ of this form
is \defi{inflated} from $G/H$.
\end{wsdefinition}

\begin{problem}[Hilbert's Theorem 90]
\label{kummer:prob:h90}
\leavevmode
\begin{enumerate}
  \item\label{kummer:it:h90-finite} Let $L/k$ be a finite Galois extension with group $G$, and
        $\xi\colon G \to L^\times$ a $1$-cocycle. Show that there is a $c \in L^\times$
        with
        \[
        b \coloneqq \sum_{h \in G} \xi_h \cdot h(c) \neq 0,
        \]
        and that then $\xi_g = b/g(b)$ for every $g \in G$. Deduce that $H^1(G, L^\times) = 1$.
        \hint{compute $g(b)$ with the cocycle condition.}
\begin{solution}
% Write your solution here.
\end{solution}
  \item\label{kummer:it:h90-cyclic} Suppose moreover that $G = \langle \sigma \rangle$ is cyclic.
        Show that if $a \in L^\times$ has $N_{L/k}(a) = 1$, then $a = \sigma(b)/b$ for some
        $b \in L^\times$.
\begin{solution}
% Write your solution here.
\end{solution}
  \item\label{kummer:it:h90-circle} Let $(x, y)$ be a rational point of the circle
        $x^2 + y^2 = 1$. Apply \ref{kummer:it:h90-cyclic} to $\Q(i)/\Q$ to find $u, v \in \Q$,
        not both $0$, with
        \[
        (x, y) = \Bigl(\frac{u^2 - v^2}{u^2 + v^2}, \frac{2uv}{u^2 + v^2}\Bigr).
        \]
\begin{solution}
% Write your solution here.
\end{solution}
  \item\label{kummer:it:h90-limit} Let $\xi\colon \Gal_k \to (\ksep)^\times$ be a continuous
        $1$-cocycle.
        \begin{enumerate}
          \item\label{kummer:it:h90-limit-factor} Show that there is a finite Galois extension
                $L$ of $k$ inside $\ksep$ such that $\xi$ takes values in $L^\times$ and
                $\xi_{gh} = \xi_g$ for all $g \in \Gal_k$ and $h \in \Gal_L$. So $\xi$ is
                inflated from a $1$-cocycle $\Gal(L/k) \to L^\times$.\hint{$\Gal_k$ is
                compact, and $\xi^{-1}(1)$ is open.}
\begin{solution}
% Write your solution here.
\end{solution}
          \item\label{kummer:it:h90-limit-zero} Deduce that $H^1(k, \Gm) = 1$.
\begin{solution}
% Write your solution here.
\end{solution}
        \end{enumerate}
  \item\label{kummer:it:h90-additive} ($\star$) Let $L/k$ be finite Galois with group $G$, and
        $\xi\colon G \to L$ a $1$-cocycle for the additive group of $L$. Show that
        $H^1(G, L) = 0$, by a variant of \ref{kummer:it:h90-finite}.
        \hint{take $c \in L$ with $\Tr_{L/k}(c) \neq 0$.}
\begin{solution}
% Write your solution here.
\end{solution}
\end{enumerate}
\end{problem}


\section{The Kummer sequence}
\label{kummer:sec:kseq}

This section is the exercise in \cite{Poonen02}*{\S 3}, with the explicit formula of
\cite{Silverman09}*{Propositions~B.2.5(c) and VIII.2.2}. Hilbert's Theorem 90,
$H^1(k, \Gm) = 1$, is \cref{kummer:prob:h90}\ref{kummer:it:h90-limit}.

\begin{wsproposition}[{Finite multiplicative groups \cite{Milne22}*{Exercise~1-3}}]
\label{kummer:prop:cyclic}
Every finite subgroup of the multiplicative group of a field is cyclic.
\end{wsproposition}

\begin{wsproposition}[{The long exact sequence \cite{Silverman09}*{Propositions~B.1.2 and B.2.3}, \cite{Poonen02}*{Theorem~1}}]
\label{kummer:prop:les}
Let $0 \to A \to B \xrightarrow{\psi} C \to 0$ be an exact sequence of discrete $G$-modules.
There is an exact sequence
\[
0 \to A^G \to B^G \to C^G \xrightarrow{\delta} H^1(G, A) \to H^1(G, B) \to H^1(G, C),
\]
where, for $c \in C^G$, $\delta(c)$ is the class of $g \mapsto gb - b$ for any
$b \in B$ with $\psi(b) = c$.
\end{wsproposition}

\begin{wsproposition}[{Inflation--restriction \cite{Silverman09}*{Proposition~B.2.4}}]
\label{kummer:prop:inf-res}
Let $A$ be a discrete $\Gal_k$-module and $L/k$ a finite Galois extension inside $\ksep$. The
sequence
\[
0 \to H^1(\Gal(L/k), A^{\Gal_L}) \xrightarrow{\Inf} H^1(\Gal_k, A) \xrightarrow{\Res} H^1(\Gal_L, A)
\]
is exact.
\end{wsproposition}

\begin{problem}[The Kummer sequence]
\label{kummer:prob:kseq}
Recall that $\Char k \nmid m$, so $m \neq 0$ in $k$.
\begin{enumerate}
  \item\label{kummer:it:kseq-exact} Show that $\mu_m(\ksep)$ is cyclic of order $m$, and that
        \[
        1 \to \mu_m(\ksep) \to (\ksep)^\times \xrightarrow{x \mapsto x^m} (\ksep)^\times \to 1
        \]
        is an exact sequence of discrete $\Gal_k$-modules. Deduce from \cref{kummer:prop:les} that
        its connecting map induces an isomorphism
        $\delta\colon k^\times/k^{\times m} \to H^1(k, \mu_m)$.
\begin{solution}
% Write your solution here.
\end{solution}
  \item\label{kummer:it:kseq-delta} Let $a \in k^\times$ and $\alpha \in \ksep$ with
        $\alpha^m = a$. Check directly, without \cref{kummer:prop:les}, that
        $\xi_g \coloneqq g(\alpha)/\alpha$ is a continuous $1$-cocycle with values in
        $\mu_m(\ksep)$, whose class does not depend on the choice of $\alpha$. Show that
        $a \mapsto [\xi]$ is a homomorphism $k^\times \to H^1(k, \mu_m)$ with kernel
        $k^{\times m}$. This is the map $\delta$ of \ref{kummer:it:kseq-exact}.
\begin{solution}
% Write your solution here.
\end{solution}
  \item\label{kummer:it:kseq-onto} Show directly from Hilbert 90 that every continuous
        $1$-cocycle $\xi\colon \Gal_k \to \mu_m(\ksep)$ is cohomologous to one of the form of
        \ref{kummer:it:kseq-delta}.
\begin{solution}
% Write your solution here.
\end{solution}
  \item\label{kummer:it:kseq-quadratic} Let $m = 2$, with $\Char k \neq 2$. Show that
        $H^1(k, \mu_2) = \Hom_{\mathrm{conts}}(\Gal_k, \{\pm 1\})$, and that for
        $d \in k^\times - k^{\times 2}$, $\delta(d)$ is the homomorphism
        $\chi_d\colon g \mapsto g(\sqrt d)/\sqrt d$, with kernel $\Gal_{k(\sqrt d)}$. Deduce
        that $d \mapsto k(\sqrt d)$ is a bijection from the nontrivial elements of
        $k^\times/k^{\times 2}$ to the quadratic extensions of $k$ inside $\ksep$.
\begin{solution}
% Write your solution here.
\end{solution}
  \item\label{kummer:it:kseq-res} Let $L$ be an extension of $k$ inside $\ksep$. Show that the
        restriction $H^1(k, \mu_m) \to H^1(L, \mu_m)$ corresponds, under $\delta$, to the
        map $k^\times/k^{\times m} \to L^\times/L^{\times m}$ induced by inclusion, and
        that its kernel corresponds to $(k^\times \cap L^{\times m})/k^{\times m}$. Compute
        the kernel for $\Q(i)/\Q$ and $m = 2$, and compare it with \cref{kummer:prop:inf-res}.
\begin{solution}
% Write your solution here.
\end{solution}
\end{enumerate}
\end{problem}


\section{Kummer theory}
\label{kummer:sec:kummer}

This section is \cite{Silverman09}*{Exercise~8.4}, in the form of
\cite{Milne22}*{Theorem~5.30 and Remark~5.32}, together with an example that shows why
the roots of unity are needed. In it, $\mu_m(\ksep) \subseteq k$, except in
item~\ref{kummer:it:kth-cube}.

\begin{wsproposition}[{Characters of a finite abelian group \cite{Milne22}*{discussion before Theorem~5.30}}]
\label{kummer:prop:characters}
If $B$ is a finite abelian group of exponent dividing $m$, then $\Hom(B, \mu_m(\ksep))$ has the
same order as $B$.
\end{wsproposition}

\begin{problem}[Abelian extensions of exponent dividing $m$]
\label{kummer:prob:kth}
\leavevmode
\begin{enumerate}
  \item\label{kummer:it:kth-cyclic} Suppose that $\mu_m(\ksep) \subseteq k$. Show that $\delta$ is an
        isomorphism
        \[
        \delta\colon k^\times/k^{\times m} \to \Hom_{\mathrm{conts}}(\Gal_k, \mu_m(\ksep)).
        \]
        For
        $a \in k^\times$ and $\alpha^m = a$, show that $\ker \delta(a) = \Gal_{k(\alpha)}$,
        and deduce that $k(\alpha)/k$ is Galois and cyclic, of degree the order of $a$ in
        $k^\times/k^{\times m}$.
\begin{solution}
% Write your solution here.
\end{solution}
  \item\label{kummer:it:kth-pairing} Suppose that $\mu_m(\ksep) \subseteq k$. Let $\Delta$ be a subgroup
        of $k^\times$ containing $k^{\times m}$, with $\Delta/k^{\times m}$ finite, and
        $L \coloneqq k(\Delta^{1/m})$ the extension of $k$ inside $\ksep$ generated by the
        $m$-th roots of the elements of $\Delta$. Show that $L/k$ is finite and Galois, and
        that
        \[
        \langle\ ,\ \rangle\colon \Gal(L/k) \times \Delta/k^{\times m} \to \mu_m(\ksep), \qquad
        \langle g, a \rangle \coloneqq g(\alpha)/\alpha \quad (\alpha^m = a),
        \]
        is a well defined bilinear pairing, and that it is perfect: the induced maps
        \[
        \Gal(L/k) \to \Hom(\Delta/k^{\times m}, \mu_m(\ksep)), \qquad
        \Delta/k^{\times m} \to \Hom(\Gal(L/k), \mu_m(\ksep))
        \]
        are isomorphisms. Deduce that
        $\Gal(L/k)$ is abelian of exponent dividing $m$ and that
        $[L : k] = (\Delta : k^{\times m})$.\hint{show that both maps are injective,
        then count with \cref{kummer:prop:characters}.}
\begin{solution}
% Write your solution here.
\end{solution}
  \item\label{kummer:it:kth-bijection} Suppose that $\mu_m(\ksep) \subseteq k$. Show that
        $\Delta \mapsto k(\Delta^{1/m})$ is a bijection from the subgroups $\Delta$ as in
        \ref{kummer:it:kth-pairing} to the finite abelian extensions of $k$ inside $\ksep$ of
        exponent dividing $m$, with inverse
        $L \mapsto \Delta_L \coloneqq k^\times \cap L^{\times m}$.\hint{for $L$ given,
        compute $\Delta_L/k^{\times m}$ with \cref{kummer:prob:kseq}\ref{kummer:it:kseq-res} and
        \cref{kummer:prop:inf-res}.}
\begin{solution}
% Write your solution here.
\end{solution}
  \item\label{kummer:it:kth-cube} ($\star$) Let $k = \Q$, $m = 3$, $\zeta \in \Qbar$ a primitive
        cube root of unity, and $F \coloneqq \Q(\zeta)$. Here $\mu_3(\Qbar) \not\subseteq \Q$.
        \begin{enumerate}
          \item\label{kummer:it:kth-cube-fails} Show that $\Q(\sqrt[3]2)/\Q$ is not Galois,
                although $2 \notin \Q^{\times 3}$. Which step of \ref{kummer:it:kth-cyclic} fails?
\begin{solution}
% Write your solution here.
\end{solution}
          \item\label{kummer:it:kth-cube-inj} Show that $\Q^\times/\Q^{\times 3} \to
                F^\times/F^{\times 3}$ is injective.\hint{\cref{kummer:prop:inf-res} and
                \cref{kummer:prob:cocycles}\ref{kummer:it:coc-order}.}
\begin{solution}
% Write your solution here.
\end{solution}
          \item\label{kummer:it:kth-cube-s3} Let $a \in \Q^\times - \Q^{\times 3}$, $\alpha^3 = a$,
                and $L \coloneqq F(\alpha)$. Show that $L/\Q$ is Galois of degree $6$, and
                that conjugation by any $\tau \in \Gal(L/\Q)$ with $\tau(\zeta) = \zeta^{-1}$
                acts on $\Gal(L/F) \cong \mu_3(F)$, the isomorphism of
                \ref{kummer:it:kth-cyclic} over $F$, as inversion. Deduce that
                $\Gal(L/\Q) \cong S_3$. Check with \Magma:
\begin{magma}
R<x> := PolynomialRing(Rationals());
G := GaloisGroup(x^3 - 2);
#G, IsAbelian(G);                   // 6 false
F<zeta> := CyclotomicField(3);
S<y> := PolynomialRing(F);
#GaloisGroup(y^3 - 2);              // 3
\end{magma}
\begin{solution}
% Write your solution here.
\end{solution}
        \end{enumerate}
\end{enumerate}
\end{problem}


\section{Computing \texorpdfstring{$k^\times/k^{\times m}$}{k^x/k^xm}}
\label{kummer:sec:compute}

This section computes $k^\times/k^{\times m} \cong H^1(k, \mu_m)$ for finite fields,
$\Qp$ and $\Q$, and then cuts out a finite subgroup by conditions at the primes, as in
\cite{Poonen02}*{\S\S 9--10}. The last item is \cite{Silverman09}*{proof of Proposition~VIII.1.6}, with a \Magma\ check.

\begin{wslemma}[{Hensel's lemma \cite{Silverman09}*{Lemma~IV.1.2}}]
\label{kummer:lem:hensel}
Let $F \in \Z_p[w]$ and $a \in \Z_p$ with $F(a) \in p^n\Z_p$ for some $n \geq 1$ and
$F'(a) \in \Z_p^\times$. Then $F$ has a root $b \in \Z_p$ with $b \equiv a \pmod{p^n}$.
\end{wslemma}

\begin{wsdefinition}[{Unramified classes \cite{Poonen02}*{\S\S 4 and 9}}]
\label{kummer:def:unramified}
Let $p$ be a prime. Embed $\Qbar$ in an algebraic closure $\Qbar_p$ of $\Qp$, so that
restriction is an injection $\Gal_{\Qp} \coloneqq \Gal(\Qbar_p/\Qp) \hookrightarrow \Gal_\Q$.
Let $\Qp^{\unr} \subseteq \Qbar_p$ be the maximal unramified extension of $\Qp$, the union of
its finite unramified extensions. A finite extension of $\Qp$ inside $\Qbar_p$ is
unramified if and only if it lies in $\Qp^{\unr}$. The \defi{inertia group} is
$I_p \coloneqq \Gal(\Qbar_p/\Qp^{\unr}) \subseteq \Gal_{\Qp}$. A class
$\xi \in H^1(\Q, A)$, for a commutative algebraic group $A$ over $\Q$, is \defi{unramified
at $p$} if its image in $H^1(I_p, A(\Qbar_p))$ is trivial. For a
finite set $S$ of primes, $H^1(\Q, A; S)$ is the subgroup of the classes unramified at
every prime $p \notin S$.
\end{wsdefinition}

\begin{wsproposition}[{Unramified Kummer extensions \cite{Silverman09}*{proof of Proposition~VIII.1.6}}]
\label{kummer:prop:kummer-unramified}
Let $p$ be a prime with $p \nmid m$, $a \in \Qp^\times$ and $\alpha \in \Qbar_p$ with
$\alpha^m = a$. Then $\Qp(\alpha)/\Qp$ is unramified if and only if $v_p(a) \equiv 0
\pmod m$.
\end{wsproposition}

\begin{problem}[Finite and local fields]
\label{kummer:prob:local}
\leavevmode
\begin{enumerate}
  \item\label{kummer:it:loc-finite} Let $q$ be a power of a prime not dividing $m$. Show that
        $\F_q^\times/\F_q^{\times m}$ is cyclic of order $\gcd(m, q - 1)$. If
        $m \mid q - 1$, so that $\mu_m(\Fbar_q) \subseteq \F_q$, compare the subgroups of
        $\F_q^\times/\F_q^{\times m}$ with the extensions of $\F_q$ of degree dividing $m$,
        as \cref{kummer:prob:kth}\ref{kummer:it:kth-bijection} predicts.
\begin{solution}
% Write your solution here.
\end{solution}
  \item\label{kummer:it:loc-odd} Let $p$ be odd and $u \in \Z_p^\times$ a unit whose reduction is
        not a square in $\F_p$. Show that $\Qp^\times/\Qp^{\times 2}$ has order $4$, with
        representatives $1, u, p, up$. List the quadratic extensions of $\Qp$, and say
        which one is unramified.\hint{every element of $\Qp^\times$ is $p^nv$ with $n \in
        \Z$ and $v \in \Z_p^\times$, and \cref{kummer:lem:hensel} decides which units are
        squares.}
\begin{solution}
% Write your solution here.
\end{solution}
  \item\label{kummer:it:loc-two} Show that a unit $v \in \Z_2^\times$ is a square if and only if
        $v \equiv 1 \pmod 8$. Deduce that $\Q_2^\times/\Q_2^{\times 2}$ has order $8$, with
        representatives $\pm 1, \pm 2, \pm 5, \pm 10$. \Magma\ counts the quadratic
        extensions of $\Q_p$:
\begin{magma}
#AllExtensions(pAdicRing(3, 20), 2);    // 3
#AllExtensions(pAdicRing(5, 20), 2);    // 3
#AllExtensions(pAdicRing(2, 20), 2);    // 7
\end{magma}
        Check these against your answers.\hint{for $v \equiv 1 \pmod 8$, apply
        \cref{kummer:lem:hensel} to $F(w) = w^2 + w - (v - 1)/4$, then look at $(2w + 1)^2$.}
\begin{solution}
% Write your solution here.
\end{solution}
\end{enumerate}
\end{problem}

\begin{wsdefinition}[{The group $k(S, m)$ \cite{Silverman09}*{Theorem~X.1.1(c)}}]
\label{kummer:def:kSm}
Let $k$ be a number field, $S$ a finite set of places of $k$ containing the archimedean
ones, and
\[
k(S, m) \coloneqq \{a \in k^\times/k^{\times m} : v(a) \equiv 0 \pmod m \text{ for all }
v \notin S\},
\]
where $v$ runs over the normalized valuations of $k$. Let $\calO_S$ be the ring of $S$-integers, the
$a \in k$ with $v(a) \geq 0$ for all $v \notin S$, and $\Cl(\calO_S)$ its ideal class
group. For $k = \Q$ we list only the primes in $S$, and in \cite{Poonen02}*{\S 10},
$\Q(S, 2)$ is written $\langle -1, S \rangle$.
\end{wsdefinition}

\begin{wsproposition}[{Finiteness \cite{Silverman09}*{proof of Proposition~VIII.1.6}}]
\label{kummer:prop:finiteness}
The group $\Cl(\calO_S)$ is finite, and $\calO_S^\times$ is finitely generated
(Dirichlet's $S$-unit theorem).
\end{wsproposition}

\begin{problem}[Classes unramified outside $S$]
\label{kummer:prob:global}
\leavevmode
\begin{enumerate}
  \item\label{kummer:it:glob-basis} Show that $\Q^\times/\Q^{\times 2}$ is an $\F_2$-vector space
        with basis the classes of $-1$ and of the primes. In particular, it is infinite.
\begin{solution}
% Write your solution here.
\end{solution}
  \item\label{kummer:it:glob-S} Let $S$ be a finite set of primes. Show that $\Q(S, 2)$ is the
        subgroup generated by $-1$ and the primes in $S$, of order $2^{\#S + 1}$.
\begin{solution}
% Write your solution here.
\end{solution}
  \item\label{kummer:it:glob-unram} Let $p$ be an odd prime and $a \in \Q^\times$. Show that
        $\delta(a) \in H^1(\Q, \mu_2)$ is unramified at $p$ if and only if $v_p(a)$ is
        even. Deduce that $\Q(S, 2) = \delta^{-1}H^1(\Q, \mu_2; S)$ for every finite set
        $S$ of primes containing $2$.
\begin{solution}
% Write your solution here.
\end{solution}
  \item\label{kummer:it:glob-number} ($\star$) Let $k$ be a number field and $S$ as in
        \cref{kummer:def:kSm}. Construct an exact
        sequence\hint{the prime ideals of $\calO_S$ are those of the places outside $S$.
        If $a$ represents an element of $k(S, m)$, then $a\calO_S = \mathfrak a^m$ for a
        fractional ideal $\mathfrak a$ of $\calO_S$.}
        \[
        1 \to \calO_S^\times/\calO_S^{\times m} \to k(S, m) \to \Cl(\calO_S)[m] \to 0,
        \]
        and deduce that $k(S, m)$ is finite. For $k = \Q(\sqrt{-5})$, $m = 2$ and $S$ the
        archimedean place, show that $k(S, 2) = \langle -1, 2 \rangle$, of order $4$, and
        check with \Magma:
\begin{magma}
K<s> := QuadraticField(-5);
OK := Integers(K);
#ClassGroup(OK);                    // 2
G, m := pSelmerGroup(2, {PowerIdeal(OK)|});
#G;                                 // 4
[K | G.i @@ m : i in [1..Ngens(G)]];  // [ -1, 2 ]
\end{magma}
\begin{solution}
% Write your solution here.
\end{solution}
\end{enumerate}
\end{problem}


% \bib, bibdiv, biblist are defined by the amsrefs package.
\begin{bibdiv}
\begin{biblist}

\bib{Magma}{article}{
      author={Bosma, Wieb},
      author={Cannon, John},
      author={Playoust, Catherine},
       title={The {M}agma algebra system. {I}. {T}he user language},
        date={1997},
        ISSN={0747-7171},
     journal={J. Symbolic Comput.},
      volume={24},
      number={3-4},
       pages={235\ndash 265},
         url={http://dx.doi.org/10.1006/jsco.1996.0125},
        note={Computational algebra and number theory (London, 1993)},
      review={\MR{MR1484478}},
}

\bib{Hartshorne77}{book}{
      author={Hartshorne, Robin},
       title={Algebraic geometry},
      series={Graduate Texts in Mathematics},
   publisher={Springer-Verlag, New York-Heidelberg},
        date={1977},
      volume={No. 52},
        ISBN={0-387-90244-9},
      review={\MR{463157}},
}

\bib{lmfdb}{misc}{
      author={{LMFDB Collaboration}, The},
       title={The {L}-functions and modular forms database},
         how={\url{https://www.lmfdb.org}},
        date={2025},
        note={[Online; accessed 21 July 2025]},
}

\bib{Milne20}{book}{
      author={Milne, James~S.},
       title={Elliptic curves},
     edition={Second},
   publisher={World Scientific Publishing Co. Pte. Ltd., Hackensack, NJ},
        date={2020},
}

\bib{Milne22}{book}{
      author={Milne, James~S.},
       title={Fields and {G}alois theory},
   publisher={Kea Books, Ann Arbor, MI},
        date={2022},
}

\bib{Neukirch}{book}{
      author={Neukirch, J\"{u}rgen},
       title={Algebraic number theory},
      series={Grundlehren der Mathematischen Wissenschaften [Fundamental
  Principles of Mathematical Sciences]},
   publisher={Springer-Verlag, Berlin},
        date={1999},
      volume={322},
        ISBN={3-540-65399-6},
         url={https://doi.org/10.1007/978-3-662-03983-0},
        note={Translated from the 1992 German original and with a note by
  Norbert Schappacher, With a foreword by G. Harder},
      review={\MR{1697859}},
}

\bib{Poonen02}{misc}{
      author={Poonen, Bjorn},
       title={The {S}elmer group, the {S}hafarevich--{T}ate group, and the weak
  {M}ordell--{W}eil theorem},
        date={2002},
        note={Lecture notes, Arizona Winter School 1999, February 20, 2002
  version},
}

\bib{Poonen06}{misc}{
      author={Poonen, Bjorn},
       title={Lectures on rational points on curves},
        date={2006},
        note={Lecture notes, March 5, 2006 version},
}

\bib{Poonen2017}{book}{
      author={Poonen, Bjorn},
       title={Rational points on varieties},
      series={Graduate Studies in Mathematics},
   publisher={American Mathematical Society, Providence, RI},
        date={2017},
      volume={186},
        ISBN={978-1-4704-3773-2},
         url={https://doi.org/10.1090/gsm/186},
      review={\MR{3729254}},
}

\bib{Selmer-51}{article}{
      author={Selmer, Ernst~S.},
       title={The {D}iophantine equation {$ax^3+by^3+cz^3=0$}},
        date={1951},
        ISSN={0001-5962,1871-2509},
     journal={Acta Math.},
      volume={85},
       pages={203\ndash 362 (1 plate)},
         url={https://doi.org/10.1007/BF02395746},
      review={\MR{41871}},
}

\bib{Silverman09}{book}{
      author={Silverman, Joseph~H.},
       title={The arithmetic of elliptic curves},
     edition={Second},
      series={Graduate Texts in Mathematics},
   publisher={Springer, Dordrecht},
        date={2009},
      volume={106},
        ISBN={978-0-387-09493-9},
         url={https://doi.org/10.1007/978-0-387-09494-6},
      review={\MR{2514094}},
}

\bib{Voight21}{book}{
      author={Voight, John},
       title={Quaternion algebras},
      series={Graduate Texts in Mathematics},
   publisher={Springer, Cham},
        date={2021},
      volume={288},
        ISBN={978-3-030-56692-0; 978-3-030-56694-4},
         url={https://doi.org/10.1007/978-3-030-56694-4},
      review={\MR{4279905}},
}

\end{biblist}
\end{bibdiv}

\end{document}
